\newpage
\subsection*{Organization of the Appendices}
\begin{itemize}[leftmargin=*]
\item Appendix \ref{appendix:background} introduces relevant notation and background.
\item Appendix \ref{sec:spiked_sample_covariance} states our main results for the general nonlinear spiked covariance model
\[\vK=\vG^\top \vG,\]
formalizing the discussion in Section \ref{sec:proof_overview}. These results
are divided into two subsections: Appendix \ref{subsec:nonasymptotic} gives a
``no outliers'' statement for $\vK$ and a deterministic equivalent approximation
for its resolvent, under minimal asymptotic assumptions.
Appendix \ref{subsec:spikes} then states the main characterizations of spike
eigenvalues/eigenvectors in an asymptotic setting with a spiked eigenstructure.
\item Appendix \ref{appendix:resolvent} develops a general fluctuation averaging
lemma for the sample covariance model, and proves the results of
Appendix \ref{subsec:nonasymptotic}.
\item Appendix \ref{appendix:spike} proves the
results of Appendix \ref{subsec:spikes}.
\item Finally, Appendix \ref{sec:CK_spike_proof} proves the results of Section \ref{sec:spike_CK}
on propagation of spiked eigenstructure through the layers of a neural network, and
Appendix \ref{sec:trained_CK_proof} proves the results of Section
\ref{sec:trained_CK} on the eigenstructure of the CK after gradient descent
training.
\end{itemize}

\bigskip

\section{Notations and background}\label{appendix:background}

\subsection{Stochastic domination}

We use the following standard notation for stochastic domination of random
variables, see e.g.\ \cite[Definition 2.4]{erdHos2013averaging}:
For random variables $X \equiv X(u)$ and $Y \equiv Y(u) \geq 0$ depending implicitly on $N$
and a parameter $u \in U_N$, as $N \to \infty$, we write
\[X \prec Y \text{ or } X=\SD{Y} \text{ uniformly over } u \in U_N\]
if, for any fixed $\eps,D>0$ and all large $N$,
\[\sup_{u\in U_N}\P \Big[|X(u)|>N^\eps Y(u)\Big]<N^{-D}.\]
Throughout, ``for all large $N$'' means for all $N \geq N_0$ where
$N_0$ may depend on $\eps,D$, any quantities that are constant in the context
of the statement, and convergence rates of the spike eigenvalues and
empirical spectral measures in the given assumptions.

If $X=\bI\{\cE\}$ is the indicator of an event $\cE \equiv \cE_N$, then
$\bI\{\cE\} \prec 0$ means $\P[\cE]<N^{-D}$ for any fixed $D>0$ and all large
$N$. If $X$ and $Y$ are both deterministic, then $X \prec Y$ means
$|X| \leq N^\eps Y$ (deterministically) for any $\eps>0$ and all large $N$.
For an event $\cE \equiv \cE_N$, we will write
\[X=\SDcE{Y}\]
as shorthand for $X \cdot \bI\{\cE\} \prec Y$.

We will use the following basic properties often implicitly.

\begin{proposition}\label{prop:domination}
Suppose $X \prec Y$ uniformly over $u \in U_N$.
\begin{enumerate}[label=(\alph*)]
\item If $|U_N|\leq N^C$ for a constant $C>0$, then
for any fixed $\eps,D>0$ and all large $N$,
\[\P \Big[\text{there exists } u\in U_N \text{ with } |X(u)|
\geq N^\eps Y(u)\Big] \leq N^{-D}.\]
\item If $|U_N| \leq N^C$ for a constant $C>0$, then
$\sum_{u \in U_N} X(u) \prec \sum_{u \in U_N} Y(u)$.
\item If $|U_N| \leq C$ for a constant $C>0$, then
$\prod_{u \in U_N} X(u) \prec \prod_{u \in U_N} Y(u)$.
\item If $Y$ is deterministic, and $\E[X^2] \leq N^C$ and $Y \geq N^{-C}$
for a constant $C>0$, then also $\E[|X|] \prec Y$ uniformly over $u \in U_N$.
\end{enumerate}
\end{proposition}
\begin{proof}
The first three statements follow from a union bound over $U_N$. For the last
statement, for any fixed $\eps>0$, observe that
\[\E|X| \leq N^{\eps/2}Y+\E\Big[|X|\bI\{|X|>N^{\eps/2}Y\}\Big]
\leq N^{\eps/2}Y+\E[X^2]^{1/2}\P[|X|>N^{\eps/2}Y]^{1/2}.\]
Applying $\E[X^2] \leq N^C$, $Y \geq N^{-C}$, and
$\P[|X|>N^{\eps/2}Y]<N^{-D}$ for sufficiently large $D>0$ shows that the
second term is less than $N^{\eps/2}Y$ for all large $N$, hence $\E|X|<N^\eps
Y$.
\end{proof}

\subsection{Deformed Marcenko-Pastur law}

For a probability measure $\nu$ supported on $[0,\infty)$ and an aspect ratio
parameter $\gamma>0$, consider the deformed Marcenko–Pastur measure
\[\mu=\rho^{\MP}_{\gamma}\boxtimes \nu\]
and its ``companion'' probability measure
\[\tilde \mu=\gamma\mu+(1-\gamma)\delta_0.\]
Here, $\mu$ and $\tilde\mu$ represent the limit eigenvalue distributions of
$\vG^\top \vG \in \R^{n \times n}$ and
$\vG\vG^\top \in \R^{N \times N}$ respectively, when
$\vG=\frac{1}{\sqrt{N}}[\vg_1,\ldots,\vg_N] \in \R^{N
\times n}$ has i.i.d.\ rows with mean 0 and covariance $\vSigma$, and
$n,N \to \infty$ with $n/N \to \gamma$ and 
$\frac{1}{n}\sum_{i=1}^n \delta_{\lambda_i(\vSigma)} \to \nu$ weakly.

These measures $\mu,\tilde\mu$ may be defined by their Stieltjes transform
\begin{equation}\label{eq:stieltjes}
m(z)=\int \frac{1}{x-z}d\mu(x), \qquad
\tilde m(z)=\int \frac{1}{x-z}d\tilde \mu(x)
\end{equation}
where $\tilde m(z)=\gamma m(z)+(1-\gamma)(-1/z)$. 
By the results of \citep{marchenko1967,silverstein1995empirical},
for any $z \in \C^+$, $m(z)$ and $\tilde m(z)$ are the unique roots in
$\{m \in \C:\gamma m+(1-\gamma)(-1/z) \in\C^+\}$ and $\C^+$, respectively,
to the \textit{Marcenko-Pastur equations}
\begin{equation}\label{eq:MPeq}
m(z)=\int \frac{1}{\lambda(1-\gamma-\gamma z m(z))-z}\,d\nu(\lambda), \quad
z=-\frac{1}{\tilde m(z)}+\gamma \int \frac{\lambda}{1+\lambda
\tilde m(z)}d\nu(\lambda).
\end{equation}
We define $m(z),\tilde m(z)$ via (\ref{eq:stieltjes}) also on the full domains
$\C \setminus \supp{\mu}$ and $\C \setminus \supp{\tilde \mu}$ respectively,
where the support sets $\supp{\mu}$ and $\supp{\tilde \mu}$ may differ only at
the single point $\{0\}$.

In the setting $\vSigma=\vI$ (and $\nu=\delta_1$), the law $\mu=\rho_\gamma^\MP$
is the standard Marcenko-Pastur law, with explicit density function with respect
to Lebesgue measure
\[d\rho_\gamma^\MP(\lambda)=\frac{1}{2\pi}\frac{\sqrt{(\lambda_{+}-\lambda)(\lambda-\lambda_{-})}}{\gamma\lambda}\cdot\bI_{\lambda
\in [\lambda_{-},\lambda_+]}\,d\lambda, \qquad \lambda_{\pm}:=(1\pm\sqrt{\gamma})^2\]
for $\gamma \leq 1$, and an additional point mass $(1-1/\gamma)$ at 0
when $\gamma>1$.

In general, $\mu$ and $\tilde \mu$ do not have analytically
explicit densities. However, $\supp{\tilde \mu}$ is explicitly
characterized in \cite{silverstein1995analysis}, and we review this
characterization here: Define
\begin{equation}\label{eq:zdomain}
\cT=\{0\} \cup \{-1/\lambda:\lambda \in \supp{\nu}\}.
\end{equation}
For $\tilde m \in \C \setminus \cT$, define
\begin{equation}\label{eq:inv_z}
z(\tilde m)={-}\frac{1}{\tilde m}+\gamma \int\frac{\lambda}{1+\lambda \tilde m}d\nu(\lambda).
\end{equation}	
In light of the second equation of (\ref{eq:MPeq}), this may be understood as
a formal inverse of $\tilde m(z)$.
From~\cite[Theorems 4.1 and 4.2]{silverstein1995analysis},
we have the following properties.

\begin{proposition}\label{prop:inv_z}
$\tilde m(\cdot)$ defines a bijection from $\{z \in \R \setminus \supp{\tilde
\mu}\}$ to $\{\tilde m \in \R \setminus \cT:z'(\tilde m)>0\}$, whose inverse
function is $z(\cdot)$. In particular, $x \in \R$ does not belong to
$\supp{\tilde \mu}$ if and only if there exists
$\tilde m \in \R \setminus \cT$
such that $z'(\tilde m)>0$ and $z(\tilde m)=x$.
\end{proposition}

\subsection{Additional notation}

For a probability measure $\mu$, its support is the closed set
\[\supp{\mu}=\{x \in \R:\mu(O)>0 \text{ for any open neighborhood } O \ni x\}.\]
We write $\dist{x,A}=\inf\{|x-y|:y \in A\}$ and define the $\eps$-neighborhood
\[\supp{\mu}+(-\eps,\eps)=\{x \in \R:\dist{x,\supp{\mu}}<\eps\}.\]
We write $\delta_x$ for the probability measure given by a point mass at
$x \in \R$, $a\mu_0+(1-a)\mu_1$ for the convex combination of $\mu_0,\mu_1$,
and $a \otimes \mu \oplus b$ for
the law of $a\mathsf{x}+b$ when $\mathsf{x} \sim \mu$.

For vectors, $\|\vv\| \equiv \|\vv\|_2$ is the Euclidean norm. For matrices,
$\|\vM\|$ is the
operator norm $\sup_{\vv:\|\vv\|=1} \|\vM\vv\|$, $\|\vM\|_F$ is the
Frobenius norm $(\Tr \vM^\top\vM)^{1/2}$, $\Tr$ is the (unnormalized)
matrix trace, and $\vA \odot \vB$ is the entrywise (Hadamard) product.
We write $\diag(\vv)$ for the diagonal matrix with vector $\vv$ along the
main diagonal, and $\vI_n$ for the $n \times n$ identity matrix.

%%%%%%%%%%%%%%%%%%%%%%%%%%%%%%%%%%%%%%%%%%%%%%%%%%%%%%%%%%%%%%%%%%%%%%%%%%%%
%%%%%%%%%%%%%%%%%%%%%%%%%%%%%%%%%%%%%%%%%%%%%%%%%%%%%%%%%%%%%%%%%%%%%%%%%%%%

\section{Results for the nonlinear spiked covariance model}\label{sec:spiked_sample_covariance}

\subsection{Deterministic equivalent for the resolvent}\label{subsec:nonasymptotic}

We consider the sample covariance and Gram matrix
\[\vK=\vG^\top\vG \in \R^{n \times n}, \quad
\widetilde \vK=\vG\vG^\top \in \R^{N \times N}, \quad \text{ where } \quad
\vG=\frac{1}{\sqrt{N}}[\vg_1,\ldots,\vg_N] \in \R^{N \times n}.\]
The following are our basic assumptions, where we
recall that $\bI\{\cE\} \prec 0$ means $\P[\cE] \leq N^{-D}$ for any fixed $D>0$
and all large $N$.
\begin{assumption}\label{assump:G}
\phantom{}
The rows of $\vG$ are independent and satisfy
$\E[\vg_i]=\mathbf{0}$ and $\E[\vg_i\vg_i^\top]=\vSigma$ for all $i\in[N]$,
such that:
\begin{enumerate}[label=(\alph*)]
\item There exist constants $C,c>0$ such that $c<n/N<C$ and
$\|\vSigma\|<C$.
\item There exists a constant $B>0$ such that
$\bI\{\|\vK\|>B\} \prec 0$.
\item Uniformly over deterministic matrices $\vA \in \C^{n \times n}$ and over
$i \neq j \in [N]$,
\[\vg_i^\top \vA \vg_i - \E[\vg_i^\top \vA \vg_i]
\prec \|\vA\|_F, \qquad \vg_i^\top \vA \vg_j \prec \|\vA\|_F.\]
\item For any integer $\alpha>0$, there exists a constant $C=C(\alpha)>0$
such that $\E[\|\vg_i\|^\alpha] \leq N^C$.
\end{enumerate}
\end{assumption}

Denote the finite-$N$
dimension ratio and empirical eigenvalue distribution of $\vSigma$ by
\begin{equation}\label{eq:gammanuN}
\gamma_N=\frac{n}{N}, \qquad \nu_N=\frac{1}{n}\sum_{i=1}^n
\delta_{\lambda_i(\vSigma)}.
\end{equation}
Let
\[\mu_N=\rho_{\gamma_N}^\MP \boxtimes \nu_N, 
\qquad \tilde \mu_N=\gamma_N\mu_N+(1-\gamma_N)\delta_0.\]
Denote the Stieltjes transforms of $\mu_N,\tilde \mu_N$ by
$m_N(z),\tilde m_N(z)$. These are characterized exactly as in \eqref{eq:MPeq}
with $(\gamma_N,\nu_N)$ in place of $(\gamma,\nu)$.

We first establish that with high probability,
$\vK$ and $\widetilde \vK$ have no outlier eigenvalues far from the support set
\begin{equation}\label{eq:support}
\cS_N=\supp{\mu_N} \cup \{0\}=\supp{\tilde \mu_N} \cup \{0\}.
\end{equation}
\begin{theorem}\label{thm:no_outlier}
Suppose Assumption \ref{assump:G} holds. Then for any fixed $\eps>0$,
\begin{align*}
\bI\Big\{\vK \text{ has an eigenvalue outside }
\cS_N+(-\eps,\eps)\Big\} \prec 0.
\end{align*}
\end{theorem}
In asymptotic settings where $\nu_N \to \nu$ and $\mu_N \to \mu$ weakly and
$\vSigma$ has no spike eigenvalues, this set $\cS_N$ will converge to
$\cS:=\supp{\mu} \cup \{0\}$. In general, $\cS_N$ may contain intervals around
spike eigenvalues of $\vK$ that are separated from $\supp{\mu} \cup \{0\}$
if $\vSigma$ has a spiked structure, and this will be clarified in the
subsequent section.

Next, we establish a deterministic equivalent approximation for the resolvent of
$\vK$, for spectral arguments separated from this support set $\cS_N$.
Let us denote by
\[\vR(z)=(\vK-z\vI)^{-1}, \qquad
m_{\vK}(z)=\frac{1}{n} \Tr \vR(z)\]
the resolvent and Stieltjes transform of $\vK$ for $z\not\in\supp{\mu_N}$.
For any $\eps>0$, define the domain
\begin{equation}\label{eq:U_e}
U_N(\eps)=\Big\{z \in \C:\,|z| \leq \eps^{-1},\,
\dist{z,\cS_N} \geq \eps\Big\}.
\end{equation}

\begin{theorem}\label{thm:deterministic_equ}
Suppose Assumption \ref{assump:G} holds. Then for any fixed $\eps>0$,
uniformly over $z \in U_N(\eps)$ and over deterministic matrices
$\vA \in \C^{n \times n}$, we have
\[m_{\vK}(z)-m_N(z) \prec \frac{1}{N},
\qquad \Tr \Big[\vR(z)\vA-(-z\tilde m_N(z)\vSigma-z\vI)^{-1}\vA\Big]
\prec \frac{1}{\sqrt{N}}\,\|\vA\|_F.\]
\end{theorem}
For spectral arguments $z \in \C \setminus \R_+$ separated from the positive
real line, such a result has been shown recently in
\cite{chouard2022quantitative,schroder2023deterministic} (using different proof
techniques). We use Theorem \ref{thm:no_outlier} as an input to establish this
approximation also for spectral arguments in $\R_+ \setminus \cS_N$, as such a
result (and its extension to a generalized resolvent) is needed
for our analysis of spiked eigenstructure to follow.

\subsection{Spike eigenvalues and eigenvectors}\label{subsec:spikes}

Now we consider an asymptotic setting with a specific spiked structure for the
population covariance matrix $\vSigma$, having a fixed number of spikes outside the
support of the weak limit of its spectral law.
This assumption is summarized as follows.

\begin{assumption}\label{assum:spike}
$\vSigma$ has eigenvalues $\lambda_1(\vSigma),\ldots,\lambda_n(\vSigma)$ (not
necessarily ordered by magnitude) where,
for a fixed integer $r \geq 0$, as $N \to \infty$:
\begin{enumerate}[label=(\alph*)]
\item $n/N \to \gamma \in (0,\infty)$.
\item There exists a probability measure $\nu$ with compact support in
$(0,\infty)$, such that
\[\frac{1}{n-r}\sum_{i=r+1}^n \delta_{\lambda_i(\vSigma)} \to \nu \text{
weakly.}\]
Furthermore, for any fixed $\eps>0$ and all large $N$,
\[\lambda_i(\vSigma) \in \supp{\nu}+(-\eps,\eps) \text{ for all }
i \geq r+1.\]
\item There exist distinct values $\lambda_1,\ldots,\lambda_r>0$ with
$\lambda_1,\ldots,\lambda_r \not\in\supp{\nu}$ such that
\[\lambda_i(\vSigma) \to \lambda_i \text{ for all } i=1,\ldots,r.\]
\end{enumerate}
\end{assumption}

Under this assumption, we analyze the outlier singular values of $\vG$ and their
corresponding singular vectors. Let
\[\gamma_{N,0}=\frac{n-r}{N}, \qquad \nu_{N,0}=\frac{1}{n-r}\sum_{i=r+1}^n
\delta_{\lambda_i(\vSigma)}\]
be the finite-$N$ aspect ratio and population spectral measure corresponding to
the bulk component of $\vSigma$. Define the laws
\[\mu_{N,0}=\rho_{\gamma_{N,0}}^{\MP} \boxtimes\nu_{N,0}, \qquad
\tilde \mu_{N,0}=\gamma_{N,0}\mu_{N,0}+(1-\gamma_{N,0})\delta_0\]
and let $m_{N,0}(z),\tilde m_{N,0}(z)$ be their Stieltjes transforms.
In the setting of Assumption \ref{assum:spike}, we note that
$\mu_{N,0} \to \mu=\rho_\gamma^{\MP} \boxtimes \nu$ and
$\tilde \mu_{N,0} \to \tilde \mu=\gamma \mu+(1-\gamma)\delta_0$ weakly as
$N \to \infty$, where the Stieltjes transforms $m(z),\tilde m(z)$ of these
limits $\mu,\tilde \mu$ are characterized by \eqref{eq:MPeq}.

Denote the limit support set
\begin{equation}\label{eq:limitsupport}
\cS=\supp{\mu} \cup \{0\}=\supp{\tilde \mu} \cup \{0\}.
\end{equation}
Under Assumption \ref{assum:spike} when $r=0$, i.e.\ $\vSigma$ does
not have spike eigenvalues, the following is a
corollary of Theorem \ref{thm:no_outlier}. A similar ``no outlier'' statement
has been shown for linearly defined sample covariance models in
\citep{bai1998no}.

\begin{coro}\label{cor:no_outlier}
Suppose Assumptions \ref{assump:G} and \ref{assum:spike} hold, where $r=0$.
Then for any fixed $\eps>0$,
\begin{align*}
\bI\Big\{\vK \text{ has an eigenvalue outside }
\cS+(-\eps,\eps)\Big\} \prec 0.
\end{align*}
\end{coro}

We now give a more quantitative version of Theorem \ref{thm:informal} stated
informally in Section \ref{sec:proof_overview}, which describes the
spike eigenvalues of $\vK=\vG^\top \vG$ and corresponding singular vectors of $\vG$ when
there are possibly spike eigenvalues in $\vSigma$. Define the domain
\[\cT_{N,0}=\{0\} \cup \{-1/\lambda:\lambda \in \supp{\nu_{N,0}}\}.\]
For $\tilde m \in \C \setminus \cT_{N,0}$, define the functions
\begin{equation}\label{eq:zN0}
z_{N,0}(\tilde m)=-\frac{1}{\tilde m}+\gamma_{N,0}
\int \frac{\lambda}{1+\lambda \tilde m}d\nu_{N,0}(\lambda), \qquad
\varphi_{N,0}(\tilde m)={-}\frac{\tilde m z_{N,0}'(\tilde m)}{z_{N,0}(\tilde m)}.
\end{equation}
We note that under Assumption \ref{assum:spike}, the domain $\cT_{N,0}$
converges in Hausdorff distance to $\cT$ as defined in \eqref{eq:zdomain}.
We will verify in the proof (c.f.\ Lemma \ref{lemma:supportcompareb})
that $z_{N,0}(\tilde m) \to z(\tilde m)$ 
and $z_{N,0}'(\tilde m) \to z'(\tilde m)$ for each fixed
$\tilde m \in \C \setminus \cT$,
where $z(\cdot)$ is as defined in \eqref{eq:inv_z}. Then also
$\varphi_{N,0}(\tilde m) \to \varphi(\tilde m)$ for the limiting function
\begin{equation}\label{eq:phi_limit}
\varphi(\tilde m)={-}\frac{\tilde m z'(\tilde m)}{z(\tilde m)}.
\end{equation}

\begin{theorem}\label{thm:spiked}
Suppose Assumptions \ref{assump:G} and \ref{assum:spike} hold. Let
\[\cI=\big\{i \in \{1,\ldots,r\}:z'(-1/\lambda_i)>0\big\}.\]
\begin{enumerate}[label=(\alph*)]
\item For any sufficiently small constant $\eps>0$ and all large $N$,
on an event $\cE \equiv \cE_N$ satisfying $\bI\{\cE^c\} \prec 0$,
there is a 1-to-1 correspondence between the eigenvalues of
$\vK$ outside $\cS+(-\eps,\eps)$ and $\{\lambda_i:i \in \cI\}$.
Denoting these eigenvalues of $\vK$ by
$\{\widehat \lambda_i:i \in \cI\}$, we have
\[\widehat \lambda_i-z_{N,0}(-1/\lambda_i(\vSigma))=\SDcE{\frac{1}{\sqrt{N}}}\]
for each $i \in \cI$, where
$z_{N,0}(-1/\lambda_i(\vSigma)) \to z(-1/\lambda_i)>0$ as $N \to \infty$.
\item On this event $\cE$, for each $i \in \cI$, let $\widehat \vv_i \in \R^n$
be a unit-norm eigenvector of $\vK$ (i.e.\ right singular vector of $\vG$)
corresponding to its eigenvalue $\widehat \lambda_i$, and let $\vv_i$ be a unit-norm
eigenvector of $\vSigma$ corresponding to $\lambda_i(\vSigma)$.
Then, uniformly over (deterministic) unit vectors $\vv\in\R^n$,
\begin{align}
|\vv^\top \widehat\vv_i|-\sqrt{\varphi_{N,0}
({-}1/\lambda_i(\vSigma))} \cdot |\vv^\top\vv_i|=\SDcE{\frac{1}{\sqrt{N}}}
\end{align}
where $\varphi_{N,0}({-}1/\lambda_i(\vSigma)) \to \varphi({-}1/\lambda_i)>0$
as $N \to \infty$. In particular, for each $i \in \cI$,
$|\vv_i^\top\widehat{\vv}_i|^2 \to \varphi({-}1/\lambda_i)$
and $\sup_{j \in [n]:j \neq i} |\vv_j^\top\widehat{\vv}_i|^2 \to 0$
almost surely as $N \to \infty$.
\item Let $\vu=\frac{1}{\sqrt{N}}(u_1,\ldots,u_N)^\top \in\R^N$ be a random
vector such that $[\vu,\vG] \in \R^{N \times (n+1)}$ has independent rows
also satisfying Assumption~\ref{assump:G}. Denote by $\E[u\vg] \in \R^n$
the common value of $\E[u_j\vg_j]$ for all $j \in [N]$.

On this event $\cE$, for each $i \in \cI$, let $\widehat \vu_i \in \R^N$ be a unit-norm
eigenvector of $\widetilde\vK$ (i.e.\ left singular vector of $\vG$)
corresponding to its eigenvalue $\widehat \lambda_i$,
and let $\vv_i$ be the eigenvector of $\vSigma$ as in part (b). Then
\begin{align}
|\vu^\top\widehat{\vu}_i|-
\frac{\sqrt{z_{N,0}(-1/\lambda_i(\vSigma))\varphi_{N,0}({-}1/\lambda_i(\vSigma))}}{\lambda_i(\vSigma)}\cdot\left|\E[u\vg]^\top\vv_i\right|=\SDcE{\frac{1}{\sqrt{N}}}.
\end{align}
\end{enumerate} 
\end{theorem}


%%%%%%%%%%%%%%%%%%%%%%%%%%%%%%%%%%%%%%%%%%%%%%%%%%%%%%%%%%%%%%%%%%%%%%%%%%%%
%%%%%%%%%%%%%%%%%%%%%%%%%%%%%%%%%%%%%%%%%%%%%%%%%%%%%%%%%%%%%%%%%%%%%%%%%%%%


\section{Analysis of the resolvent}\label{appendix:resolvent}

We prove the results of Appendix \ref{subsec:nonasymptotic}. Appendix
\ref{sec:fluctuation} first develops a fluctuation averaging lemma for
the sample covariance model. Appendix \ref{sec:nooutliers} applies this lemma
within the arguments of \cite{bai1998no}, to prove the ``no outliers'' result
of Theorem \ref{thm:no_outlier}. Appendix \ref{sec:deterministicequiv} uses
Theorem \ref{thm:no_outlier} and a second application of the fluctuation
averaging lemma to prove the deterministic equivalent
approximation of Theorem \ref{thm:deterministic_equ}.

\subsection{Fluctuation averaging lemma}\label{sec:fluctuation}

Recall the definitions
\[\vK=\vG^\top \vG, \qquad \widetilde \vK=\vG\vG^\top.\]
For $S \subset [N]$, let $\vG^{(S)} \in \R^{(N-|S|) \times n}$ be the matrix
obtained by removing the rows of $\vG$ corresponding to $i \in S$, and define
\[\vK^{(S)}={\vG^{(S)}}^\top \vG^{(S)}
=\frac{1}{N}\sum_{i \in [N] \setminus S} \vg_i\vg_i^\top \in \R^{n \times n}.\]
Then, for $\vGamma \in \C^{n \times n}$, define
\begin{equation}\label{eq:leaveoneout}
\begin{gathered}
\qquad \vR^{(S)}(\vGamma)=(\vK^{(S)}-\vGamma)^{-1},
\qquad m_{\vK}^{(S)}(\vGamma)=\frac{1}{n}\Tr \vR^{(S)}(\vGamma),\\
\tilde m_{\vK}^{(S)}(\vGamma)=\gamma_N m_{\vK}^{(S)}(\vGamma)
+(1-\gamma_N)\left({-}\frac{1}{z}\right)
=\frac{1}{N}\,\Tr\vR^{(S)}(\vGamma)+\left(1-\frac{n}{N}\right)\left({-}\frac{1}{z}\right).
\end{gathered}
\end{equation}
Importantly, these quantities are independent of $\{\vg_i:i \in S\}$.
We say that $\vR^{(S)}(\vGamma)$ exists (and hence also $m_{\vK}^{(S)},\tilde m_{\vK}^{(S)}$ exist) when $\vK^{(S)}-\vGamma$ is invertible.
For simplicity, we write $\vR=\vR^\emptyset$, $\vR^{(i)}=\vR^{(\{i\})}$,
$\vR^{(Si)}=\vR^{(S \cup \{i\})}$, and similarly for $m_{\vK}$ and $\tilde
m_{\vK}$.

\begin{lemma}\label{lemma:fluctuationavg}
Suppose Assumption \ref{assump:G} holds. Suppose also that there are constants
$C_0,c_0,\delta,\upsilon>0$, $N$-dependent
domains $U \subset \C \setminus \{0\}$ and
$\cD_\Gamma,\cD_A \subseteq \C^{n \times n}$, and $N$-dependent maps
$\Phi_N:\cD_\Gamma \times \cD_A \to (N^{-\upsilon},N^\upsilon)$ and
$\Psi_N:\cD_\Gamma \to (N^{-\upsilon},N^{1-\delta})$,
such that for any fixed $L \geq 1$, the events
\begin{align}
\cE(z,\vGamma,\vA,S)&=\Big\{\vR^{(S)}(\vGamma) \text{ exists, }
\|\vR^{(S)}(\vGamma)\vA\|_F \leq \Phi_N(\vGamma,\vA),\;
\|\vR^{(S)}(\vGamma)\|_F \leq \Psi_N(\vGamma),\nonumber\\
&\quad
\|(z^{-1}\vGamma+\tilde m_{\vK}^{(S)}(\vGamma) \vSigma)^{-1}\| \leq C_0, \text{ and }
|1+N^{-1} \vg_j^\top \vR^{(S)}(\vGamma)\vg_j| \geq c_0 \text{ for all } j \in S \Big\}\label{eq:ESdef}
\end{align}
satisfy $\bI\{\cE(z,\vGamma,\vA,S)^c\} \prec 0$
uniformly over $z \in U$, $\vGamma \in \cD_\Gamma$, $\vA \in \cD_A$,
and $S \subset [N]$ with $|S| \leq L$.

Then, denoting by $\E_{\vg_i}$ the partial expectation over only $\vg_i$
(i.e.\ conditional on $\{\vg_j\}_{j \neq i}$),
also uniformly over $z \in U$, $\vGamma \in \cD_\Gamma$, and $\vA \in \cD_A$,
\begin{equation}\label{eq:FAexpression}
\frac{1}{N}\sum_{i=1}^N (1-\E_{\vg_i})\big[\vg_i^\top \vR^{(i)}(\vGamma)\vA
(z^{-1}\vGamma+\tilde m_{\vK}^{(i)}(\vGamma)\vSigma)^{-1}\vg_i\big]
\prec \max\left(\frac{\Psi_N(\vGamma)}{N},\frac{1}{\sqrt{N}}\right)
\cdot \Phi_N(\vGamma,\vA).
\end{equation}
\end{lemma}

We remark that applying Assumption~\ref{assump:G}(c) and the conditions
of $\cE(z,\vGamma,\vA,i)$ separately to each summand of the left side of
~\eqref{eq:FAexpression} gives the naive bound
\begin{align*}
&\frac{1}{N}\sum_{i=1}^N (1-\E_{\vg_i})[\vg_i^\top \vR^{(i)}(\vGamma)\vA
(z^{-1}\vGamma+\tilde m_{\vK}^{(i)}(\vGamma)\vSigma)^{-1}\vg_i]\\
&\prec \max_{i=1}^N \|\vR^{(i)}(\vGamma)\vA\|_F
\cdot \|(z^{-1}\vGamma+\tilde m_{\vK}^{(i)}(\vGamma)\vSigma)^{-1}\|
\prec \Phi_N(\vGamma,\vA).
\end{align*}
The content of the lemma is to improve this by the additional
factor of $\max(\frac{\Psi_N(\vGamma)}{N},\frac{1}{\sqrt{N}})  \ll 1$.

In this work, we will apply Lemma \ref{lemma:fluctuationavg} only to
spectral arguments $z$ with $O(1)$-separation from $\supp{\mu_N}$
(and matrices $\vGamma=z\vI$ or a finite-rank perturbation thereof), in which
case we will take $\Psi_N(\vGamma)=C/\sqrt{N}$ for a constant $C>0$.
For full-rank matrices $\vA$ 
having bounded operator norm, we will take also
$\Phi_N(\vGamma,\vA)=C/\sqrt{N}$, whereas for finite-rank matrices $\vA$
we will take $\Phi_N(\vGamma,\vA)=C$. We state the result
here more abstractly, as it may be of independent interest to prove local
laws in this nonlinear sample covariance model for spectral arguments $z$
that approach $\supp{\mu_N}$.

In the remainder of this section, we prove Lemma \ref{lemma:fluctuationavg}.
Fix $z \in U$, $\vGamma \in \cD_\Gamma$, and $\vA \in \cD_A$, and
write as shorthand
\[\vR^{(S)}=\vR^{(S)}(\vGamma),\qquad
\tilde m^{(S)}=\tilde m_{\vK}^{(S)}(\vGamma), \qquad
\vOmega^{(S)}=(z^{-1}\vGamma+\tilde m_{\vK}^{(S)}(\vGamma)\vSigma)^{-1},\]
\[\Phi_N=\Phi_N(\vGamma,\vA), \qquad
\Psi_N=\Psi_N(\vGamma),\qquad
\cE(S)=\cE(z,\vGamma,\vA,S).\]
All subsequent instances of $\prec$ will be
implicitly uniform over $z \in U$,
$\vGamma \in \cD_\Gamma$, and $\vA \in \cD_A$.
Define the quantities, for $i \in S$, $j,k \in S \setminus \{i\}$, and
$d \geq 0$,
\begin{align*}
Y_i^{(S)}[d]&=\Tr (\vg_i\vg_i^\top-\vSigma)\vR^{(S)}\vA\vOmega^{(S)}
[\vSigma\vOmega^{(S)}]^d,\\
Z_{ijk}^{(S)}[d]&=N^{-1}\Tr (\vg_i\vg_i^\top-\vSigma)\vR^{(S)}\vg_j\vg_k^\top
\vR^{(S)}\vA\vOmega^{(S)}[\vSigma\vOmega^{(S)}]^d,\\
B_{jk}^{(S)}&=N^{-1}\vg_j^\top\vR^{(S)}\vg_k,\\
C_{jk}^{(S)}&=N^{-2}\vg_j^\top(\vR^{(S)})^2\vg_k,\\
Q_j^{(S)}&=(1+N^{-1} \vg_j^\top \vR^{(S)} \vg_j)^{-1}.
\end{align*}
For each $L \geq 1$, define also the event
\begin{equation}\label{eq:EL}
\cE_L=\bigcap_{S \subset [N]:|S| \leq L} \cE(S).
\end{equation}

\begin{lemma}\label{lemm:bounds}
For any fixed $L,D \geq 1$, uniformly over $S \subset [N]$
with $|S| \leq L$, and over $i \in S$ and
$j,k \in S \setminus \{i\}$ and $d \leq D$,
\begin{equation}\label{eq:FAbounds}
\begin{gathered}
Y_i^{(S)}[d]=\SDcES{\Phi_N},
\quad Z_{ijk}^{(S)}[d]=\SDcES{N^{-1}\Psi_N\Phi_N},\\
\quad B_{jk}^{(S)}=\SDcES{N^{-1}\Psi_N} \text{ for } j \neq k,
\quad C_{jk}^{(S)}=\SDcES{N^{-2}\Psi_N^2},
\quad Q_j^{(S)}=\SDcES{1}.
\end{gathered}
\end{equation}
Furthermore, for any $\alpha>0$, there exists a constant
$C=C(\alpha,L,D)>0$ such that
\begin{equation}\label{eq:FAEbounds}
\begin{gathered}
\E\big[|Y_i^{(S)}[d]|^\alpha\bI\{\cE(S)\}\big]<N^C, \quad
\E\big[|Z_{ijk}^{(S)}[d]|^\alpha\bI\{\cE(S)\}\big]<N^C,\\
\E\big[|B_{jk}^{(S)}|^\alpha\bI\{\cE(S)\}\big]<N^C, \quad
\E\big[|C_{jk}^{(S)}|^\alpha\bI\{\cE(S)\}\big]<N^C, \quad
\E\big[|Q_j^{(S)}|^\alpha\bI\{\cE(S)\}\big]<N^C.
\end{gathered}
\end{equation}
\end{lemma}
\begin{proof}
On the event $\cE(S)$, we have by definition $Q_j^{(S)} \leq 1/c_0$, so the
two statements for $Q_j^{(S)}$ hold immediately. The remaining statements of
\eqref{eq:FAEbounds} follow easily from Holder's inequality, the moment bounds
for $\|\vg_i\|$ in Assumption \ref{assump:G}(d),
the bound $\|\vSigma\|<C$ in Assumption \ref{assump:G}(a),
and the conditions $\|\vR^{(S)}\vA\| \leq \Phi_N \leq N^\upsilon$,
$\|\vR^{(S)}\|_F \leq \Psi_N \leq N$, and $\|\vOmega^{(S)}\| \leq C_0$
defining $\cE(S)$.

For the bounds for $B_{jk}^{(S)}$ and $C_{jk}^{(S)}$ in \eqref{eq:FAbounds},
note that when $j \neq k$, Assumption \ref{assump:G}(c) implies $B_{jk}^{(S)}
\prec N^{-1}\|\vR^{(S)}\|_F$ and $C_{jk}^{(S)} \prec N^{-2}\|(\vR^{(S)})^2\|_F
\leq N^{-2}\|\vR^{(S)}\|_F^2$.
When $j=k$, Assumption \ref{assump:G}(c) implies also
\begin{align*}
C_{jj}^{(S)} &\prec N^{-2}|\Tr \vSigma(\vR^{(S)})^2|+N^{-2}\|(\vR^{(S)})^2\|_F\\
&\leq N^{-2}\|\vSigma \vR^{(S)}\|_F \|\vR^{(S)}\|_F+N^{-2}\|\vR^{(S)}\|_F^2
\leq N^{-2}(\|\vSigma\|+1)\|\vR^{(S)}\|_F^2.
\end{align*}
Then these bounds in \eqref{eq:FAbounds} follow from the condition
$\|\vR^{(S)}\|_F \leq \Psi_N$ defining $\cE(S)$.

Finally, for the bounds for $Y_i^{(S)}[d]$ and $Z_{ijk}^{(S)}[d]$ in
\eqref{eq:FAbounds}, observe that for any matrix $\vA \in \C^{n \times n}$ 
independent of $\vg_i$, we have
$\Tr(\vg_i\vg_i^\top-\vSigma)\vA \prec \|\vA\|_F$
by Assumption \ref{assump:G}(c). Then
$Y_i^{(S)}[d] \prec \|\vR^{(S)}\vA \vOmega^{(S)}[\vSigma\vOmega^{(S)}]^d\|_F
\leq \|\vR^{(S)}\vA\|_F \cdot
\|\vOmega^{(S)}\|^{d+1}\|\vSigma\|^d$, so the bound for $Y_i^{(S)}[d]$ in
\eqref{eq:FAbounds} follows
from the conditions $\|\vR^{(S)}\vA\|_F \leq \Phi_N$ and $\|\vOmega^{(S)}\| \leq
C_0$ defining $\cE(S)$.
For $Z_{ijk}^{(S)}[d]$, similarly by Assumption \ref{assump:G}(c),
\[Z_{ijk}^{(S)} \prec N^{-1}\|\vR^{(S)}\vg_j\vg_k^\top
\vR^{(S)}\vA \vOmega^{(S)}[\vSigma\vOmega^{(S)}]^d\|_F
\leq N^{-1}\|\vR^{(S)}\vg_j\| \cdot \|\vg_k^\top
\vR^{(S)}\vA\| \cdot \|\vOmega^{(S)}\|^{d+1}\|\vSigma\|^d.\]
Applying again Assumption \ref{assump:G}(c), we have
\begin{align*}
\|\vR^{(S)}\vg_j\|_2^2
&=\vg_j^\top (\vR^{(S)})^*\vR^{(S)}\vg_j
\prec |\Tr \vSigma(\vR^{(S)})^*\vR^{(S)}|
+\|(\vR^{(S)})^*\vR^{(S)}\|_F \prec \|\vR^{(S)}\|_F^2
\end{align*}
and similarly $\|\vg_k^\top \vR^{(S)}\vA\|_2^2 \prec \|\vR^{(S)}\vA\|_F^2$.
Then the bound for $Z_{ijk}^{(S)}[d]$ in \eqref{eq:FAbounds}
follows from the conditions
$\|\vR^{(S)}\vA\|_F \leq \Phi_N$, $\|\vR^{(S)}\|_F \leq \Psi_N$,
and $\|\vOmega^{(S)}\| \leq C_0$ defining $\cE(S)$.
\end{proof}

\begin{lemma}
Fix any $L,D \geq 1$. Then there exist coefficients
$\alpha(d,d',D) \in \R$ such that the following holds:
Uniformly over $S \subset [N]$
with $|S| \leq L-1$, and over $i \in S$, $j,k \in S \setminus \{i\}$,
$l \in [N] \setminus S$, and $d \leq D$,
\begin{align}
Y_i^{(S)}[d]&=\sum_{d'=d}^{d+\lceil D/2 \rceil}
\alpha(d,d',D)\Big[C_{ll}^{(Sl)}Q_l^{(Sl)}
\Big]^{d'-d}\Big(Y_i^{(Sl)}[d']-Z_{ill}^{(Sl)}[d']Q_l^{(Sl)}\Big)
+\SDcEL{N^{-D}\Psi_N^D\Phi_N}\label{eq:FAYexpand}\\
Z_{ijk}^{(S)}[d]&=\sum_{d'=d}^{d+\lceil D/2 \rceil}
\alpha(d,d',D)\Big[C_{ll}^{(Sl)}Q_l^{(Sl)}\Big]^{d'-d}
\Big(Z_{ijk}^{(Sl)}[d']-Z_{ilk}^{(Sl)}[d']B_{lj}^{(Sl)}Q_l^{(Sl)}\nonumber\\
&\quad
-Z_{ijl}^{(Sl)}[d']B_{kl}^{(Sl)}Q_l^{(Sl)}+Z_{ill}^{(Sl)}[d']
B_{lj}^{(Sl)}B_{kl}^{(Sl)}(Q_l^{(Sl)})^2\Big)
+\SDcEL{N^{-D}\Psi_N^D\Phi_N},\label{eq:FAZexpand}\\
B_{jk}^{(S)}&=B_{jk}^{(Sl)}-B_{jl}^{(Sl)}B_{lk}^{(Sl)}Q_l^{(Sl)},\label{eq:FABexpand}\\
C_{jk}^{(S)}&=C_{jk}^{(Sl)}-B_{jl}^{(Sl)}C_{lk}^{(Sl)}Q_l^{(Sl)}
-C_{jl}^{(Sl)}B_{lk}^{(Sl)}Q_l^{(Sl)}
+B_{jl}^{(Sl)}C_{ll}^{(Sl)}B_{lk}^{(Sl)}(Q_l^{(Sl)})^2,\label{eq:FACexpand}\\
Q_j^{(S)}&=\sum_{d=1}^{\lceil D/2 \rceil}
\left(Q_j^{(Sl)}\right)^d\left[(B_{jl}^{(Sl)})^2Q_l^{(Sl)}\right]^{d-1}+\SDcEL{N^{-D}\Psi_N^D}.\label{eq:FAQexpand}
\end{align}
\end{lemma}
\begin{proof}
By the Sherman-Morrison formula, on the event $\cE_L$ where $\vR^{(S)}$ and
$\vR^{(Sl)}$ both exist, we have
\begin{equation}\label{eq:FAShermanMorrison}
\vR^{(S)}=\vR^{(Sl)}-N^{-1} \vR^{(Sl)}\vg_l\vg_l^\top \vR^{(Sl)}
\cdot Q_l^{(Sl)}.
\end{equation}
Applying this to each copy of $\vR^{(S)}$ defining $B_{jk}^{(S)}$ and
$C_{jk}^{(S)}$ yields immediately \eqref{eq:FABexpand} and \eqref{eq:FACexpand},
as well as the identities
\begin{align*}
z^{-1}\vGamma+\tilde m^{(S)}\vSigma&=z^{-1}\vGamma+\Big(N^{-1}\Tr
\vR^{(S)}+(1-\gamma_N)(-1/z)\Big)\vSigma\\
&=(z^{-1}\vGamma+\tilde m^{(Sl)}\vSigma)-C_{ll}^{(Sl)}Q_l^{(Sl)}\vSigma,\\
1+B_{jj}^{(S)}&=1+B_{jj}^{(Sl)}-(B_{jl}^{(Sl)})^2Q_l^{(Sl)}.
\end{align*}
Taking inverses and applying the expansion
\[(\vA-\vDelta)^{-1}=\sum_{d=1}^{\lceil D/2 \rceil}
\vA^{-1}(\vDelta\vA^{-1})^{d-1}
+(\vA-\vDelta)^{-1}(\vDelta\vA^{-1})^{\lceil D/2 \rceil},\]
we obtain
\begin{align}
\vOmega^{(S)}&=\sum_{d=1}^{\lceil D/2 \rceil}
\vOmega^{(Sl)}[C_{ll}^{(Sl)}Q_l^{(Sl)}
\vSigma\vOmega^{(Sl)}]^{d-1}+\vE,\label{eq:Omegaexpand}\\
Q_j^{(S)}&=\sum_{d=1}^{\lceil D/2 \rceil}
Q_j^{(Sl)}[(B_{jl}^{(Sl)})^2Q_l^{(Sl)}Q_j^{(Sl)}]^{d-1}+e,\label{eq:Qexpand}
\end{align}
for remainder terms $\vE \in \C^{n \times n}$ and $e \in \C$ satisfying,
by the bounds of Lemma \ref{lemm:bounds},
\[\|\vE\|=\SDcEL{|C_{ll}^{(Sl)}|^{D/2}}
=\SDcEL{(N^{-1}\Psi)^D}, \quad |e|=\SDcEL{|(B_{jl}^{(Sl)})^2|^{D/2}}
=\SDcEL{(N^{-1}\Psi)^D}.\]
In particular,
\eqref{eq:Qexpand} shows \eqref{eq:FAQexpand}. Applying \eqref{eq:Omegaexpand}
to the definitions of $Y_i^{(S)}[d]$ and $Z_{ijk}^{(S)}[d]$, we get
\begin{align*}
Y_i^{(S)}[d]&=\Tr (\vg_i\vg_i^\top-\vSigma)\vR^{(S)}\vA
\left(\sum_{d'=1}^{\lceil D/2 \rceil}
\vOmega^{(Sl)}[C_{ll}^{(Sl)}Q_l^{(Sl)}
\vSigma\vOmega^{(Sl)}]^{d'-1}+\vE\right)\\
&\qquad\qquad \cdot \left(\vSigma\left[\sum_{d'=1}^{\lceil D/2 \rceil}
\vOmega^{(Sl)}[C_{ll}^{(Sl)}Q_l^{(Sl)}
\vSigma\vOmega^{(Sl)}]^{d'-1}+\vE\right]\right)^d,\\
Z_{ijk}^{(S)}[d]&=\frac{1}{N}\Tr (\vg_i\vg_i^\top-\vSigma)\vR^{(S)}
\vg_j\vg_k^\top \vR^{(S)}\vA
\left(\sum_{d'=1}^{\lceil D/2 \rceil}
\vOmega^{(Sl)}[C_{ll}^{(Sl)}Q_l^{(Sl)}
\vSigma\vOmega^{(Sl)}]^{d'-1}+\vE\right)\\
&\qquad\qquad\cdot
\left(\vSigma\left[\sum_{d'=1}^{\lceil D/2 \rceil}
\vOmega^{(Sl)}[C_{ll}^{(Sl)}Q_l^{(Sl)}
\vSigma\vOmega^{(Sl)}]^{d'-1}+\vE\right]\right)^d.
\end{align*}
For any matrix $\vB \in \C^{n \times n}$ independent of $\vg_i$, observe that
$\Tr (\vg_i\vg_i^\top-\vSigma)\vR^{(S)}\vA\vB=\SDcES{\Phi_N\|\vB\|}$
and $\Tr (\vg_i\vg_i^\top-\vSigma)\vR^{(S)}\vg_j\vg_k^\top \vR^{(S)}\vA\vB
=\SDcES{\Psi_N\Phi_N\|\vB\|}$
by the same arguments as those bounding $Y_i^{(S)}[d]$ and $Z_{ijk}^{(S)}[d]$
in the proof of Lemma \ref{lemm:bounds}. Then,
expanding the above and absorbing all terms containing $\vE$ 
and all terms with combined power of $C_{ll}^{(Sl)}$ larger than $D/2$
into $\SDcES{N^{-D}\Psi_N^D \Phi_N}$ remainders, we obtain for some coefficients
$\alpha(d,d',D) \in \R$ that
\begin{align*}
Y_i^{(S)}[d]&=\Tr (\vg_i\vg_i^\top-\vSigma)\vR^{(S)}\vA
\sum_{d'=0}^{\lceil D/2 \rceil} \alpha(d,d',D)[C_{ll}^{(Sl)}Q_l^{(Sl)}]^{d'}
\vOmega^{(Sl)}[\vSigma\vOmega^{(Sl)}]^{d+d'}\\
&\hspace{1in}+\SDcES{N^{-D}\Psi_N^D\Phi_N},\\
Z_{ijk}^{(S)}[d]&=\frac{1}{N}\Tr (\vg_i\vg_i^\top-\vSigma)\vR^{(S)}
\vg_j\vg_k^\top \vR^{(S)}\vA
\sum_{d'=0}^{\lceil D/2 \rceil} \alpha(d,d',D)[C_{ll}^{(Sl)}Q_l^{(Sl)}]^{d'}
\vOmega^{(Sl)}[\vSigma\vOmega^{(Sl)}]^{d+d'}\\
&\hspace{1in}+\SDcES{N^{-D}\Psi_N^D\Phi_N}.
\end{align*}
Finally, applying the Sherman-Morrison formula \eqref{eq:FAShermanMorrison}
to expand each copy of $\vR^{(S)}$, and re-indexing the summations by
$d+d' \mapsto d'$, we get \eqref{eq:FAYexpand} and \eqref{eq:FAZexpand}.
\end{proof}

\begin{lemma}\label{lemm:FAexpansion}
Fix any $L,D \geq 1$. Uniformly over $S \subset [N]$ with $|S| \leq L$ and over
$i \in S$, the following holds: Denote $\bar S=S \setminus \{i\}$. Then there exists
a collection of monomials $\cM_{i,S}$ such that
$Y_i^{(i)}[0]$ can be approximated as
\begin{align}
Y_{i}^{(i)}[0]&=\sum_{q \in \cM_{i,S}}
q\bigg(\{Y_i^{(S)}[d]\}_{d \leq \lfloor D/2 \rfloor},
\{Z_{ijk}^{(S)}[d]\}_{j,k \in \bar S, d \leq \lfloor D/2 \rfloor},
\{B_{jk}^{(S)}\}_{j \neq k \in \bar S},\nonumber\\
&\hspace{2in}\{C_{jk}^{(S)}\}_{j,k\in \bar S},\{Q_j^{(S)}\}_{j \in \bar S}\bigg)
+\SDcEL{N^{-D}\Psi_N^D\Phi_N}.\label{eq:FAexpansionrepr}
\end{align}
Each monomial $q \in \cM_{i,S}$ is a product of a real-valued scalar coefficient
and one or more factors of the form
$Y_i^{(S)}[d]$, $Z_{ijk}^{(S)}[d]$, $B_{jk}^{(S)}$ with $j \neq k$,
$C_{jk}^{(S)}$, $Q_j^{(S)}$ for $j,k \in \bar S$ and $d \leq \lfloor D/2 \rfloor$.
We have $q=\SDcEL{\Phi_N}$ uniformly over $q \in \cM_{i,S}$,
and the number of monomials $|\cM_{i,S}|$ is most a constant depending on
$L,D$. Furthermore:
\begin{enumerate}[label=(\alph*)]
\item There is exactly one factor of the form $Y_i^{(S)}[d]$ or
$Z_{ijk}^{(S)}[d]$ appearing in $q$.
\item The number of factors $Z_{ijk}^{(S)}[d]$,
$B_{jk}^{(S)}$, and $C_{jk}^{(S)}$ appearing in $q$ is no less than the number
of distinct indices of $\bar S$ (not including $i$) that appear as lower
indices across all factors of $q$.
\end{enumerate}
\end{lemma}
\begin{proof}
We arbitrarily order the indices of $\bar S=S \setminus \{i\}$ as
$l_1,l_2,\ldots,l_{|S|-1}$. Beginning with the monomial $Y_i^{(i)}[0]$,
iteratively for $j=1,2,\ldots,|S|-1$, we
replace all factors with superscript $(il_1\ldots l_{j-1})$
by a sum of terms with superscript $(il_1\ldots l_j)$,
using the recursions \eqref{eq:FAYexpand}--\eqref{eq:FAQexpand}. It is then
direct to check that
this gives a representation of the form \eqref{eq:FAexpansionrepr}, where:
\begin{itemize}
\item Each application of \eqref{eq:FAYexpand}--\eqref{eq:FAZexpand} replaces a
factor $Y_i^{(\ldots)}[d]$ or $Z_{ijk}^{(\ldots)}[d]$ by terms
having exactly one such factor. Thus, each monomial $q \in \cM_{i,S}$
has exactly one factor $Y_i^{(S)}[d]$ or $Z_{ijk}^{(S)}[d]$.
\item The number of total applications of
\eqref{eq:FAYexpand}--\eqref{eq:FAQexpand} is bounded by a constant depending on
$L,D$, so $|\cM_{i,S}|$ and the scalar coefficient of each $q \in \cM_{i,S}$
are both bounded by constants depending on $L,D$. Then, by the bounds of
\eqref{eq:FAbounds}, each $q \in \cM_{i,S}$ satisfies
$q=\SDcEL{\Phi_N}$, and the remainder in \eqref{eq:FAexpansionrepr} is at most
$\SDcEL{N^{-D}\Psi_N^D\Phi_N}$. If $q$ has the term $Y_i^{(S)}[d]$ or
$Z_i^{(S)}[d]$, then it also has combined power of
$\{C_{jk}^{(S)}\}_{j,k \in \bar S}$ equal to $d$, and hence may be
absorbed into the remainder of \eqref{eq:FAexpansionrepr} if $d>D/2$.
\item Each term on the right side of \eqref{eq:FAYexpand}--\eqref{eq:FAQexpand}
that contains the new lower index $l$ has at least one more factor of the
form $Z_{ijk}^{(\ldots)}[d]$, $B_{jk}^{(\ldots)}$, or $C_{jk}^{(\ldots)}$ than the left side. Thus, each monomial $q \in \cM_{i,S}$ is such
that the number of distinct lower indices of $\bar S$ across all of its factors
is no greater than the number of its factors of the form
$Z_{ijk}^{(\ldots)}[d]$, $B_{jk}^{(\ldots)}$, or $C_{jk}^{(\ldots)}$.
\end{itemize}
Combining these observations yields the lemma.
\end{proof}

\begin{proof-of-lemma}[\ref{lemma:fluctuationavg}]
For each $\eps,D>0$, let us fix an even integer $L=L(\eps,D)>D/\eps$.
The assumption of this lemma guarantees
$\bI\{\cE(S)^c\} \prec 0$ uniformly over $S \subset [N]$ with $|S|
\leq L$. Since the number of such subsets is at most $N^L$, we may take a union
bound (c.f.\ Proposition \ref{prop:domination}(a))
to obtain $\bI\{\cE_L^c\} \prec 0$ for the intersection event
$\cE_L$ of \eqref{eq:EL}. Noting that
$(1-\E_{\vg_i})[\vg_i^\top \vR^{(i)}\vA\vOmega\vg_i]=Y_i^{(i)}[0]$,
to prove the lemma, it suffices to show for any $\eps,D>0$ and
all sufficiently large $N$ that
\begin{equation}\label{eq:FAgoal}
\P\left[\left(\frac{1}{N}\sum_{i=1}^N 
Y_i^{(i)}[0]\right)
\bI\{\cE_L\}>\max\left(\frac{\Psi_N}{N},\frac{1}{\sqrt{N}}\right)\Phi_N
\cdot N^\eps \right]<N^{-D}.
\end{equation}

In anticipation of applying Markov's inequality, we analyze
\begin{equation}\label{eq:EYQL}
\E\left[\left(\sum_{i=1}^{N} Y_i^{(i)}[0]\right)^L
\bI\{\cE_L\}\right]
=\sum_{i_1,\ldots,i_L=1}^{N} \underbrace{\E\left[\prod_{l=1}^L Y_{i_l}^{(i_l)}
[0] \bI\{\cE_L\}\right]}_{:=\E[m(i_1,\ldots,i_L)]}.
\end{equation}
Fix any index tuple $(i_1,\ldots,i_L)$.
Letting $S=\{i_1,\ldots,i_L\}$ be the set of distinct indices in this tuple,
we apply Lemma \ref{lemm:FAexpansion} to each term
$Y_{i_l}^{(i_l)}[0]$, with this set $S$ and with $D=L$. This gives
\begin{equation}\label{eq:mexpand}
m(i_1,\ldots,i_L)=\sum_{q^{(1)} \in \cM(i_1,S)} \ldots
\sum_{q^{(l)} \in \cM(i_l,S)} \prod_{l=1}^L q^{(l)} \cdot \bI\{\cE_L\}
+\SD{(N^{-1}\Psi_N)^L\Phi_N^L},
\end{equation}
where each $\cM(i_l,S)$ is the collection of monomials arising in the
approximation of $Y_{i_l}^{(i_l)}[0]$,
and we have applied $q^{(l)}=\SDcEL{\Phi_N}$ to bound the remainder.
Observe that by \eqref{eq:FAEbounds} and Holder's inequality, we have
$\E[|m(i_1,\ldots,i_L)|^2] \leq N^C$ and
$\E[|\prod_{l=1}^L q^{(l)} \cdot \bI\{\cE_L\}|^2] \leq N^C$
for all $q^{(1)},\ldots,q^{(L)}$ and a constant $C>0$.
By this and the given condition $\Psi_N,\Phi_N \geq N^{-\upsilon}$,
we may take expectations in \eqref{eq:mexpand} using
Proposition \ref{prop:domination}(d) to get
\begin{equation}\label{eq:miLbound}
\E[m(i_1,\ldots,i_L)]=\sum_{q^{(1)} \in \cM(i_1,S)} \ldots
\sum_{q^{(l)} \in \cM(i_l,S)} \E\left[\prod_{l=1}^L q^{(l)} \cdot
\bI\{\cE_L\}\right]+\SD{(N^{-1}\Psi_N)^L\Phi_N^L}.
\end{equation}

Now to bound $\E[\prod_{l=1}^L q^{(l)} \cdot \bI\{\cE_L\}]$,
we consider separately two cases, focusing on those indices $i_l$ which
appear exactly once in $(i_1,\ldots,i_L)$.
In the first case, suppose there is some such index $i_l$ that does not appear
as a lower index of $q^{(l')}$ for any $l' \neq l$. Fixing this set
$S=\{i_1,\ldots,i_L\}$ and index $i_l \in S$, let us introduce
\begin{align*}
\cE'&=\bigg\{\vR^{(S)} \text{ exists, } \|\vR^{(S)}\vA\|_F \leq \Phi_N,\;
\|\vR^{(S)}\|_F \leq \Psi_N,\;
\|(z^{-1}\vGamma+\tilde m^{(S)} \vSigma)^{-1}\| \leq C_0,\\
&\hspace{2in} \text{ and }
|1+N^{-1} \vg_j^\top \vR^{(S)}\vg_j| \geq c_0 \text{ for all } j \in S \setminus
\{i_l\} \bigg\}.
\end{align*}
Comparing
with the definition of $\cE(S)$ from \eqref{eq:ESdef}, observe that
only the last condition defining $\cE'$ is different (where we do not require
the bound for $j=i_l$), so that this event $\cE'$ is independent of $\vg_{i_l}$.
Then $\cE_L \subseteq \cE(S) \subseteq \cE'$, and
\begin{equation}\label{eq:eventdecomp}
\E\left[\prod_{l=1}^L q^{(l)} \cdot \bI\{\cE_L\}\right]
=\E\left[\prod_{l=1}^L q^{(l)} \cdot \bI\{\cE'\}\right]
-\E\left[\prod_{l=1}^L q^{(l)} \cdot \bI\{\cE'\}\bI\{\cE_L^c\}
\right].
\end{equation}
For the first term of \eqref{eq:eventdecomp}, observe that
both $\{q^{(l')}:l' \neq l\}$ and $\cE'$ are independent of
$\vg_{i_l}$, and only the one
factor $Y_{i_l}^{(S)}[d]$ or $Z_{i_ljk}^{(S)}[d]$ in $q^{(l)}$ depends on
$\vg_{i_l}$. Then, noting that $\E_{\vg_i}[Y_i^{(S)}[d]]=0$ and
$\E_{\vg_i}[Z_{ijk}^{(S)}[d]]=0$, the first term of \eqref{eq:eventdecomp}
is 0. For the second term of \eqref{eq:eventdecomp},
observe that all statements of \eqref{eq:FAEbounds} continue to hold with
$\cE(S)$ replaced by $\cE'$, except for the bound on $Q_{i_l}^{(S)}$.
But $Q_{i_l}^{(S)}$ appears neither in
$\{q^{(l')}:l' \neq l\}$ nor in $q^{(l)}$, so we may apply Holder's inequality
to get $\E[|\prod_{l=1}^L q^{(l)}|^2 \bI\{\cE'\}] \leq N^C$ for a constant
$C>0$. Then, applying Cauchy-Schwarz and $\bI\{\cE_L^c\} \prec 0$, the second
term of \eqref{eq:eventdecomp} is bounded by $N^{-D'}$ for any fixed constant
$D'>0$ and all large $N$. Thus,
\begin{equation}\label{eq:FAcase1}
\E\left[\prod_{l=1}^L q^{(l)} \cdot \bI\{\cE_L\}\right] \leq N^{-D'}.
\end{equation}

In the second case,
every index $i_l$ that appears exactly once in $(i_1,\ldots,i_L)$ appears
as a lower index of $q^{(l')}$ for some $l' \neq l$. Call the number
of such indices $K$. Then condition (b) of Lemma \ref{lemm:FAexpansion} implies
that the total number of factors of the forms
$Z_{ijk}^{(S)}[d]$, $B_{jk}^{(S)}$ for $j \neq k$, and $C_{jk}^{(S)}$
across all monomials $q^{(1)},\ldots,q^{(L)}$ is at least $K$. Then, 
by the bounds of Lemma \ref{lemm:bounds} and
Proposition \ref{prop:domination}(d), we have
\begin{equation}\label{eq:FAcase2}
\E\left[\prod_{l=1}^L q^{(l)} \cdot \bI\{\cE_L\}\right]
\prec (N^{-1}\Psi_N)^K\Phi_N^L.
\end{equation}

Under the given condition $\Phi_N,\Psi_N \geq N^{-\upsilon}$, we have
$N^{-D'} \leq (N^{-1}\Psi_N)^K\Phi_N^L$ for large enough $D'$. Then,
combining the two cases \eqref{eq:FAcase1} and \eqref{eq:FAcase2} and applying
this back to \eqref{eq:miLbound}, we get
\begin{equation}\label{eq:miLboundfinal}
\E[m(i_1,\ldots,i_L)] \prec (N^{-1}\Psi_N)^K\Phi_N^L
\end{equation}
where $K$ is the number of indices in $S=\{i_1,\ldots,i_L\}$ that appear exactly
once in $(i_1,\ldots,i_L)$.
Let $J$ be the number of distinct indices in $S=\{i_1,\ldots,i_L\}$ that appear
at least twice in $(i_1,\ldots,i_L)$. Then $2J+K \leq L$, and the number of
index tuples $(i_1,\ldots,i_L) \in [N]^L$ with these values of $(J,K)$ is at
most $CN^{J+K}$, for a constant $C=C(J,K)>0$. Then,
applying \eqref{eq:miLboundfinal} back to \eqref{eq:EYQL} yields
\begin{align*}
\E\left[\left(\sum_{i=1}^{N} Y_i^{(i)}[0]\right)^L
\bI\{\cE_L\}\right] &\prec \max_{J,K \geq 0:\,2J+K \leq L}
N^{J+K} \cdot (N^{-1}\Psi_N)^K\Phi_N^L\\
&=\max_{J,K \geq 0:\,2J+K \leq L}
(\sqrt{N})^{2J}\Psi_N^K\Phi_N^L \leq \max(\Psi_N,\sqrt{N})^L \Phi_N^L.
\end{align*}
Finally, by Markov's inequality, the probability in \eqref{eq:FAgoal} is at most
\begin{align*}
\max(\Psi_N,\sqrt{N})^{-L}\Phi_N^{-L}N^{-\eps L}
\cdot \E\left[\left(\sum_{i=1}^{N} Y_i^{(i)}[0]\right)^L
\bI\{\cE_L\}\right] \prec N^{-\eps L},
\end{align*}
and \eqref{eq:FAgoal} follows as desired under our initial
choice $L=L(\eps,D)>D/\eps$.
\end{proof-of-lemma}

\subsection{No eigenvalues outside the support}\label{sec:nooutliers}

We now prove Theorem \ref{thm:no_outlier}.
Let $m_N(z),\tilde m_N(z)$ be the Stieltjes
transform of the $N$-dependent deterministic measures $\mu_N,\tilde \mu_N$.
For each $z \in \C^+$, $\tilde m_N(z)$ is the unique root in $\C^+$ to the equation
\begin{align}
z&=-\frac{1}{\tilde m_N(z)}+\gamma_N \int \frac{\lambda}{1+\lambda \tilde
m_N(z)} d\nu_N(\lambda),\label{eq:tildemn}
\end{align}
and $m_N(z),\tilde m_N(z)$ are related by
$\tilde m_N(z)=\gamma_N m_N(z)+(1-\gamma_N)({-}1/z)$.
Define the discrete set
\begin{equation}\label{eq:zNdomain}
\cT_N=\{0\} \cup \{-1/\lambda:\lambda \in \supp{\nu_N}\}.
\end{equation}
On the domain $\C \setminus \cT_N$,
we may define the formal inverse of (\ref{eq:tildemn}),
\begin{align}
z_N(\tilde m)&=-\frac{1}{\tilde m}+\gamma_N \int \frac{\lambda}{1+\lambda \tilde
m} d\nu_N(\lambda),\label{eq:z_N_def}
\end{align}
which is a finite-$N$ analogue of (\ref{eq:inv_z}).
%Here $z_N$ is a meromorphic function with poles at the points of $\cT_N$.
Let $\cS_N$ be the deterministic support defined in \eqref{eq:support}, and let
$U_N(\eps)$ be the spectral domain \eqref{eq:U_e}. The following basic
properties of $\cS_N$ and $\tilde m_N(z)$ are known.

\begin{proposition}\label{prop:basic_mn}
Suppose Assumption \ref{assump:G}(a) holds, and fix any $\eps>0$. Then there
exist constants $C_0,c_0>0$, depending only on $\eps$ and the constants $C,c$ of
Assumption \ref{assump:G}(a), such that for all $x \in \cS_N$ we have $|x| \leq C$,
and for all $z=x+i\eta \in U_N(\eps)$ we have
\[c<|\tilde m_N(z)|<C, \qquad c\eta \leq |\Im \tilde m_N(z)| \leq C\eta,
\qquad \min_{\lambda \in \supp{\nu_N}}
|1+\lambda\,\tilde m_N(z)| \geq c\]
\end{proposition}
\begin{proof}
See \cite[Propositions A.3, B.1, B.2]{fan2022tracy}.
\end{proof}

Let $m_{\tilde \vK}(z)=N^{-1}\Tr (\widetilde \vK-z\vI)^{-1}$ be the Stieltjes
transform of the empirical eigenvalue distribution of $\widetilde \vK=\vG\vG^\top$.
Since $\widetilde \vK$ and $\vK=\vG^\top \vG$ have the same
eigenvalues up to $|N-n|$ 0's, we have
\begin{equation}\label{eq:m_tildeK}
m_{\tilde \vK}(z)=\gamma_N m_{\vK}(z)+(1-\gamma_N)(-1/z),
\end{equation}
so in particular $m_{\tilde \vK}$ coincides with $\tilde m_{\vK}^{(\emptyset)}$
from \eqref{eq:leaveoneout}.
We begin with a preliminary estimate for the Stieltjes transform
$m_{\tilde \vK}(z)$ when $\Im z \geq N^{-1/11}$. Similar statements have been shown
in \citep{silverstein1995strong,bai1998no}, and we provide an argument
here following ideas of \cite[Section 3]{bai1998no} for later reference.

\begin{lemma}\label{lemma:prelimestimate}
Fix any $\eps>0$, and suppose Assumption \ref{assump:G} holds.
Then, uniformly over $z=x+i\eta \in U_N(\eps)$ with $\Im z \geq N^{-1/11}$,
\[m_{\tilde \vK}(z)-\tilde m_N(z) \prec \frac{1}{\sqrt{N}\eta^4}.\]
\end{lemma}
\begin{proof}
Let $\vR^{(i)}$ and $\tilde m_{\vK}^{(i)}$ be as defined in
\eqref{eq:leaveoneout} with $\vGamma=z\vI$.
Applying the Sherman-Morrison formula
\begin{equation}\label{eq:ShermanMorrison}
\vR=\vR^{(i)}-\frac{N^{-1} \vR^{(i)}\vg_i\vg_i^\top \vR^{(i)}}
{1+N^{-1}\vg_i^\top \vR^{(i)}\vg_i},
\end{equation}
for any matrix $\vB \in \C^{n \times n}$ we have
\begin{align}
\Tr \vB=\Tr (\vK-z\vI)\vR\vB
&=-z\Tr \vR\vB+\frac{1}{N}\sum_{i=1}^N \vg_i^\top \vR\vB\vg_i\nonumber\\
&=-z\Tr \vR\vB+\frac{1}{N}\sum_{i=1}^N \frac{\vg_i^\top \vR^{(i)}\vB\vg_i}
{1+N^{-1}\vg_i^\top \vR^{(i)}\vg_i}.\label{eq:fundamentalidentity}
\end{align}

Choosing $\vB=\vI$ in \eqref{eq:fundamentalidentity},
applying $\Tr \vR=n\,m_{\vK} =N\,m_{\tilde \vK}+(n-N)(-1/z)$, and rearranging,
we obtain the identity
\begin{equation}\label{eq:mnidentity}
m_{\tilde \vK}=-\frac{1}{Nz}\sum_{i=1}^N \frac{1}{1+N^{-1} \vg_i^\top
\vR^{(i)}\vg_i}.
\end{equation}
Now fix any deterministic matrix $\vA \in \C^{n \times n}$, define
\[d_i=\frac{1}{N}\,\vg_i^\top \vR^{(i)}\vA
(\vI+m_{\tilde \vK}\vSigma)^{-1}\vg_i
-\frac{1}{N}\Tr \vR\vA(\vI+m_{\tilde \vK}\vSigma)^{-1}\vSigma,\]
and choose $\vB=\vA(\vI+m_{\tilde \vK}\vSigma)^{-1}$
in \eqref{eq:fundamentalidentity}. Then, applying also the identity
\eqref{eq:mnidentity}, we get
\begin{align}
&\Tr \vA(\vI+m_{\tilde \vK}\vSigma)^{-1}\nonumber\\
&=-z\Tr \vR\vA(\vI+m_{\tilde \vK}\vSigma)^{-1}
-zm_{\tilde \vK} \Tr \vR\vA(\vI+m_{\tilde \vK}\vSigma)^{-1}\vSigma
+\sum_{i=1}^N \frac{d_i}{1+N^{-1}\vg_i^\top \vR^{(i)}\vg_i}\nonumber\\
&=-z \Tr \vR\vA+\sum_{i=1}^N \frac{d_i}{1+N^{-1}\vg_i^\top \vR^{(i)}\vg_i}.
\label{eq:detequividentity}
\end{align}

We proceed to bound $d_i$, where (for later purposes)
we derive estimates in terms of the Frobenius norms of
$\vR,\vR\vA,\vR^{(i)},\vR^{(i)}\vA$ rather than their operator norms.
Note that Assumption \ref{assump:G}(c) implies, for any matrix
$\vB \in \C^{n \times n}$ independent of $\vg_i$,
\begin{equation}\label{eq:Agbound}
\|\vB\vg_i\|^2=\vg_i^\top \vB^*\vB \vg_i \prec
\Tr \vSigma\vB^*\vB+\|\vB^*\vB\|_F \prec \|\vB\|_F^2.
\end{equation}
We have also, by Assumption \ref{assump:G}(c) and the Sherman-Morrison
formula \eqref{eq:ShermanMorrison},
\begin{align}
N^{-1}|\Tr \vR\vB-\Tr \vR^{(i)}\vB|
&=N^{-2}|1+N^{-1}\vg_i^\top \vR^{(i)}\vg_i|^{-1}
|\vg_i^\top \vR^{(i)}\vB\vR^{(i)}\vg_i|\nonumber\\
&\prec N^{-2}|1+N^{-1}\vg_i^\top \vR^{(i)}\vg_i|^{-1}
\Big(|\Tr \vSigma\vR^{(i)}\vB\vR^{(i)}|+\|\vR^{(i)}\vB\vR^{(i)}\|_F\Big)\nonumber\\
&\prec N^{-2}|1+N^{-1}\vg_i^\top \vR^{(i)}\vg_i|^{-1}\|\vR^{(i)}\vB\|_F 
\|\vR^{(i)}\|_F.
\label{eq:mdiffestimate}
\end{align}

Define
$d_i=d_{i,1}+d_{i,2}+d_{i,3}+d_{i,4}$ where
\begin{equation}\label{eq:di14}
\begin{aligned}
d_{i,1}&=N^{-1}\vg_i^\top \vR^{(i)}\vA(\vI+m_{\tilde \vK}\vSigma)^{-1}
\vg_i-N^{-1}\vg_i^\top \vR^{(i)}\vA(\vI+\tilde m_{\vK}^{(i)}\vSigma)^{-1}\vg_i,\\
d_{i,2}&=N^{-1}\vg_i^\top \vR^{(i)}\vA(\vI+\tilde m_{\vK}^{(i)}\vSigma)^{-1}
\vg_i-N^{-1}\Tr \vSigma \vR^{(i)}\vA(\vI+\tilde m_{\vK}^{(i)}\vSigma)^{-1},\\
d_{i,3}&=N^{-1}\Tr \vSigma \vR^{(i)}\vA(\vI+\tilde m_{\vK}^{(i)}\vSigma)^{-1}
-N^{-1}\Tr \vSigma \vR\vA(\vI+\tilde m_{\vK}^{(i)}\vSigma)^{-1},\\
d_{i,4}&=N^{-1}\Tr \vSigma \vR\vA(\vI+\tilde m_{\vK}^{(i)}\vSigma)^{-1}
-N^{-1}\Tr \vSigma \vR\vA(\vI+m_{\tilde \vK}\vSigma)^{-1}.
\end{aligned}
\end{equation}
Applying the identity $\vA^{-1}-\vB^{-1}=\vA^{-1}(\vB-\vA)\vB^{-1}$,
the definition of $\tilde m_{\vK}^{(i)}$ in \eqref{eq:leaveoneout},
and the bounds \eqref{eq:Agbound} and \eqref{eq:mdiffestimate} (the latter with
$\vB=\vI$),
\begin{align}
|d_{i,1}| &\leq N^{-1} \|\vg_i^\top \vR^{(i)}\vA\|
\|(\vI+m_{\tilde \vK}\vSigma)^{-1}\| \|(\tilde m_{\vK}^{(i)}-m_{\tilde
\vK})\vSigma\| \|(\vI+\tilde m_{\vK}^{(i)}\vSigma)^{-1}\| \|\vg_i\|\nonumber\\
&\prec N^{-5/2} |1+N^{-1}\vg_i^\top \vR^{(i)}\vg_i|^{-1} 
\|\vR^{(i)}\vA\|_F \|\vR^{(i)}\|_F^2 \|(\vI+m_{\tilde \vK}\vSigma)^{-1}\|
\|(\vI+\tilde m_{\vK}^{(i)}\vSigma)^{-1}\|.\label{eq:di1}
\end{align}
Applying Assumption \ref{assump:G}(c),
\begin{align}
|d_{i,2}| &\prec N^{-1} \|\vR^{(i)}\vA(\vI+\tilde m_{\vK}^{(i)}\vSigma)^{-1}\|_F
\leq N^{-1} \|\vR^{(i)}\vA\|_F \|(\vI+\tilde m_{\vK}^{(i)}\vSigma)^{-1}\|.\label{eq:di2}
\end{align}
Applying the Sherman-Morrison identity \eqref{eq:ShermanMorrison},
$|\Tr \vu\vv^\top| \leq \|\vu\| \|\vv\|$, and \eqref{eq:Agbound},
\begin{align}
|d_{i,3}| & \leq N^{-2} |1+N^{-1}\vg_i^\top \vR^{(i)}\vg_i|^{-1}
\|\vSigma \vR^{(i)}\vg_i\| \|\vg_i^\top \vA
\vR^{(i)}(\vI+\tilde m_{\vK}^{(i)}\vSigma)^{-1}\|\nonumber\\
&\prec N^{-2} |1+N^{-1}\vg_i^\top \vR^{(i)} \vg_i|^{-1}
\|\vR^{(i)}\vA\|_F \|\vR^{(i)}\|_F \|(\vI+\tilde m_{\vK}^{(i)}\vSigma)^{-1}\|.\label{eq:di3}
\end{align}
Finally, applying $\vA^{-1}-\vB^{-1}=\vA^{-1}(\vB-\vA)\vB^{-1}$,
\eqref{eq:mdiffestimate} (with $\vB=\vI$), and 
$|\Tr \vA\vB| \leq \|\vA\|_F\|\vB\|_F \leq \sqrt{N} \|\vA\|_F\|\vB\|$,
\begin{align}
|d_{i,4}| &=N^{-1} \Big|\Tr \vSigma \vR\vA
(\vI+\tilde m_{\vK}^{(i)}\vSigma)^{-1} (\tilde m_{\vK}^{(i)}-m_{\tilde \vK})\vSigma
(\vI+m_{\tilde \vK}\vSigma)^{-1}\Big|\nonumber\\
&\prec N^{-5/2} |1+N^{-1}\vg_i^\top \vR^{(i)}\vg_i|^{-1} 
\|\vR\vA\|_F \|\vR\|_F^2
\|(\vI+\tilde m_{\vK}^{(i)}\vSigma)^{-1}\| 
\|(\vI+m_{\tilde \vK}\vSigma)^{-1}\|.\label{eq:di4}
\end{align}

For the current proof, we apply \eqref{eq:detequividentity} and the definitions
\eqref{eq:di14} with $\vA=\vI$.
Recalling $\Tr \vR=n\,m_{\vK}=N\,m_{\tilde \vK}+(n-N)(-1/z)$ and
rearranging \eqref{eq:detequividentity} with $\vA=\vI$, we get the identity
\begin{equation}\label{eq:Deltaidentity}
z_N(m_{\tilde \vK})-z={-}\frac{1}{m_{\tilde \vK}}
\cdot \frac{1}{N}\sum_{i=1}^N \frac{d_i}{1+N^{-1}\vg_i^\top
\vR^{(i)}\vg_i}
\end{equation}
where $z_N(m)=-(1/m)+N^{-1}\Tr \vSigma(\vI+m\vSigma)^{-1}$ is the function
defined in \eqref{eq:z_N_def}. For any $z=x+i\eta$ with $\eta>0$, we have
\begin{align}
|z(1+N^{-1}\vg_i^\top \vR^{(i)}\vg_i)|
&\geq \Im [z(1+N^{-1}\vg_i^\top \vR^{(i)}\vg_i)] \geq \Im z=\eta,
\label{eq:prelimbound1}\\
\max(\|\vR\|_F,\|\vR^{(i)}\|_F) &\leq N^{1/2}\max(\|\vR\|,\|\vR^{(i)}\|)
\leq N^{1/2}\eta^{-1}.\label{eq:prelimbound2}
\end{align}
Here, the second inequalities of both \eqref{eq:prelimbound1} and
\eqref{eq:prelimbound2} follow from the spectral representations of
$\vR,\vR^{(i)}$, i.e.\ writing $(\lambda_j,\vv_j)_{j=1}^n$ for the
eigenvalues and unit eigenvectors of $\vK^{(i)}$, we have
\begin{align*}
\Im [z\vg_i^\top \vR^{(i)}\vg_i]
&=\Im \left[z \vg_i^\top\left(\sum_{j=1}^n
\frac{1}{\lambda_j-z}\vv_j\vv_j^\top\right)
\vg_i\right]
=\sum_{j=1}^n \Im \frac{z}{\lambda_j-z} \cdot (\vg_i^\top \vv_j)^2\\
&=\sum_{j=1}^n \frac{\lambda_j \Im z}{|\lambda_j-z|^2} \cdot (\vg_i^\top \vv_j)^2
\geq 0,\\
\|\vR^{(i)}\|&=\left\|\sum_{j=1}^n
\frac{1}{\lambda_j-z}\vv_j\vv_j^\top\right\|=\max_{j=1}^n |\lambda_j-z|^{-1}
\leq \eta^{-1},
\end{align*}
and similarly for $\|\vR\|$. In particular, \eqref{eq:prelimbound1} and
\eqref{eq:prelimbound2} imply
\begin{equation}\label{eq:prelimbounds12}
(1+N^{-1}\vg_i^\top \vR^{(i)}\vg_i)^{-1} \prec \eta^{-1},
\qquad \|\vR\|_F,\|\vR^{(i)}\|_F \prec N^{1/2}\eta^{-1}.
\end{equation}
Next, observe that if $m(z)=\int
\frac{1}{\lambda-z}d\mu(\lambda)$ is the Stieltjes transform of any probability
measure $\mu$ supported on $[-B,B]$, then for $z=x+i\eta$ with $\eta>0$ and
$|z| \leq \eps^{-1}$, we have
\[\Im m(z)=\int \frac{\eta}{|\lambda-z|^2}d\mu(\lambda) \geq c\eta, \quad
|\Re m(z)| \leq \int \frac{|\lambda-x|}{|\lambda-z|^2}d\mu(\lambda)
\leq (C/\eta)\Im m(z)\]
for some constants $C,c>0$ depending on $\eps,B$.
Consequently, for any $\lambda \geq 0$, either
$\lambda \cdot |\Re m(z)|<1/2$ or $\lambda \cdot \Im m(z) \geq 2\eta/C$,
so $|1+\lambda m(z)| \geq \max(2,2\eta/C)$. By Assumption \ref{assump:G}(b)
and Weyl's inequality, we have $\bI\{\|\vK\|>B\} \prec 0$ and
$\bI\{\|\vK^{(i)}\|>B\} \prec 0$, and on the event where $\|\vK\|,\|\vK^{(i)}\|
\leq B$, we have that $m_{\tilde \vK},\tilde m_{\vK}^{(i)}$ are Stieltjes
transforms of probability measures supported on $[-B,B]$. Thus, this implies
\begin{equation}\label{eq:prelimbound3}
|m_{\tilde \vK}|^{-1} \leq |\Im m_{\tilde \vK}|^{-1} \prec \eta^{-1},
\quad \max(\|(\vI+m_{\tilde \vK}\vSigma)^{-1}\|,
\|(\vI+\tilde m_{\vK}^{(i)}\vSigma)^{-1}\|) \prec \eta^{-1}.
\end{equation}
Applying these bounds \eqref{eq:prelimbounds12} and
\eqref{eq:prelimbound3} to \eqref{eq:di1}--\eqref{eq:di4},
we get $d_i \prec N^{-1}\eta^{-6}+N^{-1/2}\eta^{-2} \leq 2N^{-1/2}\eta^{-2}$ for
$\eta \geq N^{-1/11}$. Then, applying these bounds 
\eqref{eq:prelimbounds12} and \eqref{eq:prelimbound3} also to
\eqref{eq:Deltaidentity}, we get
\begin{equation}\label{eq:zNestimateprelim}
z_N(m_{\tilde \vK})-z \prec \frac{1}{\sqrt{N}\eta^4}.
\end{equation}

The proof is completed by the following stability argument: When
$\eta \geq N^{-1/11}$, we have $1/(\sqrt{N} \eta^4) \ll \eta=\Im z$,
so \eqref{eq:zNestimateprelim} implies in particular that
\begin{equation}\label{eq:stabilityevent}
\bI\{z_N(m_{\tilde \vK}) \notin \C^+\} \prec 0.
\end{equation}
On the event $z_N(m_{\tilde \vK}) \in \C^+$, recalling the implicit definition
of $\tilde m_N:\C^+ \to \C^+$ by \eqref{eq:tildemn}, 
the value $\tilde m_N(z_N(m_{\tilde \vK}))$ must be the unique root
$u \in \C^+$ to the equation
\[z_N(m_{\tilde \vK})=-\frac{1}{u}+\gamma_N \int \frac{\lambda}{1+\lambda
u}d\nu_N(\lambda),\]
i.e.\ to the equation $z_N(m_{\tilde \vK})=z_N(u)$. This equation is
satisfied by $u=m_{\tilde \vK} \in \C^+$, so we deduce that
$\tilde m_N(z_N(m_{\tilde \vK}))=m_{\tilde \vK}$.
Then, applying that $z \in U_N(\eps)$ and that $\tilde m_N:\C^+ \to \C^+$
is $(4/\eps^2)$-Lipschitz over the domain $U_N(\eps/2)$,
we obtain from \eqref{eq:zNestimateprelim} that
\[\bI\{z_N(m_{\tilde \vK}) \in \C^+\}\Big(m_{\tilde \vK}-\tilde m_N(z)\Big)
=\bI\{z_N(m_{\tilde \vK}) \in \C^+\}\Big(\tilde m_N(z_N(m_{\tilde \vK}))-\tilde
m_N(z)\Big) \prec \frac{1}{\sqrt{N} \eta^4}.\]
Together with \eqref{eq:stabilityevent}, this yields the lemma.
\end{proof}

\begin{coro}\label{coro:Fnormbound}
Fix any $\eps>0$, and suppose Assumption \ref{assump:G} holds.
Then there is a constant $C>0$ such that uniformly over $z \in U_N(\eps)$
with $\Im z \geq N^{-1/11}$,
\[\bI\{\|\vR(z)\|_F>C\sqrt{N}\} \prec 0.\]
\end{coro}
\begin{proof}
Since $m_{\tilde \vK}(z)=\gamma_N m_{\vK}(z)+(1-\gamma_N)(-1/z)$ and
$\tilde m_N(z)=\gamma_N m_N(z)+(1-\gamma_N)(-1/z)$, Lemma
\ref{lemma:prelimestimate} implies also
\[m_{\vK}(z)-m_N(z) \prec \frac{1}{\sqrt{N}\eta^4} \ll \eta.\]
Observe that
$\Im m_N(z)=\int \eta/|\lambda-z|^2 d\mu_N(\lambda) \leq \eta \eps^{-2}$
for $z \in U_N(\eps)$, so
$\bI\{\Im m_{\vK}(z)>(1+\eps^{-2})\eta\} \prec 0$.
Then by the identity $\|\vR(z)\|_F^2=
\sum_i 1/|z-\lambda_i(\vK)|^2=(n/\eta)\Im m_{\vK}(z)$, we get
$\bI\{\|\vR(z)\|_F>C\sqrt{N}\} \prec 0$ for a constant $C=C(\eps)>0$, as desired.
\end{proof}

We may now apply Corollary \ref{coro:Fnormbound} and the fluctuation averaging
result of Lemma \ref{lemma:fluctuationavg} 
to improve the estimate of Lemma \ref{lemma:prelimestimate} to the following
result.

\begin{lemma}\label{lemma:improvedestimate}
Fix any $\eps>0$, and suppose Assumption \ref{assump:G} holds.
Then, uniformly over $z=x+i\eta \in U_N(\eps)$ with $\Im z \geq N^{-1/11}$,
\[m_{\tilde \vK}(z)-\tilde m_N(z) \prec \frac{1}{N}.\]
\end{lemma}
\begin{proof}
We derive an improved estimate for \eqref{eq:Deltaidentity}. First,
combining Lemma \ref{lemma:prelimestimate} with the bounds for $\tilde m_N(z)$
in Proposition \ref{prop:basic_mn}, there are constants $C_0,c_0>0$ for which
\begin{equation}\label{eq:tildembounds}
\bI\{|m_{\tilde \vK}|>C_0\} \prec 0,
\qquad \bI\{|m_{\tilde \vK}|<c_0\} \prec 0,
\qquad \bI\{\|(\vI+m_{\tilde \vK}\vSigma)^{-1}\|>C_0\} \prec 0
\end{equation}
uniformly over $z \in U_N(\eps)$ with $\Im z \geq N^{-1/11}$. Next, applying
Assumption \ref{assump:G}(c), we have also uniformly over $i \in [N]$,
\begin{align*}
N^{-1}\vg_i^\top \vR^{(i)}\vg_i
&=N^{-1}\Tr \vSigma \vR^{(i)}+\SD{N^{-1}\|\vR^{(i)}\|_F}\\
&=N^{-1}\Tr \vSigma \vR+\SD{N^{-2}|1+N^{-1}\vg_i^\top \vR^{(i)}\vg_i|^{-1}
\|\vR^{(i)}\|_F^2}+\SD{N^{-1}\|\vR^{(i)}\|_F}
\end{align*}
where the second line follows from \eqref{eq:mdiffestimate} applied with
$\vB=\vSigma$. Applying $\|\vR^{(i)}\|_F \prec N^{1/2}$ by Corollary
\ref{coro:Fnormbound} and the estimate $|1+N^{-1}\vg_i^\top \vR^{(i)}\vg_i|^{-1}
\prec \eta^{-1}$ from \eqref{eq:prelimbounds12}, this gives
\begin{equation}\label{eq:1plusgRgbound}
1+N^{-1}\vg_i^\top \vR^{(i)}\vg_i=1+N^{-1}\Tr \vSigma\vR+\SD{N^{-1/2}}.
\end{equation}
Then, applying this and 
$|1+N^{-1}\vg_i^\top \vR^{(i)}\vg_i|^{-1} \prec \eta^{-1}$ 
to \eqref{eq:mnidentity},
\[m_{\tilde \vK}=-\frac{1}{z} \cdot \frac{1}{1+N^{-1}\Tr \vSigma\vR}
+\SD{N^{-1/2}\eta^{-2}}.\]
Together with the first bound of \eqref{eq:tildembounds} and the bound $|z| \leq
\eps^{-1}$ for $z \in U_N(\eps)$, this implies
for a constant $c_0>0$ that $\bI\{|1+N^{-1}\Tr \vSigma\vR|<c_0\} \prec 0$,
and thus $\bI\{|1+N^{-1}\vg_i^\top \vR^{(i)}\vg_i|<c_0\} \prec 0$.

Applying Corollary \ref{coro:Fnormbound} and the above arguments now
for $\vK^{(S)}$ and $\vR^{(S)}$ in place of $\vK$
and $\vR$, we obtain for any fixed $L \geq 1$ and
some constants $C_0,c_0>0$, uniformly over $S \subset [N]$ with $|S| \leq L$,
over $i \in S$, and over $z \in U_N(\eps)$ with $\Im z \geq N^{-1/11}$,
\begin{equation}\label{eq:checkFAconditions}
\begin{gathered}
\bI\{|\tilde m_{\vK}^{(S)}|>C_0\} \prec 0,
\quad \bI\{|\tilde m_{\vK}^{(S)}|<c_0\} \prec 0,
\quad \bI\{\|(\vI+\tilde m_{\vK}^{(S)}\vSigma)^{-1}\|>C_0\} \prec 0,\\
\bI\{\|\vR^{(S)}\|_F>C\sqrt{N}\} \prec 0,
\quad \bI\{|1+N^{-1}\Tr \vSigma\vR^{(S)}|<c_0\} \prec 0,\\
\bI\{|1+N^{-1}\vg_i^\top \vR^{(S)}\vg_i|<c_0\} \prec 0.
\end{gathered}
\end{equation}
(We remark that a direct application of the above arguments for $\vK^{(S)}$
yields the first three estimates of \eqref{eq:checkFAconditions} for the
quantity $\frac{N}{N-|S|} \tilde m_{\vK}^{(S)}=\frac{1}{N-|S|}\Tr
\vR^{(S)}+\frac{n}{N-|S|}(-1/z)$ in place of $\tilde m_{\vK}^{(S)}$,
and the estimates for $\tilde m_{\vK}^{(S)}$ then follow for slightly modified
constants $C_0,c_0>0$ because $|S| \leq L$.)

Finally, applying \eqref{eq:1plusgRgbound} and \eqref{eq:checkFAconditions}
back to \eqref{eq:Deltaidentity} and \eqref{eq:di1}--\eqref{eq:di4} with
$\vA=\vI$, we get $|d_{i,1}|,|d_{i,3}|,|d_{i,4}| \prec N^{-1}$,
$|d_{i,2}| \prec N^{-1/2}$, and
\begin{align*}
|z_N(m_{\tilde \vK})-z| &\prec \left|\frac{1}{N}\sum_{i=1}^N \frac{d_{i,2}}
{1+N^{-1}\vg_i^\top \vR^{(i)}\vg_i}\right|+\SD{N^{-1}}\\
&=\frac{1}{N} \cdot \frac{1}{1+N^{-1}\Tr \vSigma\vR}
\cdot \left|\sum_{i=1}^N d_{i,2}\right|+\SD{N^{-1}}.
\end{align*}
The statements of \eqref{eq:checkFAconditions} verify the needed assumptions
of Lemma \ref{lemma:fluctuationavg} with $\vA=\vI$, $\vGamma=z\vI$, and
$\Phi_N=\Psi_N=C\sqrt{N}$. Then Lemma \ref{lemma:fluctuationavg}
gives $\sum_{i=1}^N d_{i,2} \prec 1$, and hence
\[|z_N(m_{\tilde \vK})-z| \prec N^{-1}.\]
The proof is then completed by the same stability argument as in the conclusion
of the proof of Lemma \ref{lemma:prelimestimate}.
\end{proof}

\begin{proof-of-theorem}[\ref{thm:no_outlier}]
We apply the idea of \cite[Section 6]{bai1998no}. Let $z=x+i\eta$, where
$\dist{x,\cS_N} \geq \eps$ and $\eta=N^{-1/11}$.
Taking imaginary part in the estimate $m_{\tilde \vK}(z)-\tilde m_N(z) \prec
N^{-1}$ of Lemma \ref{lemma:improvedestimate} and multiplying by
$\eta$ gives
\[\frac{1}{N}\sum_{j=1}^N \frac{\eta^2}{(\lambda_j(\widetilde \vK)-x)^2+\eta^2}
-\int \frac{\eta^2}{(\lambda-x)^2+\eta^2} d\tilde\mu_N(\lambda) \prec
\frac{\eta}{N}.\]
Fix any integer $P \geq 1$, and apply this instead
at the point $z=x+i\sqrt{p}\,\eta$ for each $p=1,\ldots,P$. Then
\[\frac{1}{N}\sum_{j=1}^N \frac{\eta^2}{(\lambda_j(\widetilde \vK)-x)^2+p\eta^2}
-\int \frac{\eta^2}{(\lambda-x)^2+p\eta^2} d\tilde\mu_N(\lambda) \prec
\frac{\eta}{N} \text{ for all } p=1,\ldots,P.\]
Taking successive finite differences using
\[\frac{1}{r-q+1}\left(\frac{1}{\prod_{p=q}^r (\lambda-x)^2+p\eta^2}
-\frac{1}{\prod_{p=q+1}^{r+1} (\lambda-x)^2+p\eta^2}\right)
=\frac{\eta^2}{\prod_{p=q}^{r+1} (\lambda-x)^2+p\eta^2},\]
we then obtain
\begin{equation}\label{eq:finitediffbound}
\frac{1}{N}\sum_{j=1}^N \frac{\eta^{2P}}{\prod_{p=1}^P
[(\lambda_j(\widetilde \vK)-x)^2+p\eta^2]}
-\int \frac{\eta^{2P}}{\prod_{p=1}^P[(\lambda-x)^2+p\eta^2]}
d\tilde\mu_N(\lambda) \prec \frac{\eta}{N}.
\end{equation}
Since $\dist{x,\cS_N} \geq \eps$, the second integral term of
\eqref{eq:finitediffbound} is bounded by $C\eta^{2P}$ for a
constant $C:=C(\eps,P)>0$. Thus, we get
\[\frac{1}{N}\sum_{j=1}^N \bI\{\lambda_j(\widetilde \vK) \in (x-\eta,x+\eta)\}
\leq \frac{C}{N}\sum_{j=1}^N \frac{\eta^{2P}}{\prod_{p=1}^P
[(\lambda_j(\widetilde \vK)-x)^2+p\eta^2]} \prec \frac{\eta}{N}+\eta^{2P}\]
where the first inequality holds
for a constant $C:=C(P)>0$. Finally, recalling $\eta=N^{-1/11}$ and taking
any $P \geq 6$, we get $\eta/N+\eta^{2P} \ll 1/N$, hence
\[\bI\big\{\text{there exists an eigenvalue of }\widetilde \vK
\text{ in } (x-\eta,x+\eta)\big\} \prec 0.\]
Recalling Assumption \ref{assump:G}(b) and
taking a union bound over $x$ belonging to a $\eta$-net of
$[-B,B] \setminus (\cS_N+(-\eps,\eps))$ (with cardinality at most $CN^{1/11}$),
we obtain
\[\bI\big\{\text{there exists an eigenvalue of }\widetilde \vK
\text{ in } \cS_N+(-\eps,\eps)\big\} \prec 0.\]
The theorem follows from the observation that $\vK$ has the same non-zero
eigenvalues as $\widetilde \vK$, and all 0 eigenvalues belong by definition to
$\cS_N$.
\end{proof-of-theorem}

\subsection{Deterministic equivalent for the
resolvent}\label{sec:deterministicequiv}

In this section, we prove Theorem \ref{thm:deterministic_equ}. 

\begin{lemma}\label{lemma:supportcomparea}
Suppose Assumption \ref{assump:G} holds. Let
\[\gamma_N^{(S)}=\frac{n}{N-|S|}, \qquad
\mu_N^{(S)}=\rho_{\gamma_N^{(S)}}^\MP \boxtimes \nu_N, \qquad
\tilde \mu_N^{(S)}=\gamma_N^{(S)} \mu_N^{(S)}+(1-\gamma_N^{(S)})\delta_0\]
be the analogues of $\gamma_N,\mu_N,\tilde \mu_N$ defined with the dimension
$N-|S|$ in place of $N$.
Then for any fixed $\eps>0$ and $L \geq 1$, all large $N$, and all
$S \subset [N]$ with $|S| \leq L$,
\[\supp{\tilde \mu_N^{(S)}} \subseteq \supp{\tilde \mu_N}+(-\eps,\eps)\]
\end{lemma}
\begin{proof}
Let $\cT_N$ and $z_N:\C \setminus \cT_N \to \C$ be as defined by
\eqref{eq:zNdomain} and \eqref{eq:z_N_def}. Define similarly
\[z_N^{(S)}(\tilde m)=-\frac{1}{\tilde m}+\gamma_N^{(S)}
\int \frac{\lambda}{1+\lambda \tilde m}d\nu_N(\lambda), \qquad
z_N^{(S)}:\C \setminus \cT_N \to \C.\]
We recall from Proposition \ref{prop:inv_z} that
$x \in \R \setminus \supp{\tilde\mu_N}$ if and only if there exists
$\tilde m \in \R \setminus \cT_N$ where $z_N(\tilde m)=x$ and
$z_N'(\tilde m)>0$; the analogous characterization holds for
$\R \setminus \supp{\tilde \mu_N^{(S)}}$ and $z_N^{(S)}(\tilde m)$.

Now fix any $\eps,L>0$. By Proposition \ref{prop:basic_mn}, there is a constant
$C_0>0$ such that $\supp{\tilde \mu_N^{(S)}} \subseteq [-C_0,C_0]$ for all $|S|
\leq L$ and all large $N$. Consider any $x \in [-C_0,C_0] \setminus
(\supp{\tilde \mu_N}+(-\eps,\eps))$. Then $[x-\eps/2,x+\eps/2] \subset \R
\setminus \supp{\tilde \mu_N}$, so $\tilde m_N$ is well-defined and increasing on
$[x-\eps/2,x+\eps/2]$.
Define $[\tilde m_-,\tilde m_+]=[\tilde m_N(x-\eps/2), \tilde m_N(x+\eps/2)]$.
Then Proposition \ref{prop:inv_z} implies that $z_N$ is increasing on
$[\tilde m_-,\tilde m_+]$, and $z_N([\tilde m_-,\tilde m_+])=[x-\eps/2,x+\eps/2]$.
Again by Proposition \ref{prop:basic_mn}, there is a constant $c>0$ such that,
for any such $x \in [-C_0,C_0] \setminus (\supp{\tilde \mu_N}+(-\eps,\eps))$, we
have
\[\min_{y \in [x-\eps/2,x+\eps/2]} \min_{\lambda \in
\supp{\nu_N}} |1+\lambda \tilde m_N(y)|>c.\]
This then implies that there is a constant $C>0$ for which
\[|z_N^{(S)}(\tilde m)-z_N(\tilde m)|
=|\gamma_N^{(S)}-\gamma_N| \cdot \left|\int \frac{\lambda}{1+\lambda \tilde
m}d\nu_N(\lambda)\right| \leq \frac{C}{N}<\eps/2\]
for all $\tilde m \in [\tilde m_-,\tilde m_+]$,
$|S| \leq L$, and large $N$. Then $z_N^{(S)}(\tilde m_-)<z_N(\tilde m_-)+\eps/2=x$
and $z_N^{(S)}(\tilde m_+)>z_N(\tilde m_+)-\eps/2=x$.
\cite[Theorem 4.3]{silverstein1995analysis} shows that if
$m_1,m_2 \in [\tilde m_-,\tilde m_+]$ satisfy ${z_N^{(S)}}'(m_1) \geq 0$ and
${z_N^{(S)}}'(m_2) \geq 0$, then ${z_N^{(S)}}'(m)>0$ strictly for all $m \in
[m_1,m_2]$. By this and the continuity and differentiability of
$z_N^{(S)}$ on $[\tilde m_-,\tilde m_+]$, there must be a point $\tilde m \in (\tilde
m_-,\tilde m_+)$ where $z_N^{(S)}(\tilde m)=x$ and ${z_N^{(S)}}'(\tilde m)>0$
strictly. Then Proposition \ref{prop:inv_z} implies that
$x \notin \supp{\tilde \mu_N^{(S)}}$. This holds for all
$x \in [-C_0,C_0] \setminus (\supp{\tilde \mu_N}+(-\eps,\eps))$, implying
$\supp{\tilde \mu_N^{(S)}} \subseteq \supp{\tilde \mu_N}+(-\eps,\eps)$ as desired.
\end{proof}

The following now applies Lemma \ref{lemma:supportcomparea} and
Theorem \ref{thm:no_outlier} to extend the
estimates (\ref{eq:checkFAconditions}) previously obtained over
$\{z \in U_N(\eps):\Im z \geq N^{-1/11}\}$ to all of $U_N(\eps)$.
\begin{lemma}\label{lemma:FAconditions}
Fix any $\eps>0$ and $L \geq 1$. Then for some constants $C_0,c_0>0$, uniformly
over $z \in U_N(\eps)$, $S \subset [N]$ with $|S| \leq L$, and $i \in S$,
we have
\[\begin{gathered}
\bI\{|\tilde m_{\vK}^{(S)}(z)|>C_0\} \prec 0,
\quad \bI\{|\tilde m_{\vK}^{(S)}(z)|<c_0\} \prec 0,
\quad \bI\{\|(\vI+\tilde m_{\vK}^{(S)}(z)\vSigma)^{-1}\|>C_0\} \prec 0,\\
\bI\{\|\vR^{(S)}(z)\|>C_0\} \prec 0,
\quad \bI\{|1+N^{-1}\Tr \vSigma\vR^{(S)}(z)|<c_0\} \prec 0,\\
\bI\{|1+N^{-1}\vg_i^\top \vR^{(S)}(z)\vg_i|<c_0\} \prec 0.
\end{gathered}\]
\end{lemma}
\begin{proof}
By conjugation symmetry, it suffices to show the statements for
$z \in U_N(\eps)$ with $\Im z \geq 0$. Denote for simplicity
$\vR^{(S)}=\vR^{(S)}(z)$ and $\tilde m_{\vK}^{(S)}=\tilde m_{\vK}^{(S)}(z)$.
Let $\cS_N^{(S)}=\supp{\mu_N^{(S)}} \cup \{0\}
=\supp{\tilde \mu_N^{(S)}} \cup \{0\}$ where $\mu_N^{(S)},\tilde \mu_N^{(S)}$
are as defined in Lemma \ref{lemma:supportcomparea}.
Then Theorem~\ref{thm:no_outlier} applied to $\vK^{(S)}$ guarantees that
\[\bI\{\vK^{(S)} \text{ has an eigenvalue outside } \cS_N^{(S)}+(-\eps/4,\eps/4)\} \prec 0,\]
uniformly over all $S \subset [N]$ with $|S| \leq L$. Note that
$\cS_N^{(S)}+(-\eps/4,\eps/4) \subseteq \cS_N+(-\eps/2,\eps/2)$ by
Lemma~\ref{lemma:supportcomparea}. Then, applying the bound $\|\vR^{(S)}\|
\leq 1/\dist{z,\cS_N^{(S)}}$ and the condition $z \in U_N(\eps)$, we get
\begin{equation}\label{eq:operatornormbound}
\bI\{\|\vR^{(S)}\|>2/\eps\} \prec 0.
\end{equation}

The remaining statements have already been shown for $z \in U_N(\eps)$ with
$\Im z \geq N^{-1/11}$ in \eqref{eq:checkFAconditions}. 
For $z=x+i\eta$ where $\eta \in [0,N^{-1/11}]$, define $z'=x+i N^{-1/11}$.
On the event that $\vK^{(S)}$ has no eigenvalues
outside $\cS_N+(-\eps/2,\eps/2)$, both $N^{-1}\Tr \vSigma \vR^{(S)}(z)$ and
$\tilde m_{\vK}^{(S)}(z)=N^{-1}\Tr \vR^{(S)}(z)+\gamma_N(-1/z)$ are
$C$-Lipschitz over $z \in U_N(\eps)$ for a constant $C=C(\eps)>0$,
and $N^{-1}\vg_i^\top \vR^{(S)}(z)\vg_i$ is $CN^{-1}\|\vg_i\|^2$-Lipschitz where
$N^{-1}\|\vg_i\|^2 \prec 1$ by Assumption \ref{assump:G}. Then
\[N^{-1}\Tr \vSigma \vR^{(S)}(z)-N^{-1}\Tr \vSigma \vR^{(S)}(z')
\prec N^{-1/11}, \quad
\tilde m_{\vK}^{(S)}(z)-\tilde m_{\vK}^{(S)}(z')
\prec N^{-1/11},\]
\[N^{-1}\vg_i^\top \vR^{(S)}(z)\vg_i-N^{-1}\vg_i^\top \vR^{(S)}(z')\vg_i
\prec N^{-1/11},\]
so the remaining statements of the lemma hold also for $z \in U_N(\eps)$
with $\Im z \in [0,N^{-1/11}]$.
\end{proof}

\begin{proof-of-theorem}[\ref{thm:deterministic_equ}]
Again by conjugation symmetry, it suffices to show the result for
$z \in U_N(\eps)$ with $\Im z \geq 0$. Denote for simplicity
$\vR^{(S)}=\vR^{(S)}(z)$ and $\tilde m_{\vK}^{(S)}=\tilde m_{\vK}^{(S)}(z)$.
The first estimate of Lemma \ref{lemma:FAconditions} implies
\begin{equation}\label{eq:FAconditions1}
\bI\{\|\vR^{(S)}\|_F>C\sqrt{N}\} \prec 0,
\qquad \bI\{\|\vR^{(S)}\vA\|_F>C\|\vA\|_F\} \prec 0
\end{equation}
uniformly over $z \in U_N(\eps)$ and $\vA \in \C^{n \times n}$. Then also, by
Assumption \ref{assump:G}(c) and \eqref{eq:mdiffestimate} applied with
$\vB=\vSigma$,
\begin{align}
1+N^{-1}\vg_i^\top \vR^{(i)}\vg_i
&=1+N^{-1}\Tr \vSigma \vR+\SD{N^{-2}|1+N^{-1}\vg_i^\top
\vR^{(i)}\vg_i|^{-1}\|\vR^{(i)}\|_F^2+\|\vR^{(i)}\|_F}\nonumber\\
&=1+N^{-1}\Tr \vSigma \vR+\SD{N^{-1/2}}.\label{eq:1plusgRgfull}
\end{align}

Let $d_i=d_{i,1}+d_{i,2}+d_{i,3}+d_{i,4}$ be as defined in \eqref{eq:di14} with
$\vA=\vI$. Then, applying \eqref{eq:FAconditions1}, \eqref{eq:1plusgRgfull},
and the bounds of Lemma \ref{lemma:FAconditions}, we obtain exactly as in the
proof of Lemma \ref{lemma:improvedestimate} (using again the fluctuation
averaging result of Lemma \ref{lemma:fluctuationavg}) that,
uniformly over $z \in
U_N(\eps)$, we have $|d_{i,1}|,|d_{i,3}|,|d_{i,4}| \prec N^{-1}$,
$|d_{i,2}| \prec N^{-1/2}$, and
\[|z_N(m_{\tilde \vK})-z| \prec
\frac{1}{N} \cdot \frac{1}{1+N^{-1}\Tr \vSigma\vR}
\cdot \left|\sum_{i=1}^N d_{i,2}\right|+\SD{N^{-1}}=\SD{N^{-1}}.\]
Fix any $\iota>0$. If $\Im z \geq N^{-1+\iota}$, then this
implies $\bI\{z_N(m_{\tilde \vK}) \notin \C^+\} \prec 0$. By the same
stability argument as in Lemma \ref{lemma:prelimestimate}, we get
$m_{\tilde \vK}(z)-\tilde m_N(z) \prec N^{-1}$
uniformly over $z \in U_N(\eps)$ with $\Im z \geq N^{-1+\iota}$. For
$\Im z \in [0,N^{-1+\iota}]$, on the event that all eigenvalues of $\vK$ belong
to $\cS_N+(-\eps/2,\eps/2)$, we may apply that both $m_{\tilde \vK}(z)$
and $\tilde m_N(z)$ are $C(\eps)$-Lipschitz over $z \in U_N(\eps)$ to
compare values at $z=x+i\eta$ and $z'=x+iN^{-1+\iota}$.
Applying $m_{\tilde \vK}(z')-\tilde m_N(z') \prec N^{-1}$, we then get for any
$D>0$, all $z \in U_N(\eps)$, some constant $C>0$, and all large $N$,
\[\P[|m_{\tilde \vK}(z)-\tilde m_N(z)|>CN^{-1+\iota}] \leq N^{-D}.\]
Since $\iota>0$ is arbitrary, this shows
$m_{\tilde \vK}(z)-\tilde m_N(z) \prec N^{-1}$ uniformly over $z \in U_N(\eps)$.
The bound $m_{\vK}(z)-m_N(z) \prec N^{-1}$ then follows from
$m_{\tilde \vK}(z)=\gamma_N m_{\vK}(z)+(1-\gamma_N)(-1/z)$ and
$\tilde m_N(z)=\gamma_N m_N(z)+(1-\gamma_N)(-1/z)$.

For the estimate of $\Tr \vR\vA$, we apply the definition of
$d_i=d_{i,1}+d_{i,2}+d_{i,3}+d_{i,4}$ from \eqref{eq:di14}
and the identity \eqref{eq:detequividentity} now with this matrix $\vA$. Then
\eqref{eq:detequividentity} gives
\[\Tr \Big[\vR\vA-({-}z\vI-zm_{\tilde \vK}\vSigma)^{-1}\vA\Big]
=\frac{1}{z}\sum_{i=1}^N \frac{d_i}{1+N^{-1}\vg_i^\top \vR^{(i)}\vg_i}.\]
Applying \eqref{eq:FAconditions1}, \eqref{eq:1plusgRgfull},
and the bounds of Lemma \ref{lemma:FAconditions} to
\eqref{eq:di1}--\eqref{eq:di4}, uniformly over $z \in U_N(\eps)$ and $\vA \in
\C^{n \times n}$, we have $|d_{i,1}|,|d_{i,3}|,|d_{i,4}|
\prec N^{-3/2}\|\vA\|_F$, $|d_{i,2}| \prec N^{-1}\|\vA\|_F$, and hence
\[\left|\sum_{i=1}^N \frac{d_i}{1+N^{-1}\vg_i^\top \vR^{(i)}\vg_i}\right|
\prec \frac{1}{1+N^{-1}\Tr \vSigma\vR}\left|\sum_{i=1}^N d_{i,2}\right|
+\SD{N^{-1/2}\|\vA\|_F}.\]
Finally, applying Lemma \ref{lemma:fluctuationavg} with $\vGamma=z\vI$,
$\Psi_N(\vGamma)=C\sqrt{N}$, and $\Phi_N(\vGamma,\vA)=C\|\vA\|_F$ (where we may
assume without loss of generality $\|\vA\|_F \in (N^{-\upsilon},N^\upsilon)$
by scale invariance of the desired estimate with respect to $\vA$),
we get $|\sum_i d_{i,2}| \prec N^{-1/2}\|\vA\|_F$. Thus,
\[\Tr \Big[\vR\vA-({-}z\vI-zm_{\tilde \vK}\vSigma)^{-1}\vA\Big]
\prec \frac{1}{\sqrt{N}}\,\|\vA\|_F.\]
\end{proof-of-theorem}

%%%%%%%%%%%%%%%%%%%%%%%%%%%%%%%%%%%%%%%%%%%%%%%%%%%%%%%%%%%%%%%%%%%%%%%%%%%%%%%%%%
%%%%%%%%%%%%%%%%%%%%%%%%%%%%%%%%%%%%%%%%%%%%%%%%%%%%%%%%%%%%%%%%%%%%%%%%%%%%%%%%%%

\section{Analysis of spiked eigenstructure}\label{appendix:spike}

We now consider the asymptotic setup of Appendix \ref{subsec:spikes} and
prove Corollary \ref{cor:no_outlier} and Theorem \ref{thm:spiked}.
As all the desired
statements are invariant under conjugation of $\vSigma$ by an orthogonal matrix,
we may assume without loss of generality that $\vSigma$ is diagonal and of the
form
\[\vSigma=\begin{pmatrix}
	\vSigma_r & \mathbf{0}\\
	\mathbf{0} & \vSigma_0
\end{pmatrix}, \qquad
\vSigma_r=\diag(\lambda_1(\vSigma),\ldots,\lambda_r(\vSigma)),
\qquad \vSigma_0=\diag(\lambda_{r+1}(\vSigma),\ldots,\lambda_n(\vSigma)).\]
Denote the block decomposition of $\vG$ corresponding to
$\vSigma_r,\vSigma_0$ as
\[\vG=[\vG_r,\vG_0], \qquad \vG_r \in \R^{N \times r},\qquad
\vG_0 \in \R^{N \times (n-r)}.\]
We remind the reader that $\vG_r$ and $\vG_0$ need not be independent.

\subsection{No outliers outside the limit support}

We consider first the setting of $r=0$, and prove Corollary~\ref{cor:no_outlier}
together with some uniform convergence properties of
$\tilde m_N$ and $z_N$ that will be used in the later analysis.

Recall the domain $\cT_N$ and function $z_N:\C \setminus \cT_N \to \C$ from
\eqref{eq:zNdomain} and \eqref{eq:z_N_def}, and their asymptotic
analogues $\cT$ and $z:\C
\setminus \cT \to \C$ from \eqref{eq:zdomain} and \eqref{eq:inv_z}.

\begin{lemma}\label{lemma:supportcompareb}
Suppose Assumption \ref{assump:G} holds, and Assumption \ref{assum:spike} holds
with $r=0$. Then, as $N \to \infty$,
\begin{enumerate}[label=(\alph*)]
\item $z_N(\tilde m)$ and its derivative $z_N'(\tilde m)$ converge uniformly
over compact subsets of $\C \setminus \cT$ to $z(\tilde m)$ and $z'(\tilde m)$.
\item For any $\eps>0$ and all large $N$,
\[\supp{\tilde \mu_N} \subseteq \supp{\tilde \mu}+(-\eps,\eps).\]
\item $\tilde m_N(z)$ and its derivative $\tilde m_N'(z)$ converge uniformly
over compact subsets of $\C \setminus \supp{\tilde \mu}$ to $\tilde m(z)$ and
$\tilde m'(z)$.
\end{enumerate}
\end{lemma}
\begin{proof}
For part (a), let $K \subset \C \setminus \cT$ be any fixed compact set.
Then $K$ does not intersect some sufficiently small open neighborhood of the compact
domain $\cT$. If Assumption \ref{assum:spike} holds with $r=0$, then $\cT_N$ is
contained in this open neighborhood of $\cT$ for all large $N$, so $K
\subset \C \setminus \cT_N$, and both $z_N$ and
$z$ are well-defined on $K$. The pointwise convergences $z_N(\tilde m) \to z(\tilde
m)$ and $z_N'(\tilde m) \to z'(\tilde m)$ on $K$ then follow from $\gamma_N \to
\gamma$, the weak convergence $\nu_N \to \nu$, and the uniform boundedness of the
functions $\lambda \mapsto \lambda/(1+\lambda \tilde m)$
and $\lambda \mapsto \lambda^2/(1+\lambda \tilde m)^2$ on an open
neighborhood of $\supp{\nu}$, for $\tilde m \in K$. This convergence is
furthermore uniform because $\{z_N\}$ and $\{z_N'\}$ are both equicontinuous
over $K$.

For part (b), consider any $x \notin \supp{\tilde \mu}+(-\eps,\eps)$.
Then $[x-\eps/2,x+\eps/2] \subset \R \setminus \supp{\tilde \mu}$, so $\tilde m$ is
well-defined and increasing on $[x-\eps/2,x+\eps/2]$. Let $[\tilde m_-,\tilde m_+]
=[\tilde m(x-\eps/2),\tilde m(x+\eps/2)]$. Then by Proposition~\ref{prop:inv_z},
$z'(\tilde m)>0$ for all $\tilde m \in [\tilde m_-,\tilde m_+]$, and
$z([\tilde m_-,\tilde m_+])=[x-\eps/2,x+\eps/2]$. The uniform convergence in
part (a) implies for all large $N$ that $z_N(\tilde m_-)<x$,
$z_N(\tilde m_+)>x$, and $z_N'(\tilde m)>0$ for all $\tilde m \in [\tilde
m_-,\tilde m_+]$. Then there exists $\tilde m \in [\tilde m_-,\tilde m_+]$
where $z_N(\tilde m)=x$ and $z_N'(\tilde m)>0$, implying by
Proposition \ref{prop:inv_z} that $x \notin \supp{\tilde \mu_N}$. So
$\supp{\tilde \mu_N} \subseteq \supp{\tilde \mu}+(-\eps,\eps)$ as desired.

For part (c), let $K \subset \C \setminus \supp{\tilde \mu}$ be any
fixed compact set. Then $K$ does not intersect some sufficiently small open
neighborhood of the compact set $\supp{\tilde \mu}$, so the inclusion of part
(b) implies $K \subset \C \setminus \supp{\tilde \mu_N}$ for all large $N$, and
both $\tilde m_N$ and $\tilde m$ are well-defined on $K$. The
uniform convergence $\tilde m_N(z) \to \tilde m(z)$ and $\tilde m_N'(z) \to
\tilde m'(z)$ on $K$ then follow from the weak convergence $\tilde \mu_N \to
\tilde \mu$, the uniform boundedness of the functions $\lambda \mapsto
1/(\lambda-z)$ and $\lambda \mapsto 1/(\lambda-z)^2$ on an open neighborhood of
$\supp{\tilde \mu}$ for $z \in K$,
and the equicontinuity of $\{\tilde m_N\}$ and $\{\tilde m_N'\}$ on $K$.
\end{proof}

\begin{proof-of-coro}[\ref{cor:no_outlier}]
By Lemma \ref{lemma:supportcompareb}(b), for any fixed $\eps>0$, we have
$\cS_N+(-\eps/2,\eps/2) \subseteq \cS+(-\eps,\eps)$
for all large $N$. Then by Theorem \ref{thm:no_outlier},
\begin{align*}
&\bI\{\vK \text{ has an eigenvalue in } \R \setminus (\cS+(-\eps,\eps))\\
&\leq \bI\{\vK \text{ has an eigenvalue in } \R \setminus
(\cS_N+(-\eps/2,\eps/2))\} \prec 0.
\end{align*}
\end{proof-of-coro}

\subsection{Deterministic equivalents for generalized resolvents}

We next introduce two generalized resolvents for the matrix $\vK$, and extend
Theorem \ref{thm:deterministic_equ} to establish
deterministic equivalents for these generalized resolvents.

Define the spectral domain
\[U(\eps)=\Big\{z \in \C:|z| \leq \eps^{-1},\;\dist{z,\cS} \geq \eps\Big\}\]
where $\cS$ is the limit support set defined in (\ref{eq:limitsupport}).
Given $z \in U(\eps)$ and $\alpha \in \C$, define a diagonal matrix
\begin{equation}\label{eq:Gamma}
\vGamma:=\vGamma(z,\alpha)=z\vI_n+\alpha\vV_r\vV_r^\top
=\begin{pmatrix} (z+\alpha)\vI_r & \mathbf{0} \\
\mathbf{0} & z \vI_{n-r} \end{pmatrix} \in \C^{n \times n},
\quad \vV_r=\begin{pmatrix} \vI_r \\ \mathbf{0} \end{pmatrix}
\in \R^{n \times r}.
\end{equation}
Define the first generalized resolvent
\begin{equation}\label{eq:matrix_block}
\vcR(z,\alpha)=
\begin{pmatrix}
	-\vGamma & \vG^\top\\ \vG & -\vI_N
\end{pmatrix}^{-1} \in \C^{(n+N)\times (n+N)}.
\end{equation}
This matrix inverse exists if and only if
the Schur complement $\vG^\top \vG-\vGamma=\vK-\vGamma$
for its lower right block is invertible, in which case the upper-left block of
$\vcR(z,\alpha)$ is $\vR(\vGamma)=(\vK-\vGamma)^{-1}$. The following
provides a deterministic equivalent for this block of $\vcR(z,\alpha)$.

\begin{lemma}\label{lemm:apply_fluctuation}
Under the assumptions of Theorem \ref{thm:spiked}, for any fixed $\eps>0$,
there exist $C_0,\alpha_0>0$ (depending on $\eps$) such that fixing any $\alpha
\in \C$ with $|\alpha|>\alpha_0$, the following hold:
\begin{enumerate}[label=(\alph*)]
\item The event
\[\cE=\Big\{\vcR(z,\alpha) \text{ exists and } \|\vcR(z,\alpha)\| \leq C_0
\text{ for all } z \in U(\eps)\Big\}\]
satisfies $\bI\{\cE^c\} \prec 0$.
\item Uniformly over $z \in U(\eps)$
and deterministic unit vectors $\vv_1,\vv_2 \in \R^n$,
\begin{align}
\norm{   \begin{pmatrix}
    \vv_1^\top&\mathbf{0}
    \end{pmatrix}\vcR(z,\alpha)\begin{pmatrix}
    \vv_2\\
        \mathbf{0}
    \end{pmatrix}+  \vv_1^\top \left(\vGamma+z\cdot\tilde{m}_{N,0}(z)\vSigma
\right)^{-1}\vv_2 } \prec \frac{1}{\sqrt{N}}.
\end{align} 
\end{enumerate}
\end{lemma}

In the setting of Theorem \ref{thm:spiked}(c), let
$\vu=\frac{1}{\sqrt{N}}(u_1,\ldots,u_N) \in \R^N$ be the
additional given vector for which $\{(u_j,\vg_j^\top)\}_{j=1}^N$ are
independent vectors in $\R^{n+1}$. For $z \in U(\eps)$ and $\alpha \in \C$, define
\begin{align}
\widetilde{\vSigma}&=\begin{pmatrix}
		\E[u^2] & \E[u\vg]^\top\\
		\E[u\vg] & \vSigma
	\end{pmatrix}	\in\R^{(n+1)\times (n+1)}\label{eq:tilde_Sigma},\\
\widetilde \vGamma&=\widetilde \vGamma(z,\alpha)
=\begin{pmatrix} z+\alpha & \mathbf{0} \\
\mathbf{0} & \vGamma \end{pmatrix} \in \C^{(n+1) \times (n+1)}
\nonumber
\end{align}
where $\E[u^2]$ and $\E[u\vg]$ denote the common values of $\E[u_j^2]$ and
$\E[u_j\vg_j]$ for $j=1,\ldots,N$. Define the second generalized resolvent
\begin{equation}\label{eq:tildeR}
\widetilde \vcR(z,\alpha)=
\begin{pmatrix} {-}\widetilde \vGamma & [\vu,\vG]^\top \\
[\vu,\vG] & -\vI \end{pmatrix}^{-1}
=\begin{pmatrix} {-}(z+\alpha) & \mathbf{0} & \vu^\top \\
\mathbf{0} & -\vGamma & \vG^\top \\
\vu & \vG & -\vI_N \end{pmatrix}^{-1}
\in \C^{(n+1+N)\times (n+1+N)}.
\end{equation}
We have the following
deterministic equivalent for the upper-left block of $\widetilde \vcR(z,\alpha)$,
which is analogous to Lemma \ref{lemm:apply_fluctuation}.

\begin{lemma}\label{lemm:apply_fluctuation2}
Under the assumptions of Theorem \ref{thm:spiked}(c), for any fixed $\eps>0$, there
exist $C_0,\alpha_0>0$ (depending on $\eps$) such that fixing any $\alpha \in
\C$ with $|\alpha|>\alpha_0$, the following hold:
\begin{enumerate}[label=(\alph*)]
\item The event
\[\widetilde \cE=\Big\{\widetilde \vcR(z,\alpha) \text{ exists and } \|\widetilde \vcR(z,\alpha)\| \leq C_0
\text{ for all } z \in U(\eps)\Big\}\]
satisfies $\bI\{\widetilde \cE^c\} \prec 0$.
\item Uniformly over $z \in U(\eps)$ and
deterministic unit vectors $\vv_1,\vv_2\in\R^{n+1}$,
\begin{align}
		\norm{   \begin{pmatrix}
			\vv_1^\top&\mathbf{0}
		\end{pmatrix}\widetilde \vcR(z,\alpha)\begin{pmatrix}
			\vv_2\\
			\mathbf{0}
		\end{pmatrix}+
\vv_1^\top \left(\widetilde\vGamma+z\cdot\tilde{m}_{N,0}(z)\widetilde\vSigma
\right)^{-1}\vv_2 } \prec \frac{1}{\sqrt{N}}.
\end{align}
\end{enumerate}
\end{lemma}

In the remainder of this section, we prove Lemmas \ref{lemm:apply_fluctuation}
and \ref{lemm:apply_fluctuation2}. Recall 
\[\mu_{N,0}=\rho_{\gamma_{N,0}}^{\MP} \boxtimes\nu_{N,0}, \qquad
\tilde \mu_{N,0}=\gamma_{N,0}\mu_{N,0}+(1-\gamma_{N,0})\delta_0.\]
Define the bulk components of the sample covariance and Gram matrices
\begin{equation}\label{eq:K0def}
\vK_0=\vG_0^\top \vG_0 \in \R^{(n-r) \times (n-r)},
\qquad \widetilde \vK_0=\vG_0\vG_0^\top \in \R^{N \times N}.
\end{equation}
Define also the $N$-dependent bulk spectral support and spectral domain
\begin{align}
\cS_{N,0}&=\supp{\mu_{N,0}} \cup \{0\}
=\supp{\tilde\mu_{N,0}} \cup \{0\},\nonumber\\
U_{N,0}(\eps)&=\{z \in \C:|z| \leq \eps^{-1},\dist{z,\cS_{N,0}} \geq \eps\}.\label{eq:UN0}
\end{align}
Lemma \ref{lemma:supportcompareb}(b) shows
$\cS_{N,0} \subseteq \cS+(-\eps/2,\eps/2)$ for any fixed $\eps>0$ and all large $N$,
so also $U(\eps) \subseteq U_{N,0}(\eps/2)$ for all large $N$. Thus, the results of
Appendix \ref{appendix:resolvent} applied to $\vK_0$, which hold uniformly over
$z \in U_{N,0}(\eps/2)$ for any fixed $\eps>0$,
also hold uniformly over $z \in U(\eps)$. In particular,
the following is an immediate consequence of Corollary \ref{cor:no_outlier} and
Theorem \ref{thm:deterministic_equ},
which we record here for future reference.

\begin{lemma}\label{lemm:support_G0}
Suppose Assumptions~\ref{assump:G} and \ref{assum:spike} hold.
Then for any fixed $\eps>0$,
\begin{align*}
\bI\{\vK_0 \text{ has an eigenvalue outside }
\cS+(-\eps,\eps)\} \prec 0.
\end{align*}
Furthermore, uniformly over $z \in U(\eps)$,
\[m_{\vK_0}-m_{N,0}(z)\prec 1/N, \quad  m_{\tilde\vK_0}-\tilde m_{N,0}(z)
\prec 1/N.\]
\end{lemma}

We now check that for sufficiently large $|\alpha|$, the generalized resolvent
$\vcR(z,\alpha)$ exists and has bounded operator norm with high probability.\\

\begin{proof-of-lemma}[\ref{lemm:apply_fluctuation}(a)]
Let
\[\cE'=\left\{\text{ all eigenvalues of } \vK_0 \text{ belong to }
\cS+(-\eps/2,\eps/2),\text{ and }\norm{\vG}<\sqrt{B}\right\}.\]
By Assumption \ref{assump:G}(b) and Lemma \ref{lemm:support_G0},
$\bI\{{\cE'}^c\} \prec 0$, so it suffices to show $\cE' \subseteq \cE$. 
On this event $\cE'$, for any $z \in U(\eps)$, we have that each
eigenvalue of $\vK_0$ is separated by at least $\eps/2$ from $z$.
Then,
	\begin{equation}\label{eq:block_G0}
		\vcR_0(z):=\begin{pmatrix}
		-z\vI_{n-r} & \vG_0^{\top}\\
			\vG_0 & -\vI_N
		\end{pmatrix}^{-1} \in \C^{(n-r+N) \times (n-r+N)}
	\end{equation}
exists for all $z \in U(\eps)$
because the Schur complement $\vK_0-z\vI_{n-r}$ of its lower-right block
is invertible.
Furthermore, denoting $\vR_0=(\vK_0-z\vI_{n-r})^{-1}$, we have
$\|\vR_0\| \leq 2/\eps$ and $\|\vG_0\| \leq \|\vG\|<\sqrt{B}$, so
	\begin{align}
\|\vcR_0(z)\|=\left\|\begin{pmatrix} \vR_0 & \vR_0\vG_0^\top
\\ \vG_0\vR_0 &
\vG_0\vR_0\vG_0^\top-\vI_N\end{pmatrix}\right\|
	\leq C_1\label{eq:R(z)_bound}
	\end{align}
for some constant $C_1$ depending only on $\eps,B$.

Now write $\vcR(z,\alpha)$ as defined in \eqref{eq:matrix_block} in its block
decomposition with blocks of sizes $r$ and $n-r+N$. Then the Schur complement
of the upper left block of size $r \times r$ is given by
	\begin{equation}\label{eq:schur_com}
	\vS=-\begin{pmatrix}
		\mathbf{0}&
		\vG_r^\top
	\end{pmatrix}\vcR_0(z)
	\begin{pmatrix}
		\mathbf{0}\\
		\vG_r
	\end{pmatrix}	-(\alpha+z)\vI_{r}.
	\end{equation}
	Notice that \begin{align}
	\vS\vS^*=~& |\alpha+z|^2\vI_{r}+\begin{pmatrix}
		\mathbf{0}&
		\vG_r^\top
	\end{pmatrix}\vcR_0(z)
	\begin{pmatrix}
		\mathbf{0}\\
		\vG_r
	\end{pmatrix}\begin{pmatrix}
		\mathbf{0}&
	\vG_r^\top
	\end{pmatrix}\overline{\vcR_0(z)}
	\begin{pmatrix}
		\mathbf{0}\\
		\vG_r
	\end{pmatrix}\\
	~&+(\bar \alpha+\bar z)\begin{pmatrix}
		\mathbf{0}&
		\vG_r^\top
	\end{pmatrix}\vcR_0(z)
	\begin{pmatrix}
		\mathbf{0}\\
		\vG_r
	\end{pmatrix}+ (\alpha+z)\begin{pmatrix}
		\mathbf{0}&
		\vG_r^\top
	\end{pmatrix}\overline{\vcR_0(z)}
	\begin{pmatrix}
		\mathbf{0}\\
		\vG_r
	\end{pmatrix}
	\end{align}
where the first two terms are positive semi-definite.
	Therefore, applying \eqref{eq:R(z)_bound} and $\|\vG_r\| \leq
\|\vG\|<\sqrt{B}$ on the event $\cE'$,
there exist $\alpha_0,c_0>0$ depending only on $\eps,B$, such that 
	\begin{align}
	\lambda_{\min}(\vS\vS^*) \ge
|\alpha+z|^2-2(|\alpha|+|z|)\norm{\vG_r}^2\norm{\vcR_0( z)}>c_0
	\end{align}
for any $z\in U(\eps)$ and $|\alpha|>\alpha_0$. Consequently, under the event
$\cE'$, the Schur complement $\vS$ in \eqref{eq:schur_com} is
invertible with $\|\vS^{-1}\|<c_0^{-1/2}$.
Then $\vcR(z,\alpha)$ exists, and
	\begin{align}
	\|\vcR(z,\alpha)\|
	&=\left\|\begin{pmatrix}
		\vS^{-1}& -\vS^{-1}\begin{pmatrix}
			\mathbf{0}&
			\vG_r^\top
		\end{pmatrix} \vcR_0(z)\\
		-\vcR_0(z)\begin{pmatrix}
			\mathbf{0}\\
			\vG_r
		\end{pmatrix} \vS^{-1} & \vcR_0(z)+\vcR_0(z)\begin{pmatrix}
			\mathbf{0}\\
			\vG_r
		\end{pmatrix}\vS^{-1}\begin{pmatrix}
			\mathbf{0}&
			\vG_r^\top
		\end{pmatrix}\vcR_0(z)
	\end{pmatrix}\right\| \leq C_0
	\end{align}
for a constant $C_0>0$ depending only on $\eps,B$.
This shows $\cE' \subseteq \cE$ as desired.
        \end{proof-of-lemma}

For the matrix $\vGamma=\vGamma(z,\alpha)$ in \eqref{eq:Gamma}, recall the definitions of
$\vR^{(S)}(\vGamma)$ and $\tilde m_{\vK}^{(S)}(\vGamma)$ from
\eqref{eq:leaveoneout}. The following provides an analogue of
Lemma \ref{lemma:FAconditions} for these quantities.

\begin{lemma}\label{lemm:bounds_R(z)}
Fix any $\eps>0$ and $L \geq 1$. Then there exist $C_0,c_0,\alpha_0>0$ such that
for any fixed $\alpha \in \C$ with $|\alpha|>\alpha_0$,
uniformly over $S \subset [N]$ with $|S| \leq L$, over $j \in S$, and over
$z \in U(\eps)$,
\[\begin{gathered}
\bI\{|\tilde m_{\vK}^{(S)}(\vGamma)|>C_0\} \prec 0,
\quad \bI\{|\tilde m_{\vK}^{(S)}(\vGamma)|<c_0\} \prec 0,
\quad \bI\{\|(z^{-1}\vGamma+\tilde m_{\vK}^{(S)}(\vGamma)\vSigma)^{-1}\|>C_0\} \prec 0,\\
\bI\{\|\vR^{(S)}(\vGamma)\|>C_0\} \prec 0,
\quad \bI\{|1+N^{-1}\Tr \vSigma\vR^{(S)} (\vGamma)|<c_0\} \prec 0,\\
\bI\{|1+N^{-1}\vg_j^\top \vR^{(S)}(\vGamma) \vg_j|<c_0\} \prec 0.
\end{gathered}\]
\end{lemma}
\begin{proof}
Suppose $|\alpha|$ is large enough so that Lemma \ref{lemm:apply_fluctuation}(a)
holds. Since $\vR(\vGamma)$ is the upper-left block of $\vcR(z,\alpha)$,
Lemma~\ref{lemm:apply_fluctuation}(a) applied with $\vG^{(S)}$ in place of $\vG$
shows that $\bI\{\|\vR^{(S)}(\vGamma)\|>C_0\} \prec 0$ for a constant $C_0>0$,
uniformly over $S \subset [N]$ with $|S| \leq L$ and over $z \in U(\eps)$.
For the remaining statements,
let $\vG_0^{(S)} \in \R^{(N-|S|) \times (n-r)}$ be the submatrix of $\vG_0$ with
the rows of $S$ removed, and define
\[\vK_0^{(S)}={\vG_0^{(S)}}^\top\vG_0^{(S)},
\quad \widetilde \vK_0^{(S)}={\vG_0^{(S)}}{\vG_0^{(S)}}^\top,\]
\[\vR_0^{(S)}=(\vK_0^{(S)}-z\vI_{n-r})^{-1},
\quad m_{\vK_0}^{(S)}=\frac{1}{n-r}\Tr \vR_0^{(S)},
\quad \tilde
m_{\vK_0}^{(S)}=\gamma_{N,0}m_{\vK_0}^{(S)}+(1-\gamma_{N,0})\Big({-}\frac{1}{z}\Big).\]
Then by Lemma \ref{lemm:support_G0} applied to $\vK_0^{(S)}$, also
$\bI\{\|\vR_0^{(S)}\|>C_0\} \prec 0$ for a constant $C_0>0$.

Using these bounds, we first show the comparisons
\begin{equation}\label{eq:diff_m}
	|\tilde m_{\vK}^{(S)}(\vGamma)-\tilde m_{\vK_0}^{(S)}| \prec
1/N, \quad
	\left| N^{-1}\Tr \vSigma\vR^{(S)}(\vGamma)-N^{-1}
\Tr \vSigma_0\vR_0^{(S)}\right| \prec 1/N.
\end{equation}
	For the first comparison, notice that in the decompositions into blocks
of sizes $r$ and $n-r$,
\[\frac{n-r}{n}\,m_{ \vK_0}^{(S)}= \frac{1}{n}\Tr \begin{pmatrix}
		\mathbf{0} & \mathbf{0}\\
		\mathbf{0} & \vR_0^{(S)}
	\end{pmatrix}\]
	and 
	\begin{align*}
		m_{\vK}^{(S)}(\vGamma)=~&\frac{1}{n}\Tr \vR^{(S)}(\vGamma)\\
=~&\frac{1}{n}\Tr \begin{pmatrix}
		\vG_r^{(S)\top}\vG_r^{(S)}-(\alpha+z)\vI_r & \vG_r^{(S)\top}\vG_0^{(S)}\\
		\vG_0^{(S)\top}\vG_r^{(S)} & \vK_0^{(S)}-z\vI_{N-|S|}
	\end{pmatrix}^{-1}\\=
	~&\frac{1}{n}\Tr \begin{pmatrix}
		(\vS_r^{(S)})^{-1} & -(\vS_r^{(S)})^{-1}\vG_r^{(S)\top}\vG_0^{(S)}\vR_0^{(S)}\\
		-\vR_0^{(S)}\vG_0^{(S)\top}\vG_r^{(S)}(\vS_r^{(S)})^{-1} & \vR_0^{(S)}+\vR_0^{(S)}\vG_0^{(S)\top}\vG_r^{(S)}(\vS_r^{(S)})^{-1}\vG_r^{(S)\top}\vG_0^{(S)}\vR_0^{(S)}
	\end{pmatrix}, 
	\end{align*}
	where 
	\begin{equation}\label{eq:Schur_r}
		\vS_r^{(S)}:= \vG_r^{(S)\top}\vG_r^{(S)}-(\alpha+z)\vI_r-  \vG_r^{(S)\top}\vG_0^{(S)}\vR_0^{(S)}\vG_0^{(S)\top}\vG_r^{(S)}
	\end{equation}
	is the Schur complement of the lower-right block. We have
$\snorm{(\vS_r^{(S)})^{-1}}\le \snorm{\vR^{(S)}(\vGamma)}\prec 1$,
$\snorm{\vR_0^{(S)}}\prec 1$,
and by Assumption~\ref{assump:G}, $\snorm{\vG_0^{(S)}}\prec 1$ and
$\snorm{\vG_r^{(S)}}\prec 1$. Combining these bounds and using
$|\Tr \vA| \leq r\|\vA\|$ when $\vA$ has rank $r$ (as follows from the
von Neumann trace inequality),
	\begin{align*}
		&\left|m_{ \vK}^{(S)}(\vGamma)-\frac{n-r}{n}\,m_{
\vK_0}^{(S)}\right|\\
&=\left|\frac{1}{n}\Tr\begin{pmatrix}
			(\vS_r^{(S)})^{-1} & -(\vS_r^{(S)})^{-1}\vG_r^{(S)\top}\vG_0^{(S)}\vR_0^{(S)}\\
			-\vR_0^{(S)}\vG_0^{(S)\top}\vG_r^{(S)}(\vS_r^{(S)})^{-1} & \vR_0^{(S)}\vG_0^{(S)\top}\vG_r^{(S)}(\vS_r^{(S)})^{-1}\vG_r^{(S)\top}\vG_0^{(S)}\vR_0^{(S)}
		\end{pmatrix}\right|\\
		&\le \frac{1}{n}|\Tr (\vS_r^{(S)})^{-1}|+\frac{1}{n}|\Tr \vR_0^{(S)}\vG_0^{(S)\top}\vG_r^{(S)}(\vS_r^{(S)})^{-1}\vG_r^{(S)\top}\vG_0^{(S)}\vR_0^{(S)}|\\
		&\le \frac{r}{n}\snorm{(\vS_r^{(S)})^{-1}} +
\frac{r}{n}\snorm{\vG_r^{(S)}}^2
\snorm{\vG_0^{(S)}}^2\snorm{\vR_0^{(S)}}^2\snorm{(\vS_r^{(S)})^{-1}}
\prec 1/N.
	\end{align*}
Then also $|m_{\vK}^{(S)}(\vGamma)-m_{\vK_0}^{(S)}| \prec 1/N$ since
$|m_{\vK_0}^{(S)}|\leq \|\vR_0^{(S)}\| \prec 1$ and $(n-r)/n=1+\SD{1/N}$.
Hence $|\tilde m_{\vK}^{(S)}(\vGamma)-\tilde m_{\vK_0}^{(S)}| \prec 1/N$
from the definitions $\tilde m_{\vK}^{(S)}(\vGamma)=\gamma_N
m_{\vK}^{(S)}(\vGamma)+(1-\gamma_N)(-1/z)$ and
$\tilde m_{\vK_0}^{(S)}=\gamma_{N,0}m_{\vK_0}^{(S)}+(1-\gamma_{N,0})(-1/z)$,
as $|1/z| \leq \eps$ for $z \in U(\eps)$ and $\gamma_{N,0}=\gamma_N+\SD{1/N}$.
The proof of the second comparison of \eqref{eq:diff_m} is analogous,
considering in addition
	\begin{equation}\label{eq:SigmaRcompare}
\frac{1}{n}\Tr \left(\vSigma-\begin{pmatrix}
		\mathbf{0} & \mathbf{0}\\
		\mathbf{0} & \vSigma_0
	\end{pmatrix}\right)\vR^{(S)}(\vGamma)
=\frac{1}{n}\Tr \begin{pmatrix}
		\vSigma_r & \mathbf{0}\\
		\mathbf{0} & \mathbf{0}
	\end{pmatrix}\vR^{(S)}(\vGamma)
\leq \frac{r}{n}\|\vSigma_r\|
\|\vR^{(S)}(\vGamma)\| \prec \frac{1}{N}.
\end{equation}

Now, Lemma~\ref{lemma:FAconditions} applied with $\vK_0$ shows, uniformly over $S
\subset [N]$ with $|S| \leq L$ and over $z \in U(\eps)$,
\[\bI\{|\tilde m_{\vK_0}^{(S)}|>C_0\} \prec 0, \quad
	\bI\{|\tilde m_{\vK_0}^{(S)}|<c_0\} \prec 0, \quad
\bI\{|1+N^{-1}\Tr \vSigma_0\vR_0^{(S)}|<c_0\} \prec 0,\]
which together with \eqref{eq:diff_m} implies
	\[\bI\{|\tilde m_{\vK}^{(S)}(\vGamma)|>C_0\} \prec 0, \quad
	\bI\{|\tilde m_{\vK}^{(S)}(\vGamma)|<c_0\} \prec 0, \quad
\bI\{|1+N^{-1}\Tr \vSigma\vR^{(S)}(\vGamma)|<c_0\} \prec 0\]
for adjusted constants $C_0,c_0>0$.
Also by Assumption~\ref{assump:G}, uniformly over $j \in S$,
	\begin{equation}\label{eq:concnetr_g_i}
	N^{-1}\vg_j^\top \vR^{(S)}(\vGamma)\vg_j-N^{-1}\Tr
\vSigma\vR^{(S)}(\vGamma)\prec N^{-1}\|\vR^{(S)}(\vGamma)\|_F \leq
N^{-1/2}\|\vR^{(S)}(\vGamma)\| \prec N^{-1/2},
	\end{equation}
	so $\bI\{|1+N^{-1}\vg_j^\top \vR^{(S)}(\vGamma) \vg_j|<c\} \prec 0$
for a constant $c>0$.
Lastly, from the definition of $\vGamma=\vGamma(z,\alpha)$ in \eqref{eq:Gamma}, we have
	\begin{align}
	    z^{-1}\vGamma+\tilde m_{\vK}^{(S)}(\vGamma)\vSigma&= \begin{pmatrix}
		\tilde m_{\vK}^{(S)}(\vGamma)\vSigma_r+(\frac{\alpha}{z}+1)\vI_r& \mathbf{0}\\
		\mathbf{0} & \tilde m_{\vK}^{(S)}(\vGamma)\vSigma_0+\vI_{n-r}\label{eq:decomp_gamma+sigma}
	\end{pmatrix}.
	\end{align}
By \eqref{eq:diff_m} and Lemma \ref{lemma:FAconditions}, we have
	\begin{equation}\label{eq:lowerblockImSigma}
\bI\left\{\norm{\left(\tilde m_{\vK}^{(S)}(\vGamma)\vSigma_0+\vI_{n-r}\right)^{-1}}>C\right\}\prec 0
\end{equation}
	for some constant $C>0$. We have already proved $\bI\{|\tilde
m_{\vK}^{(S)}(\vGamma)|>C_0\} \prec 0$, and applying
$\|\vSigma_r\|\le C $ under Assumption~\ref{assump:G}, we can deduce for the
smallest singular value that
\begin{equation}\label{eq:sigmamin1plusmSigma}
		\sigma_{\min} \left(\tilde
m_{\vK}^{(S)}(\vGamma)\vSigma_r+(\alpha/z+1)\vI_{r}\right)\ge
\frac{|\alpha|}{|z|}-1-|\tilde m_{\vK}^{(S)}(\vGamma)|\norm{\vSigma_r} \geq
c
\end{equation}
	on the event $\{|\tilde m_{\vK}^{(S)}(\vGamma)|\le C_0\}$,
for any $z \in U(\eps)$, $|\alpha| \geq \alpha_0$, and some $\alpha_0,c>0$
depending on $\eps,C_0$.
Thus also
\begin{equation}\label{eq:extended1plusmSigma}
\bI\{\|(z^{-1}\vGamma+m_{\tilde
\vK}^{(S)}(\vGamma)\vSigma)^{-1}\|>C\} \prec 0
\end{equation}
for a constant $C>0$, showing all statements of the lemma.
\end{proof}

\begin{proof-of-lemma}[\ref{lemm:apply_fluctuation}(b)]
	Recalling the form of $\vcR(z,\alpha)$ in \eqref{eq:matrix_block},
the quantity we wish to approximate is
	\begin{align*}
		\begin{pmatrix}
			\vv_1^\top&\mathbf{0}
		\end{pmatrix}\vcR(z,\alpha)\begin{pmatrix}
			\vv_2\\
			\mathbf{0}
		\end{pmatrix}=\vv_1^\top\vR(\vGamma)\vv_2=\vv^\top_1\left(\vK-\vGamma\right)^{-1}\vv_2.
	\end{align*}
	Analogous to \eqref{eq:fundamentalidentity} in the proof of Lemma~\ref{lemma:prelimestimate}, 
	for any matrix $\vB\in \C^{n\times n}$, we have
	\begin{align}
		\Tr \vB=
\Tr (\vK-\vGamma)\vR(\vGamma)\vB={-}\Tr \vR(\vGamma)\vB\vGamma+\frac{1}{N}\sum_{i=1}^N
\frac{\vg_i^\top\vR^{(i)}(\vGamma)\vB\vg_i}{1+N^{-1}\vg_i^\top\vR^{(i)}(\vGamma)\vg_i}.\label{eq:S-M-v1Rv2}
	\end{align}
Applying the definition
$m_{\tilde \vK}(\vGamma)=N^{-1}\Tr \vR(\vGamma) +(1-\gamma_N)(-1/z)$
and the identity \eqref{eq:S-M-v1Rv2} with $\vB=\vI$, we obtain analogously to
\eqref{eq:mnidentity} that
	\begin{align*}
m_{\tilde \vK}(\vGamma)
&=(1-\gamma_N)(-1/z)+\frac{1}{Nz}\Tr \vR(\vGamma)\vGamma
-\frac{1}{N}\Tr (z^{-1}\vGamma-\vI)\vR(\vGamma)\notag\\
&={-}\frac{1}{Nz}\sum_{i=1}^N\frac{1}{1+N^{-1}\vg_i^\top\vR^{(i)}(\vGamma)\vg_i}
-\frac{1}{N}\Tr (z^{-1}\vGamma-\vI)\vR(\vGamma).
	\end{align*}
Then, noting that $z^{-1}\vGamma-\vI$ has rank $r$ and hence
$|N^{-1}\Tr(z^{-1}\vGamma-\vI)\vR(\vGamma)|
\leq \frac{r}{N}\frac{|\alpha|}{|z|}\|\vR(\vGamma)\| \prec N^{-1}$, this gives
\begin{equation}\label{eq:S-M-TrR}
m_{\tilde \vK}(\vGamma)
={-}\frac{1}{Nz}\sum_{i=1}^N \frac{1}{1+N^{-1}\vg_i^\top\vR^{(i)}(\vGamma)\vg_i}
+\SD{N^{-1}}.
\end{equation}

Fixing the unit vectors $\vv_1,\vv_2 \in \R^n$,
let us now choose $\vA=\vv_2\vv_1^\top$ and
$\vB=\vA(z^{-1}\vGamma+m_{\tilde \vK}(\vGamma)\cdot\vSigma )^{-1}$
in \eqref{eq:S-M-v1Rv2}, and define
	\begin{align*}
d_i&=\frac{1}{N}\,\vg_i^\top \vR^{(i)}(\vGamma)\vB\vg_i
	-\frac{1}{N}\Tr \vR(\vGamma)\vB\vSigma\\
&=\frac{1}{N}\,\vg_i^\top \vR^{(i)}(\vGamma)\vA
	\left(z^{-1}\vGamma+m_{\tilde \vK}(\vGamma)\vSigma\right)^{-1}\vg_i
	-\frac{1}{N}\Tr \vR(\vGamma)\vA\left(z^{-1}\vGamma+m_{\tilde
\vK}(\vGamma)\vSigma\right)^{-1}\vSigma.
\end{align*}
Then, combining \eqref{eq:S-M-v1Rv2} and \eqref{eq:S-M-TrR}, we get
	\begin{align}
&\vv_1^\top (z^{-1}\vGamma+m_{\tilde \vK}(\vGamma)\vSigma)^{-1}\vv_2=\Tr
\vB\notag\\
&= {-}\Tr \vR(\vGamma)\vB\vGamma + \Tr
\vR(\vGamma)\vB\vSigma \cdot \left(-zm_{\tilde
\vK}(\vGamma)+\SD{N^{-1}}\right)
+\sum_{i=1}^N\frac{d_i}{1+N^{-1}\vg_i^\top\vR^{(i)}(\vGamma)\vg_i} \notag\\
&= {-}z\cdot\vv_1^\top \vR(\vGamma)\vv_2 +
\sum_{i=1}^N\frac{d_i}{1+N^{-1}\vg_i^\top\vR^{(i)}(\vGamma)\vg_i}+\SD{N^{-1}},\label{eq:J1+J2}
	\end{align} 
where the last equality applies the definition of $\vB$ to combine the first two
terms, and applies also
$|\Tr \vR(\vGamma)\vB\vSigma|
\leq \|(z^{-1}\vGamma+m_{\tilde \vK}(\vGamma)\vSigma)^{-1}\vSigma\vR(\vGamma)\|
\prec 1$ by Lemma~\ref{lemm:bounds_R(z)} to obtain the
$\SD{N^{-1}}$ remainder.

	Considering a similar decomposition as in Lemma~\ref{lemma:prelimestimate}, we define
	$d_i=d_{i,1}+d_{i,2}+d_{i,3}+d_{i,4}$ where
	\begin{equation}\label{eq:di14_gamma}
	\begin{aligned}
	d_{i,1}&=\frac{1}{N}\vg_i^\top
\vR^{(i)}(\vGamma)\vA(z^{-1}\vGamma+m_{\tilde \vK}(\vGamma)\vSigma)^{-1}
	\vg_i-\frac{1}{N}\vg_i^\top
\vR^{(i)}(\vGamma)\vA(z^{-1}\vGamma+m_{\tilde \vK}^{(i)}(\vGamma)\vSigma)^{-1}\vg_i,\\
	d_{i,2}&=\frac{1}{N}\vg_i^\top
\vR^{(i)}(\vGamma)\vA(z^{-1}\vGamma+m_{\tilde \vK}^{(i)}(\vGamma)\vSigma)^{-1}
	\vg_i-\frac{1}{N}\Tr \vSigma
\vR^{(i)}(\vGamma)\vA(z^{-1}\vGamma+m_{\tilde \vK}^{(i)}(\vGamma)\vSigma)^{-1},\\
	d_{i,3}&=\frac{1}{N}\Tr \vSigma
\vR^{(i)}(\vGamma)\vA(z^{-1}\vGamma+m_{\tilde \vK}^{(i)}(\vGamma)\vSigma)^{-1}
	-\frac{1}{N}\Tr \vSigma \vR(\vGamma)\vA(z^{-1}\vGamma+m_{\tilde
\vK}^{(i)}(\vGamma)\vSigma)^{-1},\\
	d_{i,4}&=\frac{1}{N}\Tr \vSigma \vR(\vGamma)\vA(z^{-1}\vGamma+m_{\tilde
\vK}^{(i)}(\vGamma)\vSigma)^{-1}
	-\frac{1}{N}\Tr \vSigma \vR(\vGamma)\vA(z^{-1}\vGamma+m_{\tilde
\vK}(\vGamma)\vSigma)^{-1}.
	\end{aligned}
	\end{equation}
	For $\vA=\vv_2\vv_1^\top$, by the bound
$\bI\{\|\vR^{(S)}(\vGamma)\|>C_0\} \prec 0$
from Lemma \ref{lemm:bounds_R(z)}, we have for a constant $C>0$ that
	\begin{equation}
	\bI\left\{\norm{\vR^{(S)}(\vGamma)}_F>C\sqrt{N}\right\}\prec 0, \quad \bI\left\{\norm{\vR^{(S)}(\vGamma)\vA}_F>C\right\}\prec 0
	\end{equation} 
	uniformly over $z\in U(\eps)$.
	Then, employing Lemma~\ref{lemm:bounds_R(z)} and the same bounds as
\eqref{eq:di1}--\eqref{eq:di4} from the proof of Lemma \ref{lemma:FAconditions}
(where here, the bounds for $\|\vR^{(i)}\vA\|_F,\|\vR\vA\|_F$ are
improved by a factor of $N^{-1/2}$ because $\vA$ is low-rank),
we conclude that $|d_{i,1}|,|d_{i,3}|,|d_{i,4}|
	\prec N^{-3/2}$ and $|d_{i,2}| \prec N^{-1}$.  
Hence, applying also $1+N^{-1}\vg_i^\top \vR^{(i)}(\vGamma)\vg_i
=1+N^{-1}\Tr \vSigma\vR(\vGamma)+\SD{N^{-1/2}}$ as follows from
\eqref{eq:concnetr_g_i} and the bound \eqref{eq:mdiffestimate},
	\[\sum_{i=1}^N\frac{d_i}{1+N^{-1}\vg_i^\top\vR^{(i)}(\vGamma)\vg_i}
=\frac{1}{1+N^{-1}\Tr \vSigma\vR(\vGamma)}
\cdot \sum_{i=1}^N d_{i,2}+\SD{N^{-1/2}}.\]
	By Lemma~\ref{lemma:fluctuationavg} applied with
$\Psi_N(\vGamma)=C\sqrt{N}$ and $\Phi_N(\vGamma,\vA)=C$ for a constant $C>0$,
we have $|\sum_i d_{i,2}| \prec N^{-1/2} $. Thus the above quantity is of size
$\SD{N^{-1/2}}$, so applying this back to \eqref{eq:J1+J2},
	\begin{align*}
		\vv_1^\top(\vGamma+zm_{\tilde \vK}(\vGamma)\cdot\vSigma
)^{-1}\vv_2+\vv_1^\top\vR(\vGamma)\vv_2 \prec N^{-1/2}.
	\end{align*}
	Finally, from \eqref{eq:diff_m} and Lemma~\ref{lemm:support_G0} 
we have $m_{\tilde \vK}(\vGamma)=\tilde m_{N,0}(z)+\SD{N^{-1}}$, and applying
this above completes the proof.
\end{proof-of-lemma}

\begin{proof-of-lemma}[\ref{lemm:apply_fluctuation2}]
The proof is similar to Lemma \ref{lemm:apply_fluctuation}, replacing
$r$ and $n$ throughout by $r+1$ and $n+1$,
$\vG_r^{(S)}$ by $[\vu^{(S)},\vG_r^{(S)}]$,
$\vSigma$ by $\widetilde\vSigma$, and
$\vR^{(S)}(\vGamma)$ and $\tilde m_{\vK}^{(S)}(\vGamma)$ by
\[\vR^{(S)}(\widetilde \vGamma)
=\big([\vu^{(S)},\vG^{(S)}]^\top [\vu^{(S)},\vG^{(S)}]-\widetilde
\vGamma\big)^{-1},
\quad \tilde m_{\vK}^{(S)}(\widetilde \vGamma)
=\frac{1}{N}\Tr \vR^{(S)}(\widetilde
\vGamma)+\Big(1-\frac{n+1}{N}\Big)\left({-}\frac{1}{z}\right).\]
The only difference here is that $\widetilde \vSigma$ is no longer
diagonal, leading to the following minor modifications of the preceding
proof: The bound 
\[\frac{1}{n+1}\Tr\left(\widetilde\vSigma
-\begin{pmatrix} \mathbf{0} & \mathbf{0} & \mathbf{0} \\
\mathbf{0} & \mathbf{0} & \mathbf{0} \\
\mathbf{0} & \mathbf{0} & \vSigma_0 \end{pmatrix}\right)
\vR^{(S)}(\widetilde\vGamma) \prec \frac{1}{N}\]
analogous to \eqref{eq:SigmaRcompare} follows upon noting that
(with $\E[u\vg]^\top=\begin{pmatrix} \E[u\vg_r]^\top & \E[u\vg_0]^\top
\end{pmatrix}$)
\[\widetilde \vSigma-\begin{pmatrix} \mathbf{0} & \mathbf{0} & \mathbf{0} \\
\mathbf{0} & \mathbf{0} & \mathbf{0} \\
\mathbf{0} & \mathbf{0} & \vSigma_0 \end{pmatrix}
=\begin{pmatrix} \E[u^2] & \E[u\vg_r]^\top & \E[u\vg_0]^\top \\
\E[u\vg_r] & \vSigma_r & \mathbf{0} \\
\E[u\vg_0] & \mathbf{0} & \mathbf{0} \end{pmatrix}\]
still is of low rank, with rank at most $r+2$.
Writing as shorthand $\tilde m_{\vK}^{(S)}=\tilde
m_{\vK}^{(S)}(\widetilde \vGamma)$, the bound
	\[\bI\{\|(z^{-1}\widetilde\vGamma+\tilde m_{\vK}^{(S)}
\widetilde\vSigma)^{-1}\|>C_0\} \prec 0\]
analogous to \eqref{eq:extended1plusmSigma} follows from
	\[(z^{-1}\widetilde\vGamma+\tilde m_{\vK}^{(S)}\widetilde\vSigma)^{-1}=~\begin{pmatrix}
        \tilde m_{\vK}^{(S)}\E[u^2]
+\frac{\alpha}{z}+1 & \tilde m_{\vK}^{(S)}
\E[u\vg_r]^\top & \tilde m_{\vK}^{(S)} \E[u\vg_0]^\top \\
\tilde m_{\vK}^{(S)} \E[u\vg_r] &
\tilde m_{\vK}^{(S)}\vSigma_r
+(\frac{\alpha}{z}+1)\vI_r & \mathbf{0} \\
	\tilde m_{\vK}^{(S)} \E[u\vg_0] & \mathbf{0} &
\tilde m_{\vK}^{(S)}\vSigma_0+\vI
\end{pmatrix}^{-1},\]
the bound $\bI\{\|\tilde m_{\vK}^{(S)}\vSigma_0+\vI\|^{-1}>C\} \prec 0$ for
the lower-right block as follows from
\eqref{eq:lowerblockImSigma}, and the bound for the inverse of
its Schur-complement
\begin{align*}
&\bI\Bigg\{\Bigg\|\Bigg[\begin{pmatrix}
 \tilde m_{\vK}^{(S)}\E[u^2]
+\frac{\alpha}{z}+1 & \tilde m_{\vK}^{(S)}\E[u\vg_r]^\top \\
\tilde m_{\vK}^{(S)}\E[u\vg_r] & \tilde m_{\vK}^{(S)}\vSigma_r
+(\frac{\alpha}{z}+1)\vI_r \end{pmatrix}\\
&\hspace{1in}
-\begin{pmatrix} \tilde m_{\vK}^{(S)}\E[u\vg_0]^\top \\ \mathbf{0} \end{pmatrix}
(\tilde m_{\vK}^{(S)}\vSigma_0+\vI)^{-1}
\begin{pmatrix} \tilde m_{\vK}^{(S)}\E[u\vg_0] & \mathbf{0}
\end{pmatrix}\Bigg]^{-1}
\Bigg\|>C\Bigg\} \prec 0
\end{align*}
which holds uniformly over $z \in U(\eps)$ for any $|\alpha|>\alpha_0$ 
sufficiently large, by an argument analogous to \eqref{eq:sigmamin1plusmSigma}.
The remainder of the proof
is identical to that of Lemma \ref{lemm:apply_fluctuation}, and we omit the
details. 
\end{proof-of-lemma}

\subsection{Analysis of outliers}

Let $\vV_r,\vGamma(z,\alpha),\vcR(z,\alpha),\widetilde
\vcR(z,\alpha)$ be as defined in the preceding section.
Consider the decomposition of
$\widetilde \vcR(z,\alpha)$ as in \eqref{eq:tildeR}
into its blocks of dimensions 1, $n$, and $N$, and define
\begin{align}
\widetilde \cR_{11}(z,\alpha)&:=\begin{pmatrix} 1 \\ 0 \\ 0 \end{pmatrix}^\top
\widetilde \vcR(z,\alpha) \begin{pmatrix} 1 \\ 0 \\ 0 \end{pmatrix}
=\frac{1}{-z-\alpha+\vu^\top \vu-\vu^\top
\vG(\vG^\top\vG-\vGamma(z,\alpha))^{-1} \vG^\top\vu},
\label{eq:R11}\\
\widetilde \vcR_{1V}(z,\alpha)&:=\begin{pmatrix} 1 \\ 0 \\ 0 \end{pmatrix}^\top
\widetilde \vcR(z,\alpha) \begin{pmatrix} 0 \\ \vV_r \\ 0 \end{pmatrix}
={-}\widetilde \cR_{11}(z,\alpha) \cdot
\vu^\top \vG\Big(\vG^\top \vG-\vGamma(z,\alpha)\Big)^{-1}\vV_r,
\label{eq:R1V}
\end{align}
where the second equalities follow from block matrix inversion of
the lower $2 \times 2$ blocks of $\widetilde \vcR(z,\alpha)$,
followed by block matrix inversion of the full matrix $\widetilde
\vcR(z,\alpha)$. Set
	\begin{equation}\label{eq:eigen-vector-iden}
	\vM_{\vK}(z,\alpha) = 
\vI_r+\alpha \begin{pmatrix} \vV_r^\top & \mathbf{0} \end{pmatrix}
\vcR(z,\alpha)\begin{pmatrix} \vV_r \\ \mathbf{0} \end{pmatrix}.
	\end{equation} 

\begin{prop}\label{prop:master}
Fix any $\eps>0$ and any $\alpha \in \R$ sufficiently large that satisfies
Lemmas \ref{lemm:apply_fluctuation} and \ref{lemm:apply_fluctuation2}.
Then on the event $\cE \cap \widetilde \cE$ of
Lemmas \ref{lemm:apply_fluctuation} and \ref{lemm:apply_fluctuation2},
for all sufficiently large $N$,
\begin{enumerate}[label=(\alph*)]
\item $\widehat \lambda
\in U(\eps) \cap \R$ is an eigenvalue of $\vG^\top \vG$ if and only if
$\det \vM_{\vK}(\widehat\lambda,\alpha)=0$, and its multiplicity as an
eigenvalue of $\vG^\top \vG$ equals the dimension of
$\ker\vM_{\vK}(\widehat\lambda,\alpha)$.
\item Let $\widehat\vv\in\R^n$ be a unit eigenvector of $\vG^\top \vG$ (i.e.\
right singular vector of $\vG$)
corresponding to an eigenvalue $\widehat\lambda \in U(\eps) \cap \R$. Then
$\vV_r^\top\widehat\vv$ is a non-zero vector in $\ker
\vM_{\vK}(\widehat\lambda,\alpha)$, and
\begin{align}\label{eq:norm_vvr_v}
\frac{1}{\alpha^2} = \widehat{\vv}^\top \vV_r
\begin{pmatrix} \vV_r \\ \mathbf{0} \end{pmatrix}^\top
\vcR(\widehat\lambda,\alpha) \begin{pmatrix} \vI_n & \mathbf{0} \\
\mathbf{0} & \mathbf{0} \end{pmatrix} 
\vcR(\widehat\lambda,\alpha) \begin{pmatrix} \vV_r \\ \mathbf{0} \end{pmatrix}
\vV_r^\top \widehat{\vv}
\end{align}
For any vector $\vv\in\R^n$, we have
\begin{equation}\label{eq:vvhat}
\vv^\top \widehat\vv+\alpha\begin{pmatrix} \vv^\top &\mathbf{0}
\end{pmatrix}\vcR(\widehat\lambda,\alpha)\begin{pmatrix} \vV_r \\ \mathbf{0} \end{pmatrix}
\vV_r^\top \widehat\vv=0.
\end{equation}
\item Let $\vu$ be as in Theorem \ref{thm:spiked}(c), and let $\hat{\vu} \in \R^N$ be
a unit eigenvector of $\vG\vG^\top$ (i.e.\ left singular vector of $\vG$)
corresponding to the eigenvalue $\widehat\lambda \in U(\eps) \cap \R$. Then
\begin{equation}\label{eq:uuhat}
	\vu^\top \widehat\vu=\frac{\alpha }{\widehat\lambda^{1/2}\,\widetilde
\cR_{11}(\widehat\lambda,\alpha)}\widetilde\vcR_{1V}(\widehat\lambda,\alpha) \vV_r^\top\widehat\vv.
\end{equation}
\end{enumerate}
\end{prop}
\begin{proof}
For part (a), note that if $\widehat\lambda$ is an eigenvalue of $\vG^\top \vG$,
i.e.\ $\widehat\lambda^{1/2}$ is a singular value of $\vG$ with left and
right unit singular vectors $\widehat \vu$ and $\widehat\vv$, then
\[0=\begin{pmatrix} -\widehat\lambda \vI_n & \vG^\top \\ \vG & -\vI_N \end{pmatrix}
\begin{pmatrix} \widehat\vv \\ \widehat\lambda^{1/2}\widehat\vu \end{pmatrix}\]
which implies, for any $\alpha \in \R$,
\[{-}\alpha\begin{pmatrix} \vV_r \\ \mathbf{0} \end{pmatrix}
\cdot \vV_r^\top \widehat\vv=\begin{pmatrix} -\widehat\lambda\vI_n-\alpha\vV_r\vV_r^\top & \vG^\top \\ \vG & -\vI_N \end{pmatrix}
\begin{pmatrix} \widehat\vv \\ \widehat\lambda^{1/2}\widehat\vu \end{pmatrix}.\]
Fixing $\alpha \in \R$ large enough, on the event $\cE$
of Lemma \ref{lemm:apply_fluctuation}, the generalized resolvent
\[\vcR(\widehat\lambda,\alpha)=\begin{pmatrix} -\widehat\lambda\vI_n-\alpha\vV_r\vV_r^\top & \vG^\top \\ \vG & -\vI_N
\end{pmatrix}^{-1}\]
exists, and multiplying both sides by $\vcR(\widehat\lambda,\alpha)$ gives
\begin{equation}\label{eq:eigequation}
\begin{pmatrix} \widehat\vv \\ \widehat\lambda^{1/2}\widehat\vu \end{pmatrix}
=-\alpha \vcR(\widehat\lambda,\alpha)\begin{pmatrix}\vV_r \\ \mathbf{0} \end{pmatrix}
\cdot \vV_r^\top \widehat\vv.
\end{equation}
Then, multiplying by $(\vV_r^\top\;\mathbf{0})$ on both sides and re-arranging, we get
$\vM_{\vK}(\widehat\lambda,\alpha) \cdot \vV_r^\top \widehat \vv=0$.

We remark
that if $\widehat\lambda$ is
an eigenvalue of multiplicity $k$, and $\vG$ has corresponding linearly
independent left singular vectors $\widehat\vu_1,\ldots,\widehat\vu_k$
and right singular vectors $\widehat\vv_1,\ldots,\widehat\vv_k$, then the
vectors $\{(\widehat \vv_j,\widehat\lambda^{1/2}\widehat \vu_j)\}_{j=1}^k$ on the left side
of \eqref{eq:eigequation} are linearly independent, implying that the vectors
$\{\vV_r^\top \widehat\vv_j\}_{j=1}^k$ on the right side must also be
(non-zero and) linearly independent vectors in $\ker \vM_{\vK}(\widehat\lambda,\alpha)$.
Conversely, if $\{\vy_j\}_{j=1}^k$ are linearly independent vectors in
$\ker \vM_{\vK}(\widehat\lambda,\alpha)$, then defining
\[\begin{pmatrix} \widehat\vv_j \\ \widehat\lambda^{1/2}\widehat \vu_j \end{pmatrix}
=-\alpha \vcR(\widehat\lambda,\alpha)\begin{pmatrix} \vV_r \\ \mathbf{0} \end{pmatrix}
\cdot \vy_j\]
and multiplying by $(\vV_r^\top\; \mathbf{0})$, we must have
$\vV_r^\top \widehat \vv_j=({-}\vM_{\vK}(\widehat\lambda,\alpha)+\vI)\vy_j=\vy_j$. Thus the
pairs $(\widehat\vv_j,\widehat\lambda^{1/2}\widehat\vu_j)$ are linearly independent vectors
satisfying \eqref{eq:eigequation}, and multiplying by
$\vcR(\widehat\lambda,\alpha)^{-1}$ and rearranging shows that
$\widehat\lambda^{1/2}$ must be a singular
value of $\vG$ with multiplicity at least $k$, with corresponding
singular vectors $\{(\widehat\vv_j,\widehat\vu_j)\}_{j=1}^k$.
This establishes part (a).

For part (b), the above argument has shown $\vV_r^\top \widehat \vv \in \ker
\vM_{\vK}(\widehat\lambda,\alpha)$. Multiplying
\eqref{eq:eigequation} on the left by
$$\begin{pmatrix} \vI_n & \mathbf{0} \\
	\mathbf{0} & \mathbf{0} \end{pmatrix}$$
and taking the squared norm (noting that $\widehat\lambda,\alpha$
and $\vcR(\widehat\lambda,\alpha)$ here are real) shows \eqref{eq:norm_vvr_v}.
Multiplying \eqref{eq:eigequation} on the left by $(\vv^\top\;\mathbf{0})$ shows
\eqref{eq:vvhat}. For part (c),
multiplying (\ref{eq:eigequation}) by $(\mathbf{0}\;\vu^\top)$, we have
\begin{align*}
\widehat\lambda^{1/2}\,\vu^\top \widehat\vu&=-\alpha\begin{pmatrix} \mathbf{0}\\ \vu \end{pmatrix}^\top
\vcR(\widehat\lambda,\alpha)\begin{pmatrix} \vV_r \\ \mathbf{0} \end{pmatrix} \cdot \vV_r^\top \widehat\vv
=-\alpha \vu^\top \vG \Big(\vG^\top \vG-\vGamma(\widehat\lambda,\alpha)\Big)^{-1}\vV_r
\cdot \vV_r^\top \widehat\vv
\end{align*}
where the second equality follows from the block matrix inversion of
$\vcR(\widehat\lambda,\alpha)$.
Then, recalling the forms of $\widetilde
\cR_{11}$ and $\widetilde \vcR_{1V}$ from \eqref{eq:R11} and \eqref{eq:R1V},
this gives
\[\vu^\top \widehat\vu=\frac{\alpha \widetilde
\vcR_{1V}(\widehat\lambda,\alpha)}{\widehat\lambda^{1/2}\,\widetilde
\cR_{11}(\widehat\lambda,\alpha)} \cdot \vV_r^\top \widehat\vv\]
which is \eqref{eq:uuhat}.
\end{proof}

For notational convenience, let us now introduce the shorthand
\[\psi_{N,0}(z)=z\tilde m_{N,0}(z), \qquad \psi(z)=z\tilde m(z).\]
By Lemma \ref{lemm:apply_fluctuation}(b) applied with $(\vv_1,\vv_2)$ being the
columns of $\vV_r$,
we see that $\vM_{\vK}(z,\alpha)$ is well-approximated by the (deterministic,
$N$-dependent) matrix
\begin{align}\label{eq:M(z)}
\vM_N(z,\alpha):=\vI_r-\alpha\Big((\alpha+z)\vI_r+\psi_{N,0}(z)
\diag(\lambda_1(\vSigma),\ldots,\lambda_r(\vSigma))\Big)^{-1}.
\end{align} 
To show Theorem \ref{thm:spiked}(a), we translate this approximation into a
comparison of the roots of $0=\det \vM_{\vK}(z,\alpha)$ and $0=\det
\vM_N(z,\alpha)$,
where the latter are explicitly given by $z_{N,0}(-1/\lambda_i(\vSigma))$ for
the function $z_{N,0}(\cdot)$ defined in~\eqref{eq:zN0}.\\



\begin{proof-of-theorem}[\ref{thm:spiked}(a)]
Let us fix any $\eps>0$ and $\alpha \in \R$ satisfying Lemmas
\ref{lemm:apply_fluctuation} and \ref{lemm:apply_fluctuation2}, and denote 
\[f_{N,i}(z,\alpha)=1-\frac{\alpha}{\alpha+z+\psi_{N,0}(z)\lambda_i(\vSigma)}\]
for each $i \in [r]$.
Then $\det\vM_N(z,\alpha)=\prod_{i=1}^r f_{N,i}(z,\alpha)$. Define also the
limiting functions
\[f_i(z,\alpha)=1-\frac{\alpha}{\alpha+z+\psi(z)\lambda_i},
\qquad \vM(z,\alpha)=\vI_r-\alpha\Big((\alpha+z)\vI_r+\psi(z)
\diag(\lambda_1,\ldots,\lambda_r)\Big)^{-1}\]
so $\det \vM(z,\alpha)=\prod_{i=1}^r f_i(z,\alpha)$.
We first analyze the roots of $0=\det \vM(z,\alpha)$: By the
definition $\psi(z)=z\tilde m(z)$, observe that $z \in \R
\setminus \supp{\tilde \mu}$ satisfies $0=\det \vM(z,\alpha)$ if and only if
either $z=0$ or 
\[\tilde m(z)=-1/\lambda_i \text{ for some } i \in [r].\]
(This condition is the same for any non-zero $\alpha \in \R$.)
Let $\cT=\{0\} \cup \{-1/\lambda:\lambda \in \supp{\nu}\}$ be as in
\eqref{eq:zdomain} where $\nu$ is the limit spectral law of $\vSigma_0$.
Then $-1/\lambda_i \in \R \setminus \cT$ for all $i \in [r]$ under Assumption \ref{assum:spike}, so
Proposition~\ref{prop:inv_z} implies that $\tilde m(z)=-1/\lambda_i$ holds
for some $z \in \R \setminus \supp{\tilde \mu}$ if and only if
$z'(-1/\lambda_i)>0$, i.e.\ $i \in \cI$. If $i \in \cI$, then
$\tilde m(z)=-1/\lambda_i$ holds for $z=z(-1/\lambda_i)$,
and we must have $z(-1/\lambda_i)>0$ strictly because for any $z \leq 0$,
we have $\tilde m(z)>0$
(and hence $\tilde m(z) \neq -1/\lambda_i$)
by the definition $\tilde m(z)=\int \frac{1}{x-z}d\tilde \mu(x)$.
Thus the roots of $0=\det \vM(z,\alpha)$ in $\R \setminus \cS=\R \setminus
(\supp{\tilde \mu} \cup \{0\})$ --- and hence
in $U(\eps) \cap \R$ for any sufficiently small $\eps>0$ --- are given
precisely by
\[z_i:=z(-1/\lambda_i) \text{ for } i \in \cI.\]
Since $\tilde m'(z)=\int \frac{1}{(x-z)^2}d\tilde \mu(x)>0$
for all $z \in \R \setminus \cS$, and $\{\lambda_i:i \in \cI\}$ are
distinct by assumption, these values $\{z_i:i \in \cI\}$ are simple roots of
$0=\det \vM(z_i,\alpha)$. Then
$(\det \vM)'(z_i,\alpha) \neq 0$ where $(\det \vM)'$ denotes the derivative in
$z$.

Lemma \ref{lemma:supportcompareb}(c) implies
$\tilde m_{N,0}(z) \to \tilde m(z)$ and $\tilde m_{N,0}'(z) \to \tilde m'(z)$
uniformly over $z \in U(\eps)$. Since also $\lambda_i(\vSigma) \to \lambda_i$,
we have $\det \vM_N(z,\alpha) \to \det \vM(z,\alpha)$ and
$(\det \vM_N)'(z,\alpha) \to (\det \vM)'(z,\alpha)$ uniformly over
$z \in U(\eps)$. This, together with the above condition
$(\det \vM)'(z_i,\alpha) \neq 0$, imply
that for all large $N$, the roots $z_{N,i} \in U(\eps) \cap \R$
of $0=\det \vM_N(z,\alpha)$
are in 1-to-1 correspondence with, and converge to, the above roots $z_i \in
U(\eps) \cap \R$ of
$0=\det \vM(z,\alpha)$. We note that $0=\det\vM_N(z,\alpha)$ if and only if
either $z=0$ or
\begin{equation}\label{eq:equation_z_i}
\tilde{m}_{N,0}(z)=-1/\lambda_i(\vSigma) \text{ for some } i \in [r].
\end{equation}
For each $i \in \cI$, we have $\lambda_i(\vSigma) \to \lambda_i$ where
$z'(-1/\lambda_i)>0$. Recall from Lemma \ref{lemma:supportcompareb}(a)
that $z_{N,0}(\tilde m) \to z(\tilde m)$ and $z_{N,0}'(\tilde m) \to z'(\tilde m)$
uniformly over compact subsets of $\R \setminus \cT$.
Then $z_{N,0}'(-1/\lambda_i(\vSigma)) \to z'(-1/\lambda_i)$, so
also $z_{N,0}'(-1/\lambda_i(\vSigma))>0$ for all large $N$. Then
Proposition~\ref{prop:inv_z} implies that \eqref{eq:equation_z_i} holds
for $z_{N,i}:=z_{N,0}(-1/\lambda_i(\vSigma))$. We have
$z_{N,i} \to z_i=z(-1/\lambda_i)$, so these must be the
roots of $\det \vM_N(z,\alpha)$ in $U(\eps) \cap \R$. Thus we have shown that
for any sufficiently small $\eps>0$ and all large $N$,
the roots $z \in U(\eps) \cap \R$ of $0=\det \vM_N(z,\alpha)$ are precisely
the values
\[z_{N,i}:=z_{N,0}(-1/\lambda_i(\vSigma)) \text{ for } i \in \cI,\]
and $z_{N,i} \to z_i>0$ for each $i \in \cI$.

Finally, we apply Lemma \ref{lemm:apply_fluctuation}(b) with $(\vv_1,\vv_2)$
being the columns of $\vV_r$. On the event $\cE$ of
Lemma \ref{lemm:apply_fluctuation}(a), we have
\begin{equation}\label{eq:MKlipschitz}
\norm{\vM_{\vK}(z,\alpha)} \leq C, \qquad
	\norm{\vM_{\vK}(z,\alpha)-\vM_{\vK}(z',\alpha)} \leq C|z-z'|
\end{equation}
for some $C>0$ and all $z,z' \in U(\eps/2)$.
Also $|\tilde{m}_{N,0}(z)|,|\tilde{m}_{N,0}'(z)|<C$
for a constant $C>0$, all $z \in U(\eps)$, and all large $N$, and thus
\begin{equation}\label{eq:MNlipschitz}
\norm{\vM_N(z,\alpha)} \leq C, \qquad
\norm{\vM_N(z,\alpha)-\vM_N(z',\alpha)} \leq C|z-z'|
\end{equation}
for some $C>0$ and all $z,z' \in U(\eps/2)$.
Then, applying Lemma \ref{lemm:apply_fluctuation}(b) and the Lipschitz
bounds of \eqref{eq:MKlipschitz} and \eqref{eq:MNlipschitz} to take
a union bound over a sufficiently fine covering net of $U(\eps/2)$, we get
\begin{equation}\label{eq:MKuniformapprox}
\sup_{z \in U(\eps/2)} \norm{\vM_N(z,\alpha)-\vM_{\vK}(z,\alpha)} \prec 1/\sqrt{N}.
\end{equation}
Applying also the first bounds of \eqref{eq:MKlipschitz}
and \eqref{eq:MNlipschitz}, this gives
\begin{equation}\label{eq:detMcompare}
\sup_{z \in U(\eps/2)} |\det \vM_N(z,\alpha)-\det \vM_{\vK}(z,\alpha)|
\prec 1/\sqrt{N}.
\end{equation}
Since $\det \vM_N(z,\alpha)$ and $\det \vM_{\vK}(z,\alpha)$ are both holomorphic
over $z \in U(\eps/2)$ on this event $\cE$,
the Cauchy integral formula then implies
\[\sup_{z \in U(\eps)} |(\det \vM_N)'(z,\alpha)-(\det \vM_{\vK})'(z,\alpha)|
\prec 1/\sqrt{N}.\]
In particular, combining with the uniform convergence statements
$\det \vM_N(z,\alpha) \to \det \vM(z,\alpha)$ and $(\det \vM_N)'(z,\alpha) \to
(\det \vM)'(z,\alpha)$ over $z \in U(\eps)$ as argued above,
this shows that on an event $\cE$ satisfying $\bI\{\cE^c\} \prec
0$ and for some $\delta_N \to 0$, we have
\[\sup_{z \in U(\eps) \cap \R} |\det \vM(z,\alpha)-\det \vM_{\vK}(z,\alpha)|,
|(\det \vM)'(z,\alpha)-(\det \vM_{\vK})'(z,\alpha)|<\delta_N.\]
Thus, on this event $\cE$ and as $N \to \infty$, the roots $\widehat \lambda_i
\in U(\eps) \cap \R$ of $0=\det
\vM_{\vK}(z,\alpha)$ are also in 1-to-1 correspondence with, and converge to,
the roots $z_i \in U(\eps) \cap \R$ of $0=\det \vM(z,\alpha)$.
Furthermore, the condition $(\det \vM)'(z_i,\alpha) \neq 0$ implies that
$|(\det \vM_N)'(z,\alpha)|$ and $|(\det \vM_{\vK})'(z,\alpha)|$ are bounded away
from 0 in a neighborhood of each such root $z_i$, so \eqref{eq:detMcompare}
then implies that the corresponding roots $\widehat \lambda_i$ and $z_{N,i}$ of
$0=\det \vM_{\vK}(z,\alpha)$ and $0=\det \vM_N(z,\alpha)$ satisfy
\[|\widehat \lambda_i-z_{N,i}| \prec 1/\sqrt{N}.\]

Proposition \ref{prop:master} shows that on this event $\cE$,
these roots $\{\widehat \lambda_i:i \in \cI\}$ are
precisely the eigenvalues of $\vG^\top \vG$ in $U(\eps) \cap \R$.
By the definition of $\vM(z_i,\alpha)$,
each root $z_i$ of $\det \vM(z_i,\alpha)$ is such that $\ker \vM(z_i,\alpha)$
has dimension 1. Since $\bI\{\cE\}(\widehat\lambda_i-z_i) \to 0$, we have
$\bI\{\cE\}\|\vM_{\vK}(\widehat \lambda_i,\alpha)-\vM(z_i,\alpha)\| \to 0$,
so $\ker \vM_{\vK}(\widehat \lambda_i,\alpha)$ also has dimension 1 on this
event $\cE$ for all large $N$. Then Proposition \ref{prop:master} implies that
the eigenvalues $\{\widehat \lambda_i:i \in \cI\}$ of $\vG^\top \vG$ are
simple, and thus in 1-to-1 correspondence with $\{\lambda_i:i \in \cI\}$.
This proves part (a) of the theorem.
 
\end{proof-of-theorem}

\begin{lemma}\label{lemm:R(z,alph)^2}
Under the assumptions of Theorem
\ref{thm:spiked}, for any fixed $\eps>0$, there exists $\alpha_0>0$ such
that fixing any $\alpha \in \C$ with $|\alpha|>\alpha_0$,
uniformly over $z \in U(\eps)$,
\begin{multline*}
		\Bigg\|  \begin{pmatrix} \vV_r \\ \mathbf{0} \end{pmatrix}^\top
		\vcR(z,\alpha) 
		\begin{pmatrix} \vI_n & \mathbf{0}  \\
		\mathbf{0}   & \mathbf{0} 
		\end{pmatrix} 
		\vcR(z,\alpha) \begin{pmatrix} \vV_r \\ \mathbf{0} \end{pmatrix} \\
		 -\left((\alpha+z)\vI_r+\psi_{N,0}(z)\vSigma_r
		\right)^{-2}\left(\vI_r+\psi_{N,0}'(z)\vSigma_r\right) \Bigg\|\prec\frac{1}{\sqrt{N}}.
\end{multline*}

\end{lemma}
\begin{proof}
Fix any $\alpha \in \C$ satisfying Lemma \ref{lemm:apply_fluctuation}, and
denote
\[f_N(z,\alpha):=\begin{pmatrix}
	\vV_r^\top &\mathbf{0}
\end{pmatrix}\vcR(z,\alpha)\begin{pmatrix}
	\vV_r
	\\\mathbf{0}
\end{pmatrix}, \qquad
g_N(z,\alpha):=-\left((\alpha+z)\vI_r+\psi_{N,0}(z)\vSigma_r  \right)^{-1}.\]
Applying Lemma \ref{lemm:apply_fluctuation}(b) and
the Lipschitz continuity statements of
\eqref{eq:MKlipschitz} and \eqref{eq:MNlipschitz} to take a union bound over a
sufficiently fine covering net of $U(\eps/2)$, we have
\[\sup_{z \in U(\eps/2)} \|f_N(z,\alpha)-g_N(z,\alpha)\| \prec 1/\sqrt{N}.\]
Then by the Cauchy integral formula,
$\sup_{z \in U(\eps)} \|f_N'(z,\alpha)-g_N'(z,\alpha)\| \prec 1/\sqrt{N}$
where $f_N'$ and $g_N'$ denote the entrywise derivatives in $z$. 
The lemma follows, since differentiating $\vcR(z,\alpha)$ in
\eqref{eq:matrix_block} shows
\[f_N'(z,\alpha)=\begin{pmatrix}
	\vV_r^\top &\mathbf{0}	
\end{pmatrix}\vcR(z,\alpha)
\begin{pmatrix} \vI & \mathbf{0} \\ \mathbf{0} & \mathbf{0} \end{pmatrix}
\vcR(z,\alpha)
\begin{pmatrix}
	\vV_r \\\mathbf{0},
\end{pmatrix}\]
while $g_N'(z,\alpha) = ((\alpha+z)\vI_r+\psi_{N,0}(z)\vSigma_r)^{-2}
(\vI_r+\psi_{N,0}'(z)\vSigma_r)$.
\end{proof}

\begin{proof-of-theorem}[\ref{thm:spiked}(b)]
Let $\widehat \vv_i$ be the given unit-norm eigenvector of $\vK$ with eigenvalue
$\widehat \lambda_i$. Let $z_{N,i}=z_{N,0}(-1/\lambda_i(\vSigma))$ and
$z_i=z(-1/\lambda_i)$.
Then, fixing any $\alpha \in \R$ large enough to satisfy Lemmas
\ref{lemm:apply_fluctuation} and \ref{lemm:apply_fluctuation2},
Proposition~\ref{prop:master}(b) shows that
	$\vV_r^\top\widehat{\vv}_i \in\ker
\vM_{\vK}(\widehat\lambda_i,\alpha)$.
By \eqref{eq:MKuniformapprox}, \eqref{eq:MNlipschitz}, and the bound $|\widehat
\lambda_i-z_{N,i}| \prec N^{-1/2}$ of part (a) of the theorem already proven,
we have
	\begin{align}
	\norm{\vM_{\vK}(\widehat\lambda_i,\alpha)-\vM_N(z_{N,i},\alpha)} &\le
\norm{\vM_{\vK}(\widehat\lambda_i,\alpha)-\vM_N(\widehat
\lambda_i,\alpha)}+\norm{\vM_N(\widehat\lambda_i,\alpha)-\vM_N(z_{N,i},\alpha)}
\nonumber\\
&\prec N^{-1/2}.\label{eq:MKNlambdacompare}
	\end{align}
Let $\vv_1,\ldots,\vv_r$ denote the columns of $\vV_r$, which are the unit
eigenvectors of $\vSigma$.
	Then, applying $\vV_r^\top\widehat{\vv}_i \in\ker
\vM_{\vK}(\widehat\lambda_i,\alpha)$, \eqref{eq:MKNlambdacompare}, and
the definition of $\vM_N(z,\alpha)$, and noting that
$\psi_{N,0}(z_{N,i})=z_{N,i}\tilde
m_{N,0}(z_{N,i})=-z_{N,i}/\lambda_i(\vSigma)$, we have
\[\norm{\vM_N(z_{N,i},\alpha) \cdot \vV_r^\top\widehat{\vv}_i}^2
=\sum_{j=1}^r
\left(1-\frac{\alpha}{\alpha+z_{N,i}(1-\lambda_j(\vSigma)/\lambda_i(\vSigma))}
\right)^2 (\vv_j^\top \widehat\vv_i)^2 \prec 1/N.\]
For each $j \in [r] \setminus \{i\}$, we have that
$z_{N,i}(1-\lambda_j(\vSigma)/\lambda_i(\vSigma))$ is bounded away from 0
as $N \to \infty$ because $z_{N,i} \to z_i>0$ and
$\lambda_j(\vSigma)/\lambda_i(\vSigma) \to \lambda_j/\lambda_i \neq 1$.
So this implies
	\begin{equation}\label{eq:y_q_j}
	|\vv^\top_j\widehat{\vv}_i|^2\prec 1/N\quad\text{ for all }\quad j\in [r]\setminus \{i\}.
	\end{equation}
	At the same time, applying Lemma~\ref{lemm:R(z,alph)^2} 
and $|\widehat \lambda_i-z_{N,i}| \prec N^{-1/2}$ to
bound~\eqref{eq:norm_vvr_v} in Proposition~\ref{prop:master}(b), we have 
\begin{align}
\frac{1}{\alpha^2}
&=\widehat\vv_i^\top \vV_r
\begin{pmatrix} \vV_r \\ \mathbf{0} \end{pmatrix}^\top
\vcR(z_{N,i},\alpha) \begin{pmatrix} \vI_n & \mathbf{0} \\
\mathbf{0} & \mathbf{0} \end{pmatrix} 
\vcR(z_{N,i},\alpha) \begin{pmatrix} \vV_r \\ \mathbf{0} \end{pmatrix}
\vV_r^\top \widehat\vv_i+\SD{N^{-1/2}}\notag\\
&=\widehat\vv_i^\top \vV_r
\Big((\alpha+z_{N,i})\vI_r+\psi_{N,0}(z_{N,i})\vSigma_r\Big)^{-2}
\Big(\vI_r+\psi_{N,0}'(z_{N,i})\vSigma_r\Big) \vV_r^\top \widehat\vv_i
+\SD{N^{-1/2}}\notag\\
&=|\vv^\top_i\widehat{\vv}_i|^2\cdot\frac{1+\psi_{N,0}'(z_{N,i})\lambda_i(\vSigma)}{\alpha^2}\notag\\
&\hspace{1in}+\sum_{j\neq
i} |\vv^\top_j\widehat{\vv}_i|^2
\cdot\frac{1+\psi_{N,0}'(z_{N,i})\lambda_j(\vSigma)}{(\alpha+z_{N,i}(1-\lambda_j(\vSigma)/\lambda_i(\vSigma)))^2}+\SD{N^{-1/2}}\notag\\
&=
|\vv^\top_i\widehat{\vv}_i|^2\cdot\frac{1+\psi_{N,0}'(z_{N,i})\lambda_i(\vSigma)}{\alpha^2}+\SD{N^{-1/2}}, \label{eq:x-y}
	\end{align}
	the last equality applying \eqref{eq:y_q_j}.
Observe that
\begin{align*}
1+\psi_{N,0}'(z_{N,i})\lambda_i(\vSigma)
&=1+z_{N,i}\tilde m_{N,0}'(z_{N,i})\lambda_i(\vSigma)
+\tilde m_{N,0}(z_{N,i})\lambda_i(\vSigma)\\
&=z_{N,i}\tilde m_{N,0}'(z_{N,i})\lambda_i(\vSigma)
=z_{N,i}\lambda_i(\vSigma)/z_{N,0}'(-1/\lambda_i(\vSigma)),
\end{align*}
where the last two equalities use $z_{N,i}=z_{N,0}(-1/\lambda_i(\vSigma))$
and $\tilde m_{N,0}(\cdot)$ is the inverse function of $z_{N,0}(\cdot)$.
Then, multiplying by $\alpha^2/(1+\psi_{N,0}'(z_{N,i})\lambda_i(\vSigma))$
we obtain
\[|\vv^\top_i\widehat{\vv}_i|^2=\frac{z_{N,0}'(-1/\lambda_i(\vSigma))}
{z_{N,i}\lambda_i(\vSigma)}+\SD{N^{-1/2}}
=\varphi_{N,0}({-}1/\lambda_i(\vSigma))+\SD{N^{-1/2}},\]
where we recall $\varphi_{N,0}$ from \eqref{eq:zN0}.
We have $\varphi_{N,0}(-1/\lambda_i(\vSigma)) \to \varphi(-1/\lambda_i)
=z'(-1/\lambda_i)/(\lambda_iz_i)>0$, so taking a square root gives
\begin{equation}\label{eq:vivhati}
|\vv^\top_i\widehat{\vv}_i|=\sqrt{\varphi_{N,0}({-}1/\lambda_i(\vSigma))}
+\SD{N^{-1/2}}.
\end{equation}

Finally, for any unit vector $\vv \in \R^n$, by \eqref{eq:vvhat} in
Proposition~\ref{prop:master}(b),
Lemma \ref{lemm:apply_fluctuation}(b), and the bound
$|\widehat\lambda_i-z_{N,i}| \prec N^{-1/2}$ in part (a) of the theorem
already shown, we know that 
\begin{align*}
\vv^\top\widehat{\vv}_i =~& {-}\alpha \cdot \begin{pmatrix} \vv^\top & \mathbf{0}
\end{pmatrix}
\vcR(z_{N,i},\alpha)\begin{pmatrix} \vV_r \\ \mathbf{0} \end{pmatrix}
\cdot \vV_r^\top\widehat{\vv}_i+\SD{N^{-1/2}}\\
=~& {-}\alpha
\sum_{j=1}^r\frac{\vv^\top\vv_j\cdot\vv_j^\top\widehat{\vv}_i}{\alpha+z_{N,i}+\psi_{N,0}(z_{N,i})\lambda_j(\vSigma)}+\SD{N^{-1/2}} \\
=~& {-}\alpha \sum_{j=1}^r
\frac{\vv^\top\vv_j\cdot\vv_j^\top\widehat{\vv}_i}{\alpha+z_{N,i}\cdot(1-\lambda_j(\vSigma)/\lambda_i(\vSigma))}+\SD{N^{-1/2}}.
\end{align*}
Applying \eqref{eq:y_q_j} and \eqref{eq:vivhati}, only the summand with $j=i$
contributes, and we obtain as desired
\[|\vv^\top \widehat\vv_i|=\sqrt{\varphi_{N,0}({-}1/\lambda_i(\vSigma))}
\cdot |\vv^\top \vv_i|+\SD{N^{-1/2}}.\]
\end{proof-of-theorem}

\begin{proof-of-theorem}[\ref{thm:spiked}(c)]
Applying Lemma~\ref{lemm:apply_fluctuation2}(b) and block matrix inversion of
$\widetilde \vGamma+\psi_{N,0}(z)\widetilde \vSigma$ to the definitions of
$\widetilde \cR_{11}$ and $\widetilde \vcR_{1V}$ in \eqref{eq:R11} and
\eqref{eq:R1V}, we have
\begin{align*}
\left|\widetilde
\cR_{11}(z,\alpha)+\left(z+\alpha+\psi_{N,0}(z)\cdot\E[u^2]-\psi_{N,0}(z)^2
\cdot \E[u\vg]^\top\left(\vGamma+\psi_{N,0}(z)\vSigma\right)^{-1}\E[u\vg]\right)^{-1}\right|
&\prec \frac{1}{\sqrt{N}},\\
\left\|\widetilde \vcR_{1V}(z,\alpha)-\frac{\psi_{N,0}(z) \cdot
\E[u\vg]^\top\left(\vGamma+\psi_{N,0}(z)\vSigma
\right)^{-1}\vV_r}{z+\alpha+\psi_{N,0}(z)\cdot\E[u^2]-\psi_{N,0}(z)^2
\cdot \E[u\vg]^\top\left(\vGamma+\psi_{N,0}(z)\vSigma
\right)^{-1}\E[u\vg]}\right\| &\prec \frac{1}{\sqrt{N}}.
\end{align*}
Hence, 
\begin{align*}
	\left\|\frac{ \widetilde \vcR_{1V}(z,\alpha)}{ \widetilde
\cR_{11}(z,\alpha)}+\psi_{N,0}(z) \cdot
\E[u\vg]^\top\vV_r\cdot\left((\alpha+z)\vI_r+\psi_{N,0}(z)\vSigma_r
\right)^{-1}\right\| &\prec \frac{1}{\sqrt{N}}.
\end{align*} 
Applying this and the bound
$|\widehat\lambda_i-z_{N,i}| \prec N^{-1/2}$
to Proposition~\ref{prop:master}(c),
	\begin{align*}
		\vu^\top\widehat{\vu}_i=
\frac{\alpha}{\widehat\lambda_i^{1/2}}
\frac{\widetilde \vcR_{1V}(\widehat\lambda_i,\alpha)}
{\widetilde \cR_{11}(\widehat\lambda_i,\alpha)} \cdot \vV_r^\top \widehat\vv_i
&=-\frac{\alpha}{\sqrt{z_{N,i}}}
\sum_{j=1}^r\frac{\psi_{N,0}(z_{N,i}) \cdot
\E[u\vg]^\top\vv_j\cdot\vv_j^\top\widehat{\vv}_i}{\alpha+z_{N,i}+\psi_{N,0}(z_{N,i})\lambda_j(\vSigma)}+\SD{N^{-1/2}}\\
&=-\frac{\alpha}{\sqrt{z_{N,i}}}
\sum_{j=1}^r\frac{\psi_{N,0}(z_{N,i}) \cdot
\E[u\vg]^\top\vv_j\cdot\vv_j^\top\widehat{\vv}_i}{\alpha+z_{N,i}(1-\lambda_j(\vSigma)/\lambda_i(\vSigma))}+\SD{N^{-1/2}}.
	\end{align*}
Then, applying again \eqref{eq:y_q_j} and \eqref{eq:vivhati}, only the summand
with $j=i$ contributes, and this gives
	\begin{align*}
		|\vu^\top\widehat{\vu}_i|=~&
\frac{|\E[u\vg]^\top\vv_i|\cdot |\psi_{N,0}(z_{N,i})|\sqrt{\varphi_{N,0}
(-1/\lambda_i(\vSigma))}}{\sqrt{z_{N,i}}}+\SD{N^{-1/2}}.
	\end{align*}
Recalling $\psi_{N,0}(z_{N,i})=-z_{N,i}/\lambda_i(\vSigma)$, this yields part
(c) of the theorem.
\end{proof-of-theorem} 



%%%%%%%%%%%%%%%%%%%%%%%%%%%%%%%%%%%%%%%%%%%%%%%%%%%%%%%%%%%%%%%%%%%%%%%%%%%%%%%%%%%%%
%%%%%%%%%%%%%%%%%%%%%%%%%%%%%%%%%%%%%%%%%%%%%%%%%%%%%%%%%%%%%%%%%%%%%%%%%%%%%%%%%%%%%
\section{Proofs for propagation of spiked eigenstructure in deep NNs}\label{sec:CK_spike_proof}

We next prove Theorems \ref{thm:no_outlier_ck} and \ref{thm:ck_spike}.
Appendix \ref{subsec:onelayer} first
establishes these results for a one-hidden-layer NN, $L=1$. We then apply this
result for $L=1$ inductively in Appendix \ref{subsec:propagation} to obtain
these results for general $L$. Appendix \ref{subsec:GMM_proof} proves
Corollary \ref{cor:gmm}.

\subsection{Spike analysis for one-hidden-layer CK}\label{subsec:onelayer}

Consider the setup in Section~\ref{sec:spike_CK} with a single hidden
layer $L=1$. In this setting, let us simplify notation and denote
\[\vX=\vX_0, \qquad \vW=\vW_0, \qquad d=d_0, \qquad N=d_1,\]
\[\vY=\vX_1=\frac{1}{\sqrt{N}}\sigma (\vW\vX),
\qquad \vK=\vK_1=\vY^\top \vY.\]
We denote the rows of $\vW$ and columns of $\vX$ respectively by
\[\vw_i^\top \in \R^d \text{ for } i \in [N],
\qquad \vx_\alpha \in \R^d \text{ for } \alpha \in [n].\]
We write $\E_{\vw}$ for
the expectation over a standard Gaussian vector $\vw \sim \cN(0,\vI)$ in $\R^d$.

Note that for a sufficiently large constant $B>0$
(depending on $\supp{\nu}$ and $\lambda_1,\ldots,\lambda_r$),
Assumption \ref{assump:data} implies that the event
\begin{equation}\label{eq:goodevent}
\cE(\vX)=\Big\{\|\vX\|<B,\; |\vx_\alpha^\top \vx_\beta|<\tau_n
\text{ and } |\|\vx_\alpha\|_2-1|<\tau_n
\text{ for all } \alpha \neq \beta \in [n]\Big\}
\end{equation}
holds almost surely for all large $n$. We will use throughout this
section the following argument: Since $\vW \equiv \vW^{(n)}$ is independent
of $\vX \equiv \vX^{(n)}$, and $\cE(\vX^{(n)})$ holds for all large $n$ with
probability 1 over $\{\vX^{(n)}\}_{n=1}^\infty$, to prove any almost-sure
statement, it suffices to show that the statement holds with
probability 1 over $\{\vW^{(n)}\}_{n=1}^\infty$, for any deterministic
matrices $\{\vX^{(n)}\}_{n=1}^\infty$ satisfying $\cE(\vX^{(n)})$.
Thus, we assume in the remainder of this section that
$\vX$ is deterministic and satisfies $\cE(\vX)$
for all large $n$, and write $\E,\P$ for the
expectation and probability over only the random weight matrix $\vW$.

We will apply Theorem \ref{thm:spiked} to a centered version of $\vY$,
\[\vG:=\vY-\E \vY=\frac{1}{\sqrt{N}}[\vg_1,\ldots,\vg_N]^\top,
\qquad \vg_i^\top:=\sigma(\vw_i^\top {\vX})-\E_{\vw}[\sigma(\vw^\top {\vX})].\]
Note that these rows $\vg_i^\top$ are i.i.d.\ with mean $\mathbf{0}$ and covariance
\begin{equation}\label{eq:Sigma_CK}
	\vSigma:=\E_{\vw}[\sigma(\vw^\top \vX)^\top\sigma(\vw^\top
\vX)]-\E_{\vw}[\sigma(\vw^\top \vX)]^\top\E_{\vw}[\sigma(\vw^\top
\vX)]\in\R^{n\times n}.
\end{equation}

\begin{lemma}\label{lemm:expect_rows}
Suppose Assumptions \ref{assump:NNasymptotics}, \ref{assump:data}, and
\ref{assump:sigma} hold, with $L=1$ and deterministic $\vX$. Then
\[\|\E_{\vw}[\sigma(\vw^\top \vX)]\|_2 \to 0, \qquad \|\E\vY\| \to 0.\]
\end{lemma}
\begin{proof}
Denote $\xi \sim \cN(0,1)$. Applying $\E[\sigma(\xi)]=0$,
$\E[\sigma'(\xi)\xi]=\E[\sigma''(\xi)]=0$, and a
Taylor approximation of $\sigma$, for any $\alpha \in [n]$,
	\begin{align*}
	\E_{\vw}[\sigma(\vw^\top\vx_\alpha)]=~&\E[\sigma(\norm{\vx_\alpha}\xi)]-\E[\sigma(\xi)]\\
	= ~&
\E[\sigma'(\xi)\xi(\norm{\vx_\alpha}-1)]+\E[\sigma''(\eta)\xi^2(\norm{\vx_\alpha}-1)^2]=\E[\sigma''(\eta)\xi^2(\norm{\vx_\alpha}-1)^2]
	\end{align*}
for some $\eta$ between $\xi$ and $\norm{\vx_i}\xi$.
Then, applying $|\sigma''(x)| \leq \lambda_\sigma$ and the
$\tau_n$-orthonormality of $\vX$ under $\cE(\vX)$,
\[|\E_{\vw}[\sigma(\vw^\top\vx_\alpha)]| \leq \lambda_\sigma \tau_n^2.\]
This gives $\|\E_{\vw}[\sigma(\vw^\top\vX)]\|_2 \leq \lambda_\sigma\tau_n^2\sqrt{n} \to 0$,
so also
$\norm{\E \vY}=\norm{\frac{1}{\sqrt{N}}\bI_N \cdot
\E_{\vw}[\sigma(\vw^\top \vX)]} \to 0$.
\end{proof}

Next, we recall from \cite{wang2021deformed} 
an approximation of $\vSigma$ by the linearized matrix
\begin{equation}\label{def:Phi0}
	\vSigma_{\mathrm{lin}}:=b_\sigma^2\vX^\top \vX+(1-b_\sigma^2 )\vI_n
\end{equation}
in the operator norm.

\begin{lemma}\label{lemm:phi_0}
Suppose Assumptions \ref{assump:NNasymptotics}, \ref{assump:data}, and
\ref{assump:sigma} hold, with $L=1$ and deterministic $\vX$.
\[\|\vSigma-\vSigma_{\mathrm{lin}}\| \to 0.\] 
Consequently, ordering $\lambda_1(\vSigma),\ldots,\lambda_n(\vSigma)$
in the same order as
$\lambda_1(\vX^\top\vX),\ldots,\lambda_n(\vX^\top\vX)$,
\begin{equation}\label{eq:lambda_i_diff}
	\sup_{i \in [n]} \Big|b_\sigma^2\lambda_i(\vX^\top \vX)
+(1-b_\sigma^2)-\lambda_i(\vSigma)\Big| \to 0.
\end{equation}
\end{lemma}
\begin{proof}
Denote $\xi \sim \cN(0,1)$. Let $\zeta_k(\sigma)=\E[\sigma(\xi)h_k(\xi)]$ be
the $k$-th Hermite coefficient of $\sigma$, where $h_k(x)$ is the $k$-th
Hermite polynomial normalized so that $\E[h_k(\xi)^2]=1$. Note that by Gaussian
integration by parts and the assumption $\E[\sigma''(\xi)]=0$,
\[\zeta_1(\sigma)=\E[\xi \sigma(\xi)]
=\E[\sigma'(\xi)]=b_\sigma,
\qquad \sqrt{2}\,\zeta_2(\sigma)=\E[(\xi^2-1)\sigma(\xi)]
=\E[\xi\sigma'(\xi)]=\E[\sigma''(\xi)]=0.\]
Then by \cite[Lemma 5.2]{wang2021deformed} and the first statement
of Lemma~\ref{lemm:expect_rows}, we have
\[\|\vSigma_0-\vSigma\|
\leq \|\vSigma_0-\E_{\vw}[\sigma(\vw^\top \vX)^\top \sigma(\vw^\top \vX)]\|
+\|\E_{\vw}[\sigma(\vw^\top \vX)^\top \sigma(\vw^\top \vX)]-\vSigma\| \to
0\]
where
	\[{\vSigma}_0=\zeta_1(\sigma)^2 \vX^\top \vX+  \zeta_3(\sigma)^2(
\vX^\top  \vX)^{\odot 3}+(1-\zeta_1(\sigma)^2-\zeta_3(\sigma)^2)\vI_n.\]
(Here, examination of the proof of \cite[Lemma 5.2]{wang2021deformed} shows
that the condition $\sum_\alpha (\|\vx_\alpha\|_2-1)^2 \leq B^2$ for
$(\eps,B)$-orthonormality is not
used when $\zeta_2(\sigma)=0$, and the remaining conditions of
$(\eps,B)$-orthonormality hold under $\cE(\vX)$.)
The lemma then follows upon observing that under $\cE(\vX)$,
\begin{align*}
	\norm{( \vX^\top  \vX)^{\odot 3}-\vI_n}\le~& \norm{\diag(( \vX^\top  \vX)^{\odot 3}-\vI_n)} +\norm{\offdiag(\vX^\top  \vX)^{\odot 3}}_F\\
	\le ~& \max_{\alpha\in [n]} \left|\norm{
\vx_\alpha}^6-1\right|+n\cdot\max_{\alpha\neq\beta\in[n]}|\vx_\alpha^\top\vx_\beta|^3
\leq C(\tau_n+n\tau_n^3),
\end{align*}
so that $\norm{\vSigma_{\mathrm{lin}}-\vSigma_0}=\zeta_3(\sigma)^2
\norm{( \vX^\top  \vX)^{\odot 3}-\vI_n} \to 0$ when $\lim_{n \to \infty} \tau_n
\cdot n^{1/3}=0$.
\end{proof}

Theorem \ref{thm:spiked} will provide a characterization
of outlier eigenvalues of $\vK$ that are separated from
$\cS_1=\supp{\mu_1} \cup \{0\}$, which is different from $\supp{\mu_1}$ when
$\gamma_1<1$. For $\gamma_1<1$, we augment this
statement with a small-ball argument to bound the smallest eigenvalue of $\vK$,
using the following result of \cite[Theorem 2.1]{yaskov2016controlling}.

\begin{lemma}[\cite{yaskov2016controlling}]\label{lemma:smallball}
Let $\vG=\frac{1}{\sqrt{N}}[\vg_1,\ldots,\vg_N]^\top \in \R^{N \times n}$ where
the rows $\vg_i \in \R^n$ are i.i.d.\ and equal in law to $\vg \in \R^n$. Define
\[\vSigma=\E\vg\vg^\top,
\qquad c_{\vg}=\inf_{\vv \in \R^n:\|\vv\|_2=1} \E|\vg^\top \vv|,
\qquad L_{\vg}(\delta,\iota)=\sup_{\vPi:\rank{\vPi} \geq \iota n}
\P\Big[|\vPi\vg|^2 \leq \delta\,\rank{\vPi}\Big]\]
where the latter supremum is taken over all orthogonal projections
$\vPi \in \R^{n \times n}$ with rank at least $\iota \cdot n$.

Suppose $\lambda_{\max}(\vSigma) \leq 1$, $c(\vg) \geq c$, and $n/N \leq
y$ for some constants $c>0$ and $y \in (0,1)$. Then there exist constants $s_0,\iota>0$
depending only on $(c,y)$ such that for any $\delta \in (0,1)$ and $s>0$,
\[\P[\lambda_{\min}(\vG^\top \vG) \geq (s_0-L(\delta,\iota;\vg)-s)\delta]
\geq 1-2e^{-yns^2/2}.\]
\end{lemma}

\begin{lemma}\label{lemma:outliersat0}
Suppose Assumptions \ref{assump:NNasymptotics}, \ref{assump:data}, and
\ref{assump:sigma} hold, with $L=1$ and deterministic $\vX$. Let
$\vG=\vY-\E\vY$.
\begin{enumerate}[label=(\alph*)]
\item If $\gamma_1 \geq 1$, then $0 \in \supp{\mu_1}$.
\item If $\gamma_1<1$, then there is a constant $c>0$ such that
$\lambda_{\min}(\vG^\top \vG)>c$ almost surely for all large $n$.
\end{enumerate}
\end{lemma}
\begin{proof}
If $\gamma_1>1$ strictly, then by definition
\[\mu_1=\frac{1}{\gamma_1}\,\tilde \mu_1+\frac{\gamma_1-1}{\gamma_1}\,\delta_0\]
is a mixture of $\tilde \mu_1$ and a point mass at 0, so $0 \in \supp{\mu_1}$.
If $\gamma=1$, then $\mu_1=\tilde \mu_1$. In this case,
recall from Proposition \ref{prop:inv_z}
that $\supp{\tilde \mu_1}$ is characterized by the function
\[z(\tilde m)={-}\frac{1}{\tilde m}+\gamma_1
\int \frac{\lambda}{1+\lambda \tilde m}d\nu_0(\lambda).\]
When $\gamma_1=1$, we have for all $\tilde m \in (0,\infty)$ that
\[z(\tilde m)<0, \qquad
z'(\tilde m)=\frac{1}{\tilde m^2}-\int \frac{\lambda^2}{(1+\lambda \tilde m)^2}
d\nu(\lambda)>0,\]
so $z(\tilde m)$ increases from $-\infty$ to 0 over the positive line
$\tilde m \in (0,\infty)$. Suppose by contradiction that
$0 \notin \supp{\tilde \mu}$. Then by
Proposition \ref{prop:inv_z}, there must be a point $\tilde m \in \R \setminus
\cT$ where $z(\tilde m)=0$ and $z'(\tilde m)>0$ strictly, implying that
there is an open interval $(\tilde m_-,\tilde m_+) \ni \tilde m$
on which $z(\cdot)$ increases from $z(\tilde m_-)<0$ to $z(\tilde m_+)>0$.
We must have $\tilde m<0$ by the above behavior of $z(\cdot)$ on
$(0,\infty)$, and the range $[z(\tilde m_-),z(\tilde m_+)]$ must overlap with
$[z(a),z(b)]$ for some sufficiently large $a,b \in (0,\infty)$.
But this contradicts the non-intersecting property
shown in \cite[Theorem 4.4]{silverstein1995analysis}. So
also in this case $0 \in \supp{\tilde \mu_1}=\supp{\mu_1}$, showing part (a).

For part (b), we apply Lemma \ref{lemma:smallball}. Under $\cE(\vX)$,
the condition $b_\sigma \neq 0$ implies
$c_0<\lambda_{\min}(\vSigma_{\mathrm{lin}}) \leq
\lambda_{\max}(\vSigma_{\mathrm{lin}})<C_0$ for some constants $C_0,c_0>0$.
Hence also
\begin{equation}\label{eq:Sigmabounds}
c_0<\lambda_{\min}(\vSigma) \leq \lambda_{\max}(\vSigma)<C_0
\end{equation}
for all large $n$, by Lemma \ref{lemm:phi_0}.
We may assume without loss of generality that
$\lambda_{\max}(\vSigma) \leq 1$ as needed in Lemma \ref{lemma:smallball};
otherwise, the following argument may be applied to a rescaling of $\vSigma$
and $\vG$.

To lower bound $c_{\vg}$ in Lemma \ref{lemma:smallball}, observe that
for any unit vector $\vv \in \R^n$ we have
\[\E[(\vg^\top \vv)^2]=\vv^\top \vSigma\vv>c_0.\]
Viewing $F(\vw)=\vg^\top \vv=\sigma(\vw^\top \vX)\vv=\sum_{\alpha=1}^n v_\alpha
\sigma(\vw^\top \vx_\alpha)$ as a function of $\vw  \sim \cN(0,\vI)$, we have
$\nabla F(\vw)=\sum_{\alpha=1}^n v_\alpha \sigma'(\vw^\top \vx_\alpha)
\vx_\alpha=\vX(\vv \odot \sigma'(\vw^\top \vX))$ where
$\sigma'(\cdot)$ is applied coordinatewise and $\odot$ is the coordinatewise
product. Then, applying $|\sigma'(x)| \leq \lambda_\sigma$, observe that
$\|\nabla F(\vw)\|_2 \leq \|\vX\| \cdot \|\vv \odot \sigma'(\vw^\top \vX)\|_2
\leq \lambda_\sigma \|\vX\|$, so $F(\vw)$ is $C$-Lipschitz in $\vw$
for a constant $C>0$ (not depending on $\vv$) on the event $\cE(\vX)$.
This implies by Gaussian
concentration-of-measure that $\vg^\top \vv$ is sub-gaussian,
i.e.\ for some constants $C,c>0$ and any $t>0$,
$\P[|\vg^\top \vv| \geq t] \leq Ce^{-ct^2}$.
Integrating this tail bound, for some constant $t>0$ sufficiently large,
we have $\E[(\vg^\top \vv)^2\bI_{\{|\vg^\top \vv|>t\}}]
\leq c_0/2$, and hence
\[c_0<\E[(\vg^\top \vv)^2]
\leq \E[(\vg^\top \vv)^2\bI_{\{|\vg^\top \vv| \leq t\}}]+\frac{c_0}{2}
\leq t \cdot \E|\vg^\top \vv|+\frac{c_0}{2}.\]
So $\E|\vg^\top \vv| \geq c_0/(2t)$, and hence $c_{\vg} \geq c_0/(2t)>c$ for a
constant $c>0$.

Now let $s_0,\iota>0$ be the constants depending on $(c,\gamma_1)$ in the
statement of Lemma \ref{lemma:smallball}.
By the nonlinear Hanson-Wright inequality of
\cite[Eq.\ (3.3)]{wang2021deformed}, for any orthogonal projection $\vPi \in
\R^{n \times n}$, any $t>0$, and some constant $c>0$, we have
\[\P\big[|\vg^\top \vPi \vg-\E\vg^\top \vPi \vg|>t\big] \leq
2e^{-c\min(t^2/\|\vPi\|_F^2,t/\|\vPi\|)}.\]
Here $\E \vg^\top \vPi \vg=\Tr \vPi\vSigma>c_0\rank{\vPi}$,
$\|\vPi\|_F^2=\rank{\vPi}$, and $\|\vPi\|=1$,
so applying this with $t=(c_0/2)\rank{\vPi}$ yields
\[\P[|\vPi\vg|^2 \leq (c_0/2)\rank{\vPi}] \leq 2e^{-c'\rank{\vPi}}.\]
Then, choosing $\delta=c_0/2$, we get $L_{\vg}(\delta,\iota) \to 0$ as $n \to
\infty$. Then Lemma \ref{lemma:smallball} implies $\lambda_{\min}(\vG^\top
\vG)>s_0\delta/2$ almost surely for all large $n$, as desired.
\end{proof}

The following is the main result of this section, showing that Theorems
\ref{thm:no_outlier_ck} and \ref{thm:ck_spike} hold in this setting of $L=1$.

\begin{lemma}\label{lemma:onelayer}
Theorems \ref{thm:no_outlier_ck} and
\ref{thm:ck_spike} hold for a single layer $L=1$.
Furthermore, $\vY$ is $C\tau_n$-orthonormal for some constant $C>0$, 
almost surely for all large $n$.
\end{lemma}
\begin{proof}
We condition on $\vX$ as discussed at the start of this section, and
apply Theorem \ref{thm:spiked} to the centered matrix
$\vG=\vY-\E\vY=\frac{1}{\sqrt{N}}[\vg_1,\ldots,\vg_N]^\top$.
Let us verify Assumption~\ref{assump:G} for $\vG$:
We have shown Assumption~\ref{assump:G}(a) in \eqref{eq:Sigmabounds}.
The rows $\vg_i$ are sub-gaussian as shown in the above proof of
Lemma \ref{lemma:outliersat0}(b), so Assumption~\ref{assump:G}(b) holds by
\cite[Eq.\ (5.26)]{vershynin2010introduction}, and
Assumption \ref{assump:G}(d) holds by \cite[Lemma 2]{jin2019short}.
The nonlinear Hanson-Wright inequality of
\cite[Eq.\ (3.3)]{wang2021deformed} implies
\[|\vg_i^\top\vA\vg_i-\Tr\vA\vSigma| \prec \|\vA\|_F\]
uniformly over $i \in [N]$ and deterministic matrices
$\vA \in \C^{n \times n}$. Furthermore, it is clear from the argument preceding
\cite[Eq.\ (3.3)]{wang2021deformed} that for any $i \neq j \in [N]$, the joint
vector $(\vg_i,\vg_j) \in \R^{2n}$ also satisfies Lipschitz concentration, hence
		\[\left|\begin{pmatrix}
			\vg_i^\top & \vg_j^\top
		\end{pmatrix} \vB
		\begin{pmatrix}
			\vg_i\\
			\vg_j
		\end{pmatrix}
-\Tr \vB\begin{pmatrix} \vSigma & \mathbf{0} \\
\mathbf{0} & \vSigma \end{pmatrix} \right| \prec \|\vB\|_F\]
uniformly over $i \neq j \in [N]$ and deterministic matrices $\vB \in \C^{2n
\times 2n}$. Applying this with
\[\vB=\begin{pmatrix} \mathbf{0} &\vA\\ \vA &\mathbf{0} \end{pmatrix}\]
verifies both statements of Assumption \ref{assump:G}(c).

Next, we check Assumption~\ref{assum:spike} for the population
covariance matrix $\vSigma$.
Combining Assumption~\ref{assump:data} and \eqref{eq:lambda_i_diff} from
Lemma~\ref{lemm:phi_0}, we have
\begin{align}
\frac{1}{n-r}\sum_{i=r+1}^n \delta_{\lambda_i(\vSigma)}
&\to \nu_0:=b_\sigma^2 \otimes \mu_0 \oplus (1-b_\sigma^2) \text{ weakly},
\label{eq:onelayernu}\\
\lambda_i(\vSigma) &\to -\frac{1}{s_{i,0}}:=b_\sigma^2\lambda_i+(1-b_\sigma^2)
\not\in\supp{\nu_0} \text{ for } i=1,\ldots,r.
\label{eq:onelayersi}
\end{align}
Here, the statement $-1/s_{i,0} \notin \supp{\nu_0}$ in \eqref{eq:onelayersi}
follows from the assumptions
$\lambda_i \notin \supp{\mu_0}$ and $b_\sigma \neq 0$.
This then implies by the definition of $\cT_1$
that $s_{i,0} \in \R \setminus \cT_1$, as claimed in Theorem \ref{thm:ck_spike}(a).
Furthermore, for any fixed $\eps>0$ and all large $n$,
Assumption~\ref{assump:data} and \eqref{eq:lambda_i_diff} imply also that
	\[\lambda_i(\vSigma) \in \supp{\nu_0}+(-\eps,\eps) \text{ for all }
	i \geq r+1.\]
Thus Assumption~\ref{assum:spike} holds for $\vSigma$ as $n \to \infty$.

Then we can apply Theorems~\ref{thm:deterministic_equ} and \ref{thm:spiked}
for $\bar\vK:=\vG^\top\vG$. The Stieltjes transform approximation in Theorem
\ref{thm:deterministic_equ} and Lemma \ref{lemma:supportcompareb}(c) together
imply $m_{\bar\vK}(z) \to m_1(z)$ almost surely for each fixed $z \in \C^+$,
where $m_1(z)$ is the Stieltjes transform of the measure
$\mu_1=\rho_{\gamma_1}^\MP \boxtimes \nu_0$. This implies the weak convergence
\begin{equation}\label{eq:Kbarweakconvergence}
\frac{1}{n}\sum_{i=1}^n \delta_{\lambda_i(\bar\vK)} \to \mu_1 \text{ a.s.}
\end{equation}
Theorem~\ref{thm:spiked}(a,b) further justifies:
	\begin{itemize}
\item Let $z_1(\cdot)$ and $\cI_1$ be defined by \eqref{eq:zell} and
\eqref{eq:index_set} with $\ell=1$. Then for any sufficiently small
constant $\eps>0$, almost surely for all large $n$,
there is a 1-to-1 correspondence between the eigenvalues of $\bar\vK$ outside
$\cS_1+(-\eps,\eps)$ and $\{i:i \in \cI_1\}$. Furthermore,
for each $i \in \cI_1$,
\begin{equation}\label{eq:Kbareigval}
\lambda_i(\bar\vK) \to z_1(s_{i,0})>0.
\end{equation}
almost surely as $n \to \infty$.

\item Let $\varphi_1(\cdot)$ be defined by \eqref{eq:zell} with $\ell=1$.
For each $i \in \cI_1$, let $\vv_i(\bar\vK) \in \R^n$
be a unit-norm eigenvector of $\bar\vK$
corresponding to $\lambda_i(\bar\vK)$, and for each $j \in [r]$,
let $\vv_j(\vSigma)$ be a unit-norm eigenvector of $\vSigma$
corresponding to $\lambda_j(\vSigma)$. Then
almost surely as $n \to \infty$, for each $i \in \cI_1$ and $j \in [r]$,
\begin{equation}\label{eq:Kbareigvec1}
|\vv_j(\vSigma)^\top \vv_i(\bar\vK)| \to \sqrt{\varphi_1(s_{i,0})} \cdot \bI\{i=j\}
\end{equation}
where $\varphi_1(s_{i,0})>0$. Moreover,
letting $\vv\in\R^n$ be any unit vector independent of $\vW$, almost surely
\begin{equation}\label{eq:Kbareigvec2}
|\vv^\top \vv_i(\bar\vK)|-\sqrt{\varphi_1(s_{i,0})}
\cdot|\vv^\top\vv_i(\vSigma)| \to 0.
\end{equation}
\end{itemize}

If $\gamma_1 \geq 1$, then Lemma \ref{lemma:outliersat0}(a) shows that
$\supp{\mu_1}=\supp{\mu_1} \cup \{0\}=\cS_1$. If $\gamma_1<1$, then
Lemma \ref{lemma:outliersat0}(b) shows that for any sufficiently small constant
$\eps>0$, $\bar\vK$ has no eigenvalues in $[0,\eps)$ almost surely
for all large $n$. Thus, in both cases, the first statement above in fact
establishes a 1-to-1 correspondence between $\{i:i \in \cI_1\}$ and
all eigenvalues of $\bar\vK$ outside $\supp{\mu_1}+(-\eps,\eps)$, almost surely
for all large $n$.

To translate these statements to the non-centered matrix
$\vK=\vY^\top\vY$, recall from Lemma~\ref{lemm:expect_rows} that $\|\vY-\vG\|
\to 0$ almost surely, and from Assumption \ref{assump:G}(b) verified above that 
$\bI\{\|\vG^\top \vG\|>B'\} \prec 0$ for a constant $B'>0$.
Then, almost surely as $n \to \infty$,
	\[\norm{\vK-\bar\vK} \to 0.\]
Therefore, by Weyl's inequality and \eqref{eq:Kbarweakconvergence}, the
empirical eigenvalue distribution $\widehat \mu_1$ of $\vK$ converges also to
$\mu_1$ weakly a.s., as claimed in Theorem \ref{thm:no_outlier_ck}.
Furthermore, by \eqref{eq:Kbareigval}, almost surely for all
large $n$, the eigenvalues $\widehat\lambda_{i,1}$ of $\vK$ outside
$\supp{\mu_1}+(-\eps,\eps)$ are also in 1-to-1
correspondence with $\{i:i \in \cI_1\}$, where
$\widehat\lambda_{i,1} \to z_1(s_{i,0})$ for each $i \in \cI_1$. In particular,
if $r=0$, then also $|\cI_1|=0$, so $\vK$ has no eigenvalues outside
$\supp{\mu_1}+(-\eps,\eps)$. This proves Theorem \ref{thm:no_outlier_ck}.

For each $i \in \cI_1$, let $\widehat \vv_{i,1} \in \R^n$ be a
	unit-norm eigenvector of $\vK$ 
	corresponding to $\widehat \lambda_{i,1}$.
Then by the Davis-Kahan Theorem \cite{davis1970rotation},
we may choose a sign for $\widehat \vv_{i,1}$ such that
	\[\norm{\widehat\vv_{i,1}-\vv_i(\bar\vK)} \le
\frac{\sqrt{2}\|\vK-\bar{\vK}\|}{\dist{
\widehat{\lambda}_{i,1},\spec{(\bar\vK)}\setminus \{\lambda_i(\bar\vK)\}}}.\]
We note that $\widehat{\lambda}_{i,1} \to z_1(s_{i,0})$ a.s., which is distinct
from the limit values $\{z_1(s_{j,0}):j \in \cI_1 \setminus \{i\}\}$ of
$\{\lambda_j(\bar \vK):\cI_1 \setminus \{i\}\}$
by bijectivity of the map $z_1(\cdot)$ in Proposition
\ref{prop:inv_z}. Furthermore $z_1(s_{i,0})$ falls outside
$\supp{\mu_1}+(-\eps,\eps)$ for sufficiently small $\eps>0$,
which contains all other eigenvalues of $\bar \vK$.
Thus $\dist{\widehat{\lambda}_{i,1},\spec{(\bar\vK)}\setminus
\{\lambda_i(\bar\vK)\}} \geq c$ for a constant $c>0$ almost surely for
all large $n$, so
\[\norm{\widehat\vv_{i,1}-\vv_i(\bar\vK)} \to 0 \text{ a.s.}\]
Similarly, by the convergence
$\|\vSigma-\vSigma_{\mathrm{lin}}\| \to 0$ and the assumption $b_\sigma \neq 0$,
we have $\norm{\vv_j(\vSigma)-\vv_j} \to 0$ for each $j \in [r]$, where $\vv_j$
is the unit-norm eigenvector of $\vSigma_{\mathrm{lin}}$ corresponding to
its eigenvalue
$b_\sigma^2\lambda_j(\vX^\top \vX)+(1-b_\sigma^2)$, i.e.\ the
eigenvector of $\vX^\top \vX$ corresponding to $\lambda_j(\vX^\top \vX)$.
Then \eqref{eq:Kbareigvec1} and \eqref{eq:Kbareigvec2} imply also
\[|\vv_j^\top \widehat \vv_i|^2 \to \varphi_1(s_{i,0}) \cdot \bI\{i=j\},
\qquad |\vv^\top \widehat \vv_i|^2-\varphi_1(s_{i,0}) \cdot|\vv^\top\vv_i|^2 \to 0.\]
This shows all claims of Theorem \ref{thm:ck_spike} for $L=1$.

Finally, on $\cE(\vX)$, the matrix $\vY$ is $C\tau_n$-orthonormal for a
constant $C>0$ by \cite[Lemma D.3(b)]{fan2020spectra}.
(The proof of \cite[Lemma D.3(b)]{fan2020spectra} again does not use
the condition $\sum_\alpha (\|\vx_\alpha\|_2^2-1)^2 \leq B^2$ of
$(\eps,B)$-orthonormality therein, and the remaining conditions of
$(\eps,B)$-orthonormality hold under $\cE(\vX)$.)
This shows the last claim of the lemma.
\end{proof}

\subsection{Spike analysis for multiple layers}\label{subsec:propagation}

We now prove Theorem~\ref{thm:ck_spike} by inductively applying the result for
$L=1$ through multiple layers. We follow the notations of
Section~\ref{sec:spike_CK}.\\

\begin{proof-of-theorem}[\ref{thm:ck_spike}]
Suppose inductively that Assumption \ref{assump:data}
holds with $\vX_{\ell-1}$ in place of $\vX_0$, and all
conclusions of Theorem \ref{thm:ck_spike} hold for $\vK_\ell$. The base case of
$\ell=1$ follows from Lemma \ref{lemma:onelayer}.

Then the last statement of Lemma \ref{lemma:onelayer} implies that
$\vX_\ell$ is $\tau_n'$-orthonormal almost surely for all large $n$,
for some $\tau_n'$ satisfying $\tau_n' \cdot n^{1/3} \to 0$.
Furthermore, the conclusions of Theorem \ref{thm:ck_spike}(b,c) for $\vK_\ell$
imply that statements (a) and (b) of Assumption \ref{assump:data} also
hold for $\vX_\ell$, in the following sense: Let $r_\ell=|\cI_\ell|$. Then
\[\frac{1}{n-|r_\ell|}\sum_{i \notin \cI_\ell} \delta_{\lambda_i(\vX_\ell^\top
\vX_\ell)} \to \mu_\ell \text{ weakly a.s.}\]
For any fixed $\eps>0$, almost surely for all large $n$,
$\widehat \lambda_{i,\ell}:=
\lambda_i(\vX_\ell^\top \vX_\ell) \in \supp{\mu_\ell}+(-\eps,\eps)$
for all $i \notin \cI_\ell$. Furthermore, for each $i \in \cI_\ell$,
$\widehat \lambda_{i,\ell} \to z_\ell(s_{i,\ell-1}) \notin
\supp{\mu_\ell}$.

Then we may apply Lemma \ref{lemma:onelayer} with input data
$\vX=\vX_\ell$ in place of $\vX_0$. This shows that for any fixed $\eps>0$ and all large $n$, there is a 1-to-1 correspondence between the eigenvalues $\widehat
\lambda_{i,\ell+1}$ of $\vK_{\ell+1}$ outside
$\supp{\mu_{\ell+1}}+(-\eps,\eps)$ and $\{i:i\in\cI_{\ell+1}\}$,
where
$\widehat \lambda_{i,\ell+1}\to z_{\ell+1}\left(s_{i,\ell}\right)>0$ a.s.,
and $s_{i,\ell} \in \R \setminus \cT_{\ell+1}$.
Moreover, for any unit vector $\vv \in \R^n$ independent of
$\vW_1,\ldots,\vW_{\ell+1}$,
	\[|\widehat \vv_{i,\ell+1}^\top
\vv|^2-\varphi_{\ell+1}(s_{i,\ell})\cdot|\widehat \vv_{i,\ell}^\top \vv|^2\to 0,\]
where also $\varphi_{\ell+1}(s_{i,\ell})>0$.
Then by the induction hypothesis for $|\widehat \vv_{i,\ell}^\top \vv|^2$,
	\begin{align*}
	    |\widehat \vv_{i,\ell+1}^\top \vv|^2\to ~& \prod_{k=1}^{\ell+1}
\varphi_{k}(s_{i,k-1}) \cdot |\vv_{i}^\top \vv |^2,
	\end{align*}
and specializing to $\vv=\vv_j$ for $j \in [r]$ gives
	\[|\widehat \vv_{i,\ell+1}^\top \vv_j |^2\to \prod_{k=1}^{\ell+1}
\varphi_{k}(s_{i,k-1}) \cdot \bI\{i=j\}.\]
This verifies all conclusions of Theorem \ref{thm:ck_spike}
for $\vK_{\ell+1}$, completing the induction.
\end{proof-of-theorem}

\subsection{Corollary for signal-plus-noise input data}\label{subsec:GMM_proof}

\begin{proof-of-coro}[\ref{cor:gmm}]
It is shown in \cite[Section 3.1]{benaych2012singular} that asymptotically as
$d,n \to \infty$ with $n/d \to \gamma_0$, the data matrix $\vX$ has a spike
singular value corresponding to $\theta_i$ if and only if
$\theta_i>\gamma_0^{1/4}$, in which case
\[\lambda_i(\vK_0) \to
\lambda_i:=\frac{(1+\theta_i^2)(\gamma_0+\theta_i^2)}{\theta_i^2},
\qquad |\vb_i^\top \vv_i|^2 \to
1-\frac{\gamma_0(1+\theta_i^2)}{\theta_i^2(\theta_i^2+\gamma_0)}\]
where $\vv_i$ is the unit eigenvector of the input Gram matrix
$\vK_0=\vX^\top\vX$. Thus claims (a) and (b) of Assumption \ref{assump:data}
hold with $r=|\{i:\theta_i>\gamma_0^{1/4}\}|$,
$\mu_0=\rho_{\gamma_0}^\MP$ being the standard Marcenko-Pastur law, and
$\lambda_i=(1+\theta_i^2)(\gamma_0+\theta_i^2)/\theta_i^2$
being the above values.

We note that $\vX$ is $n^{-1/2+\eps}$-orthonormal for any $\eps>0$
almost surely for all large $n$, by the given condition
$\max_{i=1}^r \|\vb_i\|_\infty<n^{-1/2+\eps}$ and the bounds, for any
$\alpha,\beta \in [n]$,
\begin{align*}
\|\vx_\alpha\|_2&=\left\|\vz_\alpha\right\|_2
+\sum_{i=1}^r \SD{\|\va_i\|_2 |\theta_i| |b_{i,\alpha}|}
=\left\|\vz_\alpha\right\|_2+\SD{n^{-1/2+\eps}}=1+\SD{n^{-1/2+\eps}},\\
\vx_\alpha^\top \vx_\beta&=\vz_\alpha^\top \vz_\beta
+\sum_{i=1}^r \SD{|\theta_i| \Big(
|\va_i^\top \vz_\alpha| |b_{i,\alpha}|
+|\va_i^\top \vz_\beta| |b_{i,\beta}|\Big)
+\theta_i^2 \|\va_i\|_2^2 |b_{i,\alpha}b_{i,\beta}|}
=\SD{n^{-1/2+\eps}}.
\end{align*}
Hence Theorem \ref{thm:ck_spike} applies, showing
that $\vK_\ell$ has an outlier eigenvalue corresponding to each input
signal $\theta_i$ if and only if $\theta_i>\gamma_0^{1/4}$ and
$i \in \cI_\ell$. The statement \eqref{eq:gmmalignment} follows from 
Theorem \ref{thm:ck_spike}(c) applied with $\vv=\vb_i$.
\end{proof-of-coro}

\section{Proofs for spiked eigenstructure of the trained CK}\label{sec:trained_CK_proof}

In this section, we prove Theorem~\ref{thm:gd_spike}. The proof is an
application of Theorem \ref{thm:spiked} as in the one-layer setting of the
preceding section, but now reversing the roles of $\vX$ and $\vW$. We abbreviate
\[\vW=\vW_{\text{trained}}, \qquad
\vY=\frac{1}{\sqrt{N}}\sigma(\tilde{\vX}\vW), \qquad
\vK=\vY\vY^\top\]
where $\vK$ is the CK matrix of interest. In contrast to the preceding section,
the theorem requires characterizing the \emph{left} spike singular vector of
$\vY$, and we will do so using Theorem \ref{thm:spiked}(c).

We first recall the following approximation of $\vW$ from \cite{ba2022high}.
\begin{prop} \label{prop:W_spike}
Under Assumption \ref{assum:GD}, set $\eta=\sum_{t=1}^T \eta_t$, and let
$\theta_1,\theta_2$ be as defined in \eqref{def:theta_12}. Then
	\begin{align}
\|\vW-\widetilde \vW\| \prec N^{-1/2} \quad \text{ where } \quad
		\widetilde\vW=\vW_0+\frac{b_\sigma\eta}{n}\vX^\top\vy\va^\top.\label{eq:tilde_W}
	\end{align}
The largest singular value $s_{\max}(\vW)$ falls outside the limit of its
empirical singular value distribution if and only if
$\theta_1>\gamma_0^{1/4}$, in which case $s_{\max}(\vW)$ 
and its unit-norm left singular vector $\vu(\vW)$ satisfy
	\begin{equation}\label{eq:trainedWspikes}
			s_{\max}(\vW)\to s_1:=
\sqrt{\frac{(1+\theta_1^2)(\gamma_0+\theta_1^2)}{\theta_1^2}}, \quad
		|\vu(\vW)^\top
\vbeta_*|^2\to\frac{\theta_2^2}{\theta_1^2}\left(1-\frac{\gamma_0+\theta_1^2}{\theta_1^2(\theta_1^2+1)}\right)
\quad \text{a.s.}
\end{equation}
\end{prop}
\begin{proof}
Notice that each gradient update matrix
$\vG_t$ of \eqref{eq:trained_W} takes the form
\begin{align*} 
	\vG_t=
   \frac{1}{n} \vX^\top \left[\left(\frac{1}{\sqrt{N}}\left(\vy -
\frac{1}{\sqrt{N}}\sigma(\vX\vW_{\!t})\va\right)\va^\top\right) \odot
\sigma'(\vX\vW_{\!t})\right], 
\end{align*}  
From the proof of \cite[Lemma 16]{ba2022high}, for each $t=1,\ldots,T$, this
matrix $\vG_t$ satisfies the same rank-one approximation
\[\Big\|\sqrt{N}\,\vG_t-\frac{b_\sigma}{n}\vX^\top\vy\va^\top\Big\|
\prec N^{-1/2}.\]
This implies \eqref{eq:tilde_W} in light of \eqref{eq:trained_W},
and the statements of \eqref{eq:trainedWspikes} then
follow from  \cite[Theorem 3]{ba2022high}.
\end{proof}

We denote the columns of $\vW \equiv
\vW_{\text{trained}} \in \R^{d \times N}$
and of the initialization $\vW_0 \in \R^{d \times N}$ by
\[\vw_i \in \R^d,\; \vw_{i,0} \in \R^d \text{ for } i \in [N]\]
respectively. Fixing a large constant $B>0$ and small constant $\eps>0$,
define the event
\[\cE(\vW)=\Big\{\|\vW\|<B,\,|\vw_i^\top \vw_j|<n^{-1/2+\eps}
\text{ and } |\|\vw_i\|_2-1|<n^{-1/2+\eps} \text{ for all } i \neq j \in
[N]\Big\}.\]

\begin{lemma}\label{lemma:event_Wt}
Under Assumption \ref{assum:GD}, for some sufficiently large constant
$B>0$ and any fixed $\eps>0$, $\cE(\vW)$ holds almost surely for all large $n$ and any fixed $T\in\N$.
\end{lemma}
\begin{proof}
By the assumption $[\vW_0]_{ij} \overset{iid}{\sim} \cN(0,1/d)$, it is
immediate to check that $\cE(\vW_0)$ holds almost surely for all large $n$.
To show that $\cE(\vW)$ holds, we apply the approximation \eqref{eq:tilde_W}.
Here, under Assumption \ref{assum:GD},
we have by standard tail bounds for Gaussian vectors and matrices that
\[\bI\{\|\vX^\top \vy\|>Cn\}
\leq \bI\{\|\vX\| \cdot (\lambda_{\sigma_*}\|\vX\vbeta_*\|+\|\vvarepsilon\|_2)>Cn\} \prec 0,
\qquad \bI\{\|\va\|_2>C\} \prec 0\]
for a sufficiently large constant $C>0$, and also
$\|\va\|_\infty \prec N^{-1/2}$. Then this implies
\begin{align}
\max_{i=1}^N \|\vw_i-\vw_{i,0}\|_2
&\leq \max_{i=1}^N \|\widetilde \vw_i-\vw_{i,0}\|_2
+\|\vW-\widetilde \vW\| \prec N^{-1/2}\label{eq:W_12_norm}
\end{align}
and $\bI\{\|\vW-\vW_0\|>C'\} \prec 0$ for a constant $C'>0$. Then $\cE(\vW)$
also holds almost surely for all large $n$, as claimed.
\end{proof}

Analogous to the argument of Appendix \ref{subsec:onelayer},
we may now condition on $\vW$, i.e.\ we assume that $\vW$ is deterministic and
satisfies $\cE(\vW)$ for all large $n$, and we
write $\E$ for the expectation
over only the randomness of the new data $(\tilde \vX,\tilde \vy)$.
Defining
\begin{equation}\label{eq:udef}
\vG=\sqrt{\frac{N}{n}}(\vY-\E\vY) \in \R^{n \times N},
\quad \vu=\frac{1}{\sqrt{n}}\tilde\vy \in \R^n
\quad \text{ where } \quad
\vY=\frac{1}{\sqrt{N}}\sigma(\tilde \vX\vW),
\end{equation}
observe that $[\vu,\vG] \in \R^{n \times (N+1)}$ has centered i.i.d.\ rows with
respect to the randomness of $(\tilde \vX,\tilde\vy)$.
We will write $\E_{\vx}$ for the expectation with respect to a standard Gaussian
vector $\vx \sim \cN(\mathbf{0},\vI_d)$.

\begin{lemma}\label{lemma:bar_sigma_app}
Suppose $\vW$ satisfies $\cE(\vW)$ for all large $n$. Then
\begin{align}
\|\E_{\vx}[\sigma(\vx^\top\vW)]\|_2 \to 0,
\qquad \|\E\vY\| \to 0,
\qquad \norm{\vSigma-\vSigma_{\mathrm{lin}}} \to 0
\end{align}
        where
\begin{align}
	\vSigma:=~&\E_{\vx}[\sigma(\vx^\top\vW)^\top\sigma(\vx^\top\vW)]-\E_{\vx}[\sigma(\vx^\top\vW)]^\top\E_{\vx}[\sigma(\vx^\top\vW)]\\
	\vSigma_{\mathrm{lin}}:=~& b_\sigma^2(\vW^\top \vW)+(1-b_\sigma^2)\vI_N.
\end{align}
\end{lemma}
\begin{proof}
The proof is the same as Lemmas~\ref{lemm:expect_rows} and~\ref{lemm:phi_0}.
\end{proof}

\begin{proof-of-theorem}[\ref{thm:gd_spike}]
We condition on $\vW$ satisfying $\cE(\vW)$ for all large $n$, and we
apply Theorem~\ref{thm:spiked}(c) for $[\vu,\vG] \in \R^{(n+1)\times N}$
(exchanging $n$ and $N$). It may be checked that
Assumption~\ref{assump:G} holds for $[\vu,\vG]$ by the same argument as
in Lemma~\ref{lemma:onelayer}.

By the convergence $\|\vSigma-\vSigma_{\mathrm{lin}}\| \to 0$ in
Lemma \ref{lemma:bar_sigma_app} and Proposition \ref{prop:W_spike},
if $\theta_1>\gamma_0^{1/4}$, then
Assumption \ref{assum:spike} holds for $\vSigma$ with $r=1$ and
\[\nu=b_\sigma^2\otimes\rho^\MP_{\gamma_0}\oplus(1-b_\sigma^2), \qquad
\lambda_1=b_\sigma^2\frac{(1+\theta_1^2)(\gamma_0+\theta_1^2)}{\theta_1^2}
+(1-b_\sigma^2) \notin \supp{\nu},\]
where $\rho_{\gamma_0}^\MP$ is the standard Marcenko-Pastur limit
for the empirical eigenvalue distribution of $\vW^\top \vW$, hence $\nu$ is
the limit empirical eigenvalue distribution of $\vSigma$, and $\lambda_1$ is the
limit of $\lambda_{\max}(\vSigma)$. If instead $\theta_1 \leq \gamma_0^{1/4}$,
then Assumption \ref{assum:spike} holds with $r=0$.

Then Theorem \ref{thm:spiked}(a,c)
characterizes the outlier eigenvalue and eigenvector of $\vG\vG^\top$, showing:
\begin{itemize}
	\item $\vG\vG^\top$ has a spike eigenvalue if and only if
$\theta_1>\gamma_0^{1/4}$ and $z'(-1/\lambda_1)>0$, where $z(\cdot)$ is defined 
by \eqref{eq:inv_z} with $\gamma=\gamma_1$ and the measure $\nu$ given above.
In this case, $\lambda_{\max}(\vG\vG^\top)\to z(-1/\lambda_1)$ almost surely.
	\item When $\theta_1>\gamma_0^{1/4}$ and $z'(-1/\lambda_1)>0$,
letting $\vu(\vG),\vv(\vSigma)$ be the leading unit-norm left singular vector
	of $\vG$ and leading unit-norm eigenvector of $\vSigma$, almost surely
			\[|\vu^\top \vu(\vG)|-
	\frac{\sqrt{z(-1/\lambda_1)\varphi({-}1/\lambda_1)}}{\lambda_1}
	\cdot\left|\E_{\vx}[\sigma_*(\vbeta_*^\top\vx)\sigma(\vx^\top\vW)]\vv(\vSigma)
	\right| \to 0.\]
where $\varphi(\cdot)$ is defined by \eqref{eq:phi_limit} also
with $\gamma=\gamma_1$ and the above measure $\nu$.
\end{itemize}

By an application of Weyl's inequality and  the Davis-Kahan Theorem as in the proof of
Theorem~\ref{thm:ck_spike}, this implies for $\vK=\vY\vY^\top$ that if
$\theta_1>\gamma_0^{1/4}$ and $z'(-1/\lambda_1)>0$, then its leading
eigenvalue $\lambda_{\max}(\vK)$ and unit eigenvector $\widehat \vu$ satisfy
\begin{equation}\label{eq:learnedKspike}
	\begin{aligned}
		&\lambda_{\max}(\vK)\to \gamma_1^{-1}z(-1/\lambda_1),\\
		&|\vu^\top \widehat\vu|-
			\frac{\sqrt{z(-1/\lambda_1)\varphi({-}1/\lambda_1)}}{\lambda_1}
			\cdot\left|\E_{\vx}[\sigma_*(\vbeta_*^\top\vx)\sigma(\vx^\top\vW)]\vv(\vW)\right|
			\to 0,
	\end{aligned}
\end{equation}
where $\vv(\vW)$ is the leading unit eigenvector of $\vSigma_{\text{lin}}$,
i.e.\ the leading right singular vector of $\vW$.
If $\theta_1 \leq \gamma_0^{1/4}$ or $z'(-1/\lambda_1) \leq 0$, then all
eigenvalues of $\vK$ converge to the support of its limiting empirical
eigenvalue law.

Finally, in the case of $\theta_1>\gamma_0^{1/4}$ and $z'(-1/\lambda_1)>0$,
we may conclude the proof by showing
\begin{equation}\label{eq:align_y_limit}
	\left\|\E_{\vx}[\sigma_*(\vbeta_*^\top\vx)\sigma(\vx^\top\vW)]
	-b_\sigma b_{\sigma_*} \vbeta_*^\top\vW\right\|_2 \to 0 \text{ a.s.}
\end{equation}
For each column $i\in [N]$, we have from \eqref{eq:W_12_norm} that
$\|\vw_i-\vw_{i,0}\|_2 \prec N^{-1/2}$, where $\vw_{i,0} \sim \cN(0,d^{-1}\vI)$
and $\vbeta_*$ is deterministic. Hence $(\vw_i,\vbeta_*)$ satisfy the approximate
orthonormality conditions $|\|\vw_i\|_2-1| \prec N^{-1/2}$, $\|\vbeta_*\|_2-1=0$,
and $|\vw_i^\top \vbeta_*| \prec N^{-1/2}$. Then
\cite[Lemma D.3(a)]{fan2020spectra} implies
\begin{align*}
	\left|\E_{\vx}[\sigma_*(\vbeta_*^\top\vx)\sigma(\vx^\top\vw_i)]-b_\sigma
b_{\sigma_*} \vbeta_*^\top\vw_i\right|\prec N^{-1}.
\end{align*} 
(We note that \cite[Lemma D.3(a)]{fan2020spectra} assumes $\sigma=\sigma_*$, but the proof is
identical for $\sigma \neq \sigma_*$ both satisfying
Assumption \ref{assump:sigma}.)
Applying this to each coordinate $i \in [N]$ yields \eqref{eq:align_y_limit}.
Observe that $\vbeta_*^\top\vW\vv(\vW)=s_{\max}(\vW) \cdot \vbeta_*^\top\vu(\vW)$
where $s_{\max}(\vW)$ and $\vu(\vW)$ are the leading singular value and left
singular vector of $\vW$, and recall from the definitions \eqref{eq:udef}
that $\vu=\frac{1}{\sqrt{n}}\tilde \vy$.
Then we can apply \eqref{eq:align_y_limit} and
Proposition~\ref{prop:W_spike} to \eqref{eq:learnedKspike} to conclude that 
\[\frac{1}{\sqrt{n}}|\tilde\vy^\top\widehat{\vu}|\to b_\sigma b_{\sigma_*}
\frac{\sqrt{z(-1/\lambda_1)\varphi(-1/\lambda_1)}}{\lambda_1}\cdot
\frac{\theta_2\sqrt{(\theta_1^4-\gamma_0)(\gamma_0+\theta_1^2)}}{\theta_1^3}
>0 \text{ a.s.}\]
\end{proof-of-theorem}

